% Applications linéaires et affines — Mathématiques 3ème — Cours signé OnBuch+
% Programme officiel MINESEC. Compiler avec : tectonic N16-applications-lineaires-affines.tex
\newcommand{\DOCMATIERE}{Mathématiques}
\newcommand{\DOCNIVEAU}{3ème}
\newcommand{\DOCMODULE}{Mathématiques – classe de 3ème}
\newcommand{\DOCLECON}{N16}
\newcommand{\DOCTITRE}{Applications linéaires et affines}
% OnBuch+ — préambule « cours signés » ENRICHI (compilé avec tectonic / XeLaTeX)
% Système visuel « Pop » d'OnBuch+ : fond crème (jamais blanc), bordures encre
% épaisses, ombres « dures » décalées, titres Archivo Black en capitales.
% Version MATHÉMATIQUES : graphes (pgfplots/tikz), boîtes pédagogiques
% (définition, propriété, méthode, attention, le savais-tu, exercice résolu,
% exercice, TP, bilan), page de garde et signature OnBuch+ ancrée en bas.
%
% Macros attendues dans main.tex AVANT \input{preamble} :
%   \DOCMATIERE  \DOCNIVEAU  \DOCMODULE  \DOCLECON (numéro)  \DOCTITRE
\documentclass[11pt]{article}

\usepackage{fontspec}
\usepackage[french]{babel}
\usepackage[a4paper,margin=2cm,top=3.4cm,bottom=2.8cm,headheight=1.4cm,footskip=1.3cm]{geometry}
\usepackage{amsmath,amssymb,amsfonts}
\usepackage{mathtools} % \xmapsto, \coloneqq... souvent utilisés par les modèles
\usepackage{cancel} % simplifications barrées
% \abs{z} (module, valeur absolue) et \norm{u} (norme), utilisés par les modèles
\providecommand{\abs}{}\let\abs\relax\DeclarePairedDelimiter{\abs}{\lvert}{\rvert}
\providecommand{\norm}{}\let\norm\relax\DeclarePairedDelimiter{\norm}{\lVert}{\rVert}
\usepackage[table]{xcolor} % [table] : fournit aussi \rowcolors (alternance de lignes)
\usepackage{pagecolor}
\usepackage{tikz}
\usetikzlibrary{arrows.meta,positioning,calc,shapes.geometric,shapes.misc,shapes.arrows,decorations.pathmorphing,decorations.pathreplacing,decorations.markings,patterns,shadows,mindmap,trees,fit,backgrounds,matrix,intersections,angles,quotes}
\usepackage{pgfplots}
\usepackage{pgf-pie} % diagrammes circulaires \pie (statistiques)
\pgfplotsset{compat=1.18}
\usepgfplotslibrary{fillbetween,groupplots} % aires entre courbes (calcul intégral)
\usepackage{enumitem}
\usepackage{fancyhdr}
% [most] charge toutes les bibliothèques tcolorbox (listings, minted, raster,
% documentation...) et gonfle inutilement la mémoire TeX sur un document long
% (~300 boîtes) : on ne charge que ce qui est réellement utilisé.
\usepackage[skins,breakable]{tcolorbox}
\usepackage{graphicx}
\usepackage{colortbl}
\usepackage{booktabs}
\usepackage{tabularx}
\usepackage{diagbox} % case d'en-tête diagonale des tableaux à double entrée
% Colonnes extensibles centrée / gauche / droite (C, L, R), souvent utilisées par les modèles
\newcolumntype{C}{>{\centering\arraybackslash}X}
\newcolumntype{L}{>{\raggedright\arraybackslash}X}
\newcolumntype{R}{>{\raggedleft\arraybackslash}X}
\usepackage{array}
\usepackage{multicol}
\usepackage{siunitx}
% parse-numbers=false : accepte aussi \SI{2,5 \times 10^{-3}}{...}
\sisetup{locale=FR,output-decimal-marker={,},group-minimum-digits=4,parse-numbers=false}
\DeclareSIUnit\ohm{\ensuremath{\symup{\Omega}}} % U+2126 absent de la police
\usepackage{qrcode}
% \degree : commande fréquemment utilisée par les modèles pour un angle ou une
% température en dehors de \SI{}{} (non fournie par les paquets chargés).
\providecommand{\degree}{^{\circ}}
% \qed (fin de démonstration), utilisé par les modèles sans amsthm
\providecommand{\qed}{\hfill$\square$}
% \barre : double barre des valeurs interdites dans un tableau de variations (tkz-tab)
\providecommand{\barre}{\ensuremath{\|}}
% environnement proof (amsthm non chargé) utilisé par les modèles
\ifdefined\proof\else\newenvironment{proof}[1][Démonstration]{\par\noindent\textit{#1.}\ }{\qed\par}\fi
\usepackage{lastpage}
\usepackage{hyperref}
\hypersetup{hidelinks}

% ============================================================
% Palette « Pop » OnBuch+ — mêmes hex que la classe P (plus_ui.dart)
% ============================================================
\definecolor{poporange}{HTML}{F59321}
\definecolor{popdark}{HTML}{E07A0C}
\definecolor{popcream}{HTML}{FFF6E8}
\definecolor{popink}{HTML}{1C1714}
\definecolor{poppurple}{HTML}{7A5AE0}
\definecolor{popgreen}{HTML}{2E9E6A}
\definecolor{poppink}{HTML}{E0457B}
\definecolor{popblue}{HTML}{2F7DE1}
\definecolor{popgold}{HTML}{C98A12}
\definecolor{popmuted}{HTML}{8A7F76}
\definecolor{popline}{HTML}{EADFD2}
\definecolor{popgray}{HTML}{8A7F76}
\colorlet{popgrey}{popgray}
% Alias tolérants (le modèle nomme parfois les couleurs différemment)
\colorlet{popwhite}{white}
\colorlet{popblack}{popink}
\colorlet{popred}{poppink}
\colorlet{popyellow}{popgold}
% Teintes claires pour fonds de tableaux / zones
\colorlet{popblueL}{popblue!10!white}
\colorlet{poporangeL}{poporange!14!white}
\colorlet{popgreenL}{popgreen!12!white}
\colorlet{poppurpleL}{poppurple!12!white}
\colorlet{poppinkL}{poppink!12!white}
\colorlet{popgoldL}{popgold!14!white}

\pagecolor{popcream}

% ============================================================
% Typographie
% ============================================================
\defaultfontfeatures{Ligatures=TeX}
\setmainfont{PlusJakartaSans-Medium.ttf}[Path=../../../fonts/,BoldFont=PlusJakartaSans-Bold.ttf]
% Maths en sans-serif (Fira Math) avec chiffres et lettres droites de Plus
% Jakarta, pour que les formules se fondent dans le texte.
\usepackage[math-style=ISO,bold-style=ISO]{unicode-math}
\setmathfont{FiraMath-Regular.otf}[Path=../../../fonts/]
\setmathfont{PlusJakartaSans-Medium.ttf}[Path=../../../fonts/,range={up/num,up/{Latin,latin},\mathup}]
\usepackage{icomma} % virgule décimale sans espace en mode math
% Caractères mathématiques Unicode absents des polices de texte (Plus Jakarta,
% Archivo Black) : rendus par leur équivalent mathématique, en texte comme en
% formule — sinon ils s'affichent en carré vide (ex. « Arithmétique dans ℤ »).
\usepackage{newunicodechar}
\newunicodechar{ℝ}{\ensuremath{\mathbb{R}}}
\newunicodechar{ℤ}{\ensuremath{\mathbb{Z}}}
\newunicodechar{ℂ}{\ensuremath{\mathbb{C}}}
\newunicodechar{ℕ}{\ensuremath{\mathbb{N}}}
\newunicodechar{ℚ}{\ensuremath{\mathbb{Q}}}
\newunicodechar{𝔻}{\ensuremath{\mathbb{D}}}
\newunicodechar{α}{\ensuremath{\alpha}}
\newunicodechar{β}{\ensuremath{\beta}}
\newunicodechar{γ}{\ensuremath{\gamma}}
\newunicodechar{δ}{\ensuremath{\delta}}
\newunicodechar{ε}{\ensuremath{\varepsilon}}
\newunicodechar{θ}{\ensuremath{\theta}}
\newunicodechar{λ}{\ensuremath{\lambda}}
\newunicodechar{μ}{\ensuremath{\mu}}
\newunicodechar{σ}{\ensuremath{\sigma}}
\newunicodechar{τ}{\ensuremath{\tau}}
\newunicodechar{φ}{\ensuremath{\varphi}}
\newunicodechar{ψ}{\ensuremath{\psi}}
\newunicodechar{ω}{\ensuremath{\omega}}
\newunicodechar{Σ}{\ensuremath{\Sigma}}
% majuscules produites par \MakeUppercase dans les titres (α-aminés -> Α)
\newunicodechar{Α}{\ensuremath{\alpha}}
\newunicodechar{Β}{\ensuremath{\beta}}
\newunicodechar{Γ}{\ensuremath{\Gamma}}
\newunicodechar{Δ}{\ensuremath{\Delta}}
\newunicodechar{Ε}{\ensuremath{\varepsilon}}
\newunicodechar{Θ}{\ensuremath{\Theta}}
\newunicodechar{Λ}{\ensuremath{\Lambda}}
\newunicodechar{Μ}{\ensuremath{\mu}}
\newunicodechar{Τ}{\ensuremath{\tau}}
\newunicodechar{Φ}{\ensuremath{\Phi}}
\newunicodechar{Ψ}{\ensuremath{\Psi}}
\newunicodechar{Ω}{\ensuremath{\Omega}}
\newunicodechar{↦}{\ensuremath{\mapsto}}
\newunicodechar{∈}{\ensuremath{\in}}
\newunicodechar{∉}{\ensuremath{\notin}}
\newunicodechar{∧}{\ensuremath{\wedge}}
\newunicodechar{⇔}{\ensuremath{\Leftrightarrow}}
\newunicodechar{⟺}{\ensuremath{\Longleftrightarrow}}
\newunicodechar{⇒}{\ensuremath{\Rightarrow}}
\newunicodechar{⟹}{\ensuremath{\Longrightarrow}}
\newunicodechar{∘}{\ensuremath{\circ}}
\newunicodechar{∪}{\ensuremath{\cup}}
\newunicodechar{∩}{\ensuremath{\cap}}
\newunicodechar{≡}{\ensuremath{\equiv}}
\newunicodechar{⊂}{\ensuremath{\subset}}
\newunicodechar{⊥}{\ensuremath{\perp}}
\newunicodechar{∃}{\ensuremath{\exists}}
\newunicodechar{∀}{\ensuremath{\forall}}
\newunicodechar{∛}{\ensuremath{\sqrt[3]{\,}}}
\newunicodechar{∜}{\ensuremath{\sqrt[4]{\,}}}
\newunicodechar{⃗}{\ensuremath{{}^{\to}}}
\newunicodechar{ⁿ}{\ifmmode^{n}\else\textsuperscript{\ensuremath{n}}\fi}
\newunicodechar{ˣ}{\ifmmode^{x}\else\textsuperscript{\ensuremath{x}}\fi}
\newunicodechar{ᵇ}{\ifmmode^{b}\else\textsuperscript{\ensuremath{b}}\fi}
\newunicodechar{ʳ}{\ifmmode^{r}\else\textsuperscript{\ensuremath{r}}\fi}
\newunicodechar{ᵘ}{\ifmmode^{u}\else\textsuperscript{\ensuremath{u}}\fi}
\newunicodechar{ʸ}{\ifmmode^{y}\else\textsuperscript{\ensuremath{y}}\fi}
\newunicodechar{ᵃ}{\ifmmode^{a}\else\textsuperscript{\ensuremath{a}}\fi}
\newunicodechar{ᵉ}{\ifmmode^{e}\else\textsuperscript{\ensuremath{e}}\fi}
\newunicodechar{ᵏ}{\ifmmode^{k}\else\textsuperscript{\ensuremath{k}}\fi}
\newunicodechar{ᵖ}{\ifmmode^{p}\else\textsuperscript{\ensuremath{p}}\fi}
\newunicodechar{ᴺ}{\ifmmode^{N}\else\textsuperscript{\ensuremath{N}}\fi}
\newunicodechar{ᶜ}{\ifmmode^{c}\else\textsuperscript{\ensuremath{c}}\fi}
\newunicodechar{ʰ}{\ifmmode^{h}\else\textsuperscript{\ensuremath{h}}\fi}
\newunicodechar{ᵠ}{\ifmmode^{q}\else\textsuperscript{\ensuremath{q}}\fi}
\newunicodechar{ₙ}{\ifmmode_{n}\else\textsubscript{\ensuremath{n}}\fi}
\newunicodechar{ₐ}{\ifmmode_{a}\else\textsubscript{\ensuremath{a}}\fi}
\newunicodechar{ᵢ}{\ifmmode_{i}\else\textsubscript{\ensuremath{i}}\fi}
\newunicodechar{ᵣ}{\ifmmode_{r}\else\textsubscript{\ensuremath{r}}\fi}
\newunicodechar{ₖ}{\ifmmode_{k}\else\textsubscript{\ensuremath{k}}\fi}
\newunicodechar{ᵦ}{\ifmmode_{\beta}\else\textsubscript{\ensuremath{\beta}}\fi}
\newunicodechar{ₓ}{\ifmmode_{x}\else\textsubscript{\ensuremath{x}}\fi}
\newfontfamily\popdisplay{ArchivoBlack-Regular.ttf}[Path=../../../fonts/]
\newfontfamily\poplogo{SpaceGrotesk-Bold.ttf}[Path=../../../fonts/]
\newfontfamily\popsemi{PlusJakartaSans-SemiBold.ttf}[Path=../../../fonts/]
\newfontfamily\popxbold{PlusJakartaSans-ExtraBold.ttf}[Path=../../../fonts/]
\newfontfamily\popxboldls{PlusJakartaSans-ExtraBold.ttf}[Path=../../../fonts/,LetterSpace=3.0]

\setlist[enumerate]{leftmargin=1.7em,itemsep=5pt,topsep=4pt}
\setlist[itemize]{leftmargin=1.7em,itemsep=4pt,topsep=4pt,label={\color{poporange}\textbullet}}
\setlength{\parindent}{0pt}
\setlength{\parskip}{5pt}
\renewcommand{\arraystretch}{1.25}

% pgfplots : style Pop par défaut
\pgfplotsset{
  every axis/.append style={axis line style={popink,line width=1pt},tick style={popink},
    grid style={popline,line width=0.5pt},label style={font=\small},tick label style={font=\footnotesize}},
  popcurve/.style={poporange,line width=1.8pt,no markers,smooth},
}
% Même style disponible dans un \draw TikZ simple (hors pgfplots)
\tikzset{popcurve/.style={poporange,line width=1.8pt}}

% ============================================================
% Pill / squiggle
% ============================================================
\newcommand{\poppill}[3]{%
  \tikz[baseline=-0.55ex]\node[draw=popink,line width=1.2pt,rounded corners=2.6pt,%
    fill=#2,inner xsep=6.5pt,inner ysep=3pt]{%
    \popxboldls\fontsize{7}{7}\selectfont\color{#3}\MakeUppercase{#1}};%
}
\newcommand{\popsquiggle}[1]{%
  \tikz[baseline=-0.4ex]{
    \draw[#1,line width=2.6pt,line cap=round]
      plot [smooth] coordinates {(0,0.09) (0.34,0.01) (0.68,0.21) (1.02,0.01) (1.36,0.10)};
  }%
}
% Mot-clé surligné (marqueur orange sous le texte)
\newcommand{\cle}[1]{{\popxbold\color{popdark}#1}}

% ============================================================
% En-tête / pied
% ============================================================
\pagestyle{fancy}
\fancyhf{}
\renewcommand{\headrulewidth}{0pt}
\renewcommand{\footrulewidth}{0pt}
\lhead{\raisebox{-0.15ex}{\poplogo\fontsize{13}{13}\selectfont\color{popink}OnBuch\color{poporange}+}}
\rhead{\poppill{\DOCMATIERE\ \textbullet\ \DOCNIVEAU\ \textbullet\ Leçon \DOCLECON}{white}{popink}}
\lfoot{\popxbold\fontsize{7.6}{7.6}\selectfont\color{popmuted}\MakeUppercase{Cours signé OnBuch+}\ \textbullet\ {\popsemi\DOCTITRE}}
\rfoot{\tikz[baseline=-0.6ex]\node[rounded corners=8.5pt,fill=popink,minimum height=17pt,minimum width=17pt,inner xsep=5pt,inner sep=0]{\popxbold\fontsize{8}{8}\selectfont\color{popcream}\thepage\,/\,\pageref*{LastPage}};}

% ============================================================
% Titres
% ============================================================
\newcommand{\chaptitre}[2]{%
  \begin{tcolorbox}[enhanced,colback=poporange,colframe=popink,boxrule=2.6pt,arc=11pt,
      drop shadow=popink,left=15pt,right=15pt,top=12pt,bottom=11pt,before skip=3pt,after skip=18pt]
    \poppill{#2}{popink}{popcream}\par\vspace{9pt}
    {\raggedright\hyphenpenalty=10000\exhyphenpenalty=10000\popdisplay\fontsize{22}{24}\selectfont\color{popcream}\MakeUppercase{#1}\par}
  \end{tcolorbox}
}
% Section numérotée : pastille orange avec le numéro + titre Archivo
\newcounter{popsec}
\newcommand{\coursec}[1]{%
  \stepcounter{popsec}\setcounter{popsub}{0}%
  \par\vspace{16pt}\noindent
  \begin{minipage}{\linewidth}
  \tikz[baseline=-0.7ex]\node[draw=popink,line width=1.6pt,fill=poporange,rounded corners=4pt,minimum size=22pt,inner sep=0]{\popdisplay\fontsize{12}{12}\selectfont\color{popcream}\thepopsec};%
  \hspace{8pt}{\popdisplay\fontsize{15}{17}\selectfont\color{popink}\MakeUppercase{#1}}\par
  \vspace{4pt}\noindent{\color{poporange}\rule{36pt}{3.6pt}}
  \end{minipage}\par\nopagebreak\vspace{8pt}
}
\newcounter{popsub}
\newcommand{\courssub}[1]{%
  \stepcounter{popsub}%
  \par\vspace{9pt}\noindent{\popxbold\fontsize{11.5}{13}\selectfont\color{popink}{\color{poporange}\thepopsec.\thepopsub}\hspace{6pt}#1}\par\nopagebreak\vspace{4pt}
}

% ============================================================
% Boîtes pédagogiques — même recette : fond blanc, bordure épaisse colorée,
% ombre dure encre, badge plein. Titre optionnel [#1].
% ============================================================
\tcbset{popbase/.style={breakable,enhanced,colback=white,boxrule=2.3pt,arc=9pt,
  shadow={3pt}{-3pt}{0pt}{popink},left=14pt,right=14pt,top=11pt,bottom=11pt,before skip=11pt,after skip=15pt,
  fontupper=\popsemi\fontsize{10.6}{14.5}\selectfont\color{popink},
  fontlower=\popsemi\fontsize{10.6}{14.5}\selectfont\color{popink}}}
% Coche dessinée (le glyphe ✓ n'existe pas dans les polices du document)
\newcommand{\popcheck}[1]{\tikz[baseline=-0.5ex]\draw[#1,line width=1.6pt,line cap=round,line join=round](-0.13,0.0)--(-0.04,-0.09)--(0.13,0.1);}
\newcommand{\popboxhead}[3]{\poppill{#1}{#2}{white}\ifx\relax#3\relax\else\hspace{6pt}{\popxbold\fontsize{10.6}{12}\selectfont\color{popink}#3}\fi\par\vspace{7pt}}

\newtcolorbox{aretenir}[1][]{popbase,colframe=popgreen,before upper={\popboxhead{À retenir}{popgreen}{#1}}}
\newtcolorbox{exemplebox}[1][]{popbase,colframe=poporange,before upper={\popboxhead{Exemple}{poporange}{#1}}}
\newtcolorbox{definition}[1][]{popbase,colframe=popblue,before upper={\popboxhead{Définition}{popblue}{#1}}}
\newtcolorbox{propriete}[1][]{popbase,colframe=poppurple,before upper={\popboxhead{Propriété}{poppurple}{#1}}}
\newtcolorbox{methode}[1][]{popbase,colframe=popgold,before upper={\popboxhead{Méthode}{popgold}{#1}}}
\newtcolorbox{attention}[1][]{popbase,colframe=poppink,before upper={\popboxhead{Attention}{poppink}{#1}}}
\newtcolorbox{experience}[1][]{popbase,colframe=popblue,colback=popblueL,before upper={\popboxhead{Expérience}{popblue}{#1}}}
% « Le savais-tu ? » : carte violette pleine (contexte, culture, Cameroun)
\newtcolorbox{savaistu}[1][]{popbase,colback=poppurple,colframe=popink,
  fontupper=\popsemi\fontsize{10.4}{14.2}\selectfont\color{white},
  before upper={\poppill{Le savais-tu\,?}{white}{poppurple}\ifx\relax#1\relax\else\hspace{6pt}{\popxbold\color{white}#1}\fi\par\vspace{7pt}}}
% Exercice résolu : énoncé en haut, \tcblower puis la solution sur fond crème
\newcounter{popexo}\newcounter{popexr}
\newtcolorbox{exoresolu}[1][]{popbase,colframe=poporange,colbacklower=popcream,
  segmentation style={popink,line width=1.2pt,dashed},
  before upper={\stepcounter{popexr}\popboxhead{Exercice résolu \thepopexr}{poporange}{#1}},
  before lower={\poppill{Solution}{popink}{popcream}\par\vspace{6pt}}}
% Exercice (non résolu en place) — la correction est dans la partie Corrigés
\newtcolorbox{exercice}[1][]{popbase,colframe=popink,
  before upper={\stepcounter{popexo}\popboxhead{Exercice \thepopexo}{popink}{#1}}}
% Corrigé
\newtcolorbox{corrige}[1][]{popbase,colframe=popgreen,colback=popgreenL,
  before upper={\popboxhead{Corrigé}{popgreen}{#1}}}
% Objectifs / prérequis / bilan
\newtcolorbox{objectifs}{popbase,colframe=popink,colback=poporangeL,
  before upper={\popboxhead{Objectifs de la leçon}{popink}{}}}
\newtcolorbox{prerequis}{popbase,colframe=popink,colback=popblueL,
  before upper={\popboxhead{Prérequis}{popblue}{}}}
\newtcolorbox{bilan}{popbase,colframe=popink,colback=white,boxrule=2.8pt,
  before upper={\popboxhead{Fiche bilan}{poporange}{L'essentiel à connaître pour l'examen}}}
% Figure encadrée avec légende
\newtcolorbox{popfigure}[1][]{enhanced,breakable=false,colback=white,colframe=popline,boxrule=1.4pt,arc=8pt,
  left=10pt,right=10pt,top=10pt,bottom=8pt,before skip=10pt,after skip=12pt,halign=center,
  lower separated=false,halign lower=center,
  fontlower=\popsemi\fontsize{9}{11.5}\selectfont\color{popmuted},
  #1}
\newcommand{\legende}[1]{\par\vspace{6pt}{\popsemi\fontsize{9}{11.5}\selectfont\color{popmuted}#1}}

% ============================================================
% Signature OnBuch+
% ============================================================
\newcommand{\poplogoinline}[1]{{\poplogo\fontsize{#1}{#1}\selectfont\color{popink}OnBuch\color{poporange}+}}
\newcommand{\signatureonbuch}{%
  \begin{tcolorbox}[enhanced,colback=popink,colframe=popink,boxrule=2.2pt,arc=10pt,
      drop shadow=poporange,left=14pt,right=14pt,top=10pt,bottom=10pt,before skip=0pt,after skip=0pt]
    \tikz[baseline=-0.7ex]\node[circle,fill=popgreen,draw=popcream,line width=1.3pt,minimum size=22pt,inner sep=0]{\popcheck{white}};%
    \hspace{9pt}\begin{minipage}[c]{0.62\linewidth}
      {\popxbold\fontsize{12}{13}\selectfont\color{popcream}Signé\ \textbullet\ {\poplogo OnBuch\color{poporange}+}}\par\vspace{2pt}
      {\raggedright\popsemi\fontsize{8}{10.5}\selectfont\color{popline}\MakeUppercase{Équipe pédagogique OnBuch+}\\\MakeUppercase{Conforme au programme officiel MINESEC}\par}
    \end{minipage}\hfill
    {\poplogo\fontsize{22}{22}\selectfont\color{popcream}OnBuch\color{poporange}+}
  \end{tcolorbox}%
}

% ============================================================
% Page de garde
% #1 titre · #2 matière · #3 niveau · #4 module · #5 n° leçon · #6 durée ·
% #7 description / objectifs (texte ou itemize)
% ============================================================
\newcommand{\pagegarde}[7]{%
  \thispagestyle{empty}
  \newgeometry{a4paper,margin=1.7cm,top=1.6cm,bottom=1.4cm}%
  \begin{tikzpicture}[remember picture,overlay]
    \fill[poporange!18] ([xshift=-3.2cm,yshift=-3.6cm]current page.north east) circle (4.6cm);
    \fill[poppurple!14] ([xshift=2.4cm,yshift=7.6cm]current page.south west) circle (3.8cm);
  \end{tikzpicture}%
  {\poplogo\fontsize{30}{30}\selectfont\color{popink}OnBuch\color{poporange}+}\hfill
  \raisebox{0.6ex}{\poppill{Cours signé \textbullet\ Premium}{popink}{popcream}}\par
  \vspace{1pt}\popsquiggle{poppurple}\par
  \vspace{8pt}
  \begin{tcolorbox}[enhanced,colback=poporange,colframe=popink,boxrule=2.8pt,arc=13pt,
      drop shadow=popink,left=18pt,right=18pt,top=12pt,bottom=13pt,after skip=10pt]
    \poppill{#2 \textbullet\ #3}{popink}{popcream}\hspace{5pt}\poppill{#4}{white}{popink}\par\vspace{8pt}
    {\popdisplay\fontsize{14}{14}\selectfont\color{popink}LEÇON #5}\par\vspace{4pt}
    {\raggedright\hyphenpenalty=10000\exhyphenpenalty=10000\popdisplay\fontsize{25}{27}\selectfont\color{popcream}\MakeUppercase{#1}\par}
    \vspace{8pt}
    {\popxbold\fontsize{9.5}{9.5}\selectfont\color{popink}\MakeUppercase{Durée conseillée : #6}}
  \end{tcolorbox}
  \begin{tcolorbox}[enhanced,colback=white,colframe=popink,boxrule=2.2pt,arc=10pt,
      drop shadow=popink,left=16pt,right=16pt,top=10pt,bottom=10pt,after skip=8pt]
    \poppill{Dans cette leçon}{poppurple}{white}\par\vspace{5pt}
    \popsemi\fontsize{9.3}{12.6}\selectfont\color{popink}#7
  \end{tcolorbox}
  \noindent\begin{minipage}[t]{0.487\linewidth}
    \begin{tcolorbox}[enhanced,colback=popgreen,colframe=popink,boxrule=2.2pt,arc=10pt,
        drop shadow=popink,left=11pt,right=11pt,top=10pt,bottom=10pt,height=3.1cm,valign=center]
      \tcbox[colback=white,colframe=popink,boxrule=1.3pt,arc=5pt,boxsep=0pt,nobeforeafter,
        left=3pt,right=3pt,top=3pt,bottom=3pt,valign=center]{\qrcode[height=1.9cm]{https://play.google.com/store/apps/details?id=com.luvvix.onbuch&hl=fr}}%
      \hspace{8pt}\begin{minipage}[c]{0.5\linewidth}
        {\popdisplay\fontsize{11}{12.5}\selectfont\color{white}\MakeUppercase{Télécharge OnBuch}}\par\vspace{4pt}
        {\raggedright\popsemi\fontsize{8.4}{11}\selectfont\color{white}Gratuit \textbullet\ Google Play\\Tuteur IA, annales, concours, résultats\par}
      \end{minipage}
    \end{tcolorbox}
  \end{minipage}\hfill
  \begin{minipage}[t]{0.487\linewidth}
    \begin{tcolorbox}[enhanced,colback=poppurple,colframe=popink,boxrule=2.2pt,arc=10pt,
        drop shadow=popink,left=11pt,right=11pt,top=10pt,bottom=10pt,height=3.1cm,valign=center]
      \tikz[baseline=-0.7ex]\node[circle,fill=white,draw=popink,line width=1.3pt,minimum size=1.6cm,inner sep=0]{\popdisplay\fontsize{24}{24}\selectfont\color{poppurple}+};%
      \hspace{8pt}\begin{minipage}[c]{0.6\linewidth}
        {\popdisplay\fontsize{11}{12.5}\selectfont\color{white}\MakeUppercase{Passe à OnBuch+}}\par\vspace{4pt}
        {\raggedright\popsemi\fontsize{8.4}{11}\selectfont\color{white}Tous les cours signés, Tuteur IA illimité, fiches PDF — onglet OnBuch+ de l'app\par}
      \end{minipage}
    \end{tcolorbox}
  \end{minipage}
  \vspace{8pt}
  \popcontenu
  \vfill
  \signatureonbuch
  \restoregeometry
  \newpage
}

% Rangée « Ce cours contient » (page de garde)
\newcommand{\poptile}[3]{% #1 couleur #2 gros texte #3 légende
  \begin{tcolorbox}[enhanced,colback=#1,colframe=popink,boxrule=2pt,arc=9pt,shadow={2.5pt}{-2.5pt}{0pt}{popink},
      left=6pt,right=6pt,top=9pt,bottom=9pt,halign=center,valign=center,height=1.75cm,width=\linewidth,nobeforeafter]
    {\popdisplay\fontsize{12.5}{13}\selectfont\color{white}\MakeUppercase{#2}}\par\vspace{3pt}
    {\centering\popsemi\fontsize{7.8}{9.5}\selectfont\color{white}#3\par}
  \end{tcolorbox}}
\newcommand{\popcontenu}{%
  \par\noindent{\popxboldls\fontsize{7.5}{7.5}\selectfont\color{popmuted}CE COURS SIGNÉ CONTIENT}\par\vspace{7pt}
  \noindent\begin{minipage}[t]{0.235\linewidth}\poptile{popblue}{Cours}{complet, illustré, pas à pas}\end{minipage}\hfill
  \begin{minipage}[t]{0.235\linewidth}\poptile{poporange}{Méthodes}{et exercices résolus}\end{minipage}\hfill
  \begin{minipage}[t]{0.235\linewidth}\poptile{poppink}{Exercices}{type examen + corrigés}\end{minipage}\hfill
  \begin{minipage}[t]{0.235\linewidth}\poptile{popgreen}{Fiche bilan}{l'essentiel à retenir}\end{minipage}\par}

% Sommaire
\newtcolorbox{sommaire}{enhanced,colback=white,colframe=popink,boxrule=2.2pt,arc=10pt,
  drop shadow=popink,left=18pt,right=18pt,top=13pt,bottom=13pt,before skip=6pt,after skip=20pt,
  before upper={\poppill{Sommaire}{poppurple}{white}\par\vspace{9pt}\popsemi\fontsize{10.6}{14.5}\selectfont\color{popink}}}

% Signature de fin de cours — poussée tout en bas de la dernière page
\newcommand{\findecours}{%
  \par\vfill\nopagebreak
  \begin{center}\popsquiggle{poporange}\end{center}\vspace{4pt}
  \signatureonbuch
}
\AtBeginDocument{\def\llbracket{\mathopen{[\mkern-3.5mu[}}\def\rrbracket{\mathclose{]\mkern-3.5mu]}}} % intervalles d'entiers (glyphe ⟦ absent de la police)
% Glyphes absents de Fira Math : redéfinis à partir de glyphes disponibles
\AtBeginDocument{%
  \def\setminus{\mathbin{\backslash}}\let\smallsetminus\setminus
  \def\vdots{\mathord{\vbox{\baselineskip3.2pt\lineskiplimit0pt\kern4pt\hbox{.}\hbox{.}\hbox{.}}}}%
  \def\ddots{\mathinner{\mkern1mu\raise6.5pt\hbox{.}\mkern2mu\raise3.6pt\hbox{.}\mkern2mu\raise0.7pt\hbox{.}\mkern1mu}}%
  \def\triangle{\mathord{\scalebox{0.9}{$\symup{\Delta}$}}}%
  \def\nabla{\mathord{\raisebox{\depth}{\scalebox{1}[-1]{$\symup{\Delta}$}}}}%
  \def\checkmark{\popcheck{popdark}}%
  \def\star{\mathbin{\ast}}\def\bigstar{\mathord{\ast}}%
  \def\diamond{\mathbin{\scalebox{0.6}{$\Diamond$}}}%
  \def\bigcup{\mathop{\mathchoice{\raisebox{-0.3em}{\scalebox{1.7}{$\cup$}}}{\scalebox{1.25}{$\cup$}}{\cup}{\cup}}\displaylimits}%
  \def\bigcap{\mathop{\mathchoice{\raisebox{-0.3em}{\scalebox{1.7}{$\cap$}}}{\scalebox{1.25}{$\cap$}}{\cap}{\cap}}\displaylimits}%
  \def\lesseqqgtr{\mathrel{\lessgtr}}\def\bigoplus{\mathop{\oplus}\displaylimits}%
}
\providecommand{\ssi}{si et seulement si} % abréviation fréquente
\AtBeginDocument{\providecommand{\wideparen}{\overparen}} % arc de cercle

%% FIGFIX-BEGIN v5 — correctifs globaux des figures (docs/onbuch-generation/tools/figfix)
\usepackage{adjustbox}\usepackage{etoolbox}\tcbuselibrary{raster}
% toute figure TikZ/pgfplots plus large que la ligne est réduite à la largeur disponible
\BeforeBeginEnvironment{tikzpicture}{\begin{adjustbox}{max width=\linewidth}}
\AfterEndEnvironment{tikzpicture}{\end{adjustbox}}
% légendes pgfplots lisibles par-dessus les courbes
\pgfplotsset{every axis/.append style={legend style={fill=white,fill opacity=0.92,text opacity=1,draw=black!15,inner sep=3pt}}}
% étiquettes d'arêtes (syntaxe edge["..."]) avec fond blanc
\tikzset{every edge quotes/.append style={fill=white,inner sep=1.2pt}}
%% FIGFIX-END
\begin{document}

\pagegarde{Applications linéaires et affines}{Mathématiques}{3ème}{Mathématiques – classe de 3ème}{N16}{4 séances (fiche de progression harmonisée)}{Cette leçon structure l'étude des applications affines : calcul d'images et d'antécédents, sens de variation, représentation graphique et cas des fonctions définies par intervalles. Elle prépare à la modélisation rigoureuse de situations concrètes.
\vspace{4pt}
\begin{multicols}{2}\small
\begin{itemize}
\item À la fin de cette leçon, je sais distinguer une application linéaire d'une application affine à sa définition et à son graphe.
\item À la fin de cette leçon, je sais calculer l'image d'un nombre et déterminer l'antécédent(s) d'un nombre par une application affine.
\item À la fin de cette leçon, je sais déterminer le sens de variation d'une application affine selon le signe de son coefficient directeur.
\item À la fin de cette leçon, je sais représenter graphiquement une application affine dans un repère orthonormé et lire des valeurs sur le graphique.
\item À la fin de cette leçon, je sais étudier une application affine définie par intervalles (calculs, variations, graphique).
\item À la fin de cette leçon, je sais modéliser une situation concrète (tarification, facturation) par une application affine ou linéaire.
\end{itemize}\end{multicols}}

\begin{sommaire}
\begin{multicols}{2}
\begin{enumerate}
\item 1. Définitions et vocabulaire : Linéaire vs Affine
\item 2. Images, antécédents et résolution d'équations
\item 3. Sens de variation et coefficient directeur
\item 4. Représentation graphique dans un repère orthonormé
\item 5. Applications affines par intervalles (Fonctions par morceaux)
\item 6. Modélisation, problèmes de synthèse et ouverture Terminale
\item Activité d'intégration
\item Méthodes et exercices résolus
\item Exercices
\item Corrigés des exercices
\item Fiche bilan
\end{enumerate}
\end{multicols}
\end{sommaire}

\begin{objectifs}
\begin{itemize}
\item À la fin de cette leçon, je sais distinguer une application linéaire d'une application affine à sa définition et à son graphe.
\item À la fin de cette leçon, je sais calculer l'image d'un nombre et déterminer l'antécédent(s) d'un nombre par une application affine.
\item À la fin de cette leçon, je sais déterminer le sens de variation d'une application affine selon le signe de son coefficient directeur.
\item À la fin de cette leçon, je sais représenter graphiquement une application affine dans un repère orthonormé et lire des valeurs sur le graphique.
\item À la fin de cette leçon, je sais étudier une application affine définie par intervalles (calculs, variations, graphique).
\item À la fin de cette leçon, je sais modéliser une situation concrète (tarification, facturation) par une application affine ou linéaire.
\item À la fin de cette leçon, je sais rédiger une démonstration ou une justification structurée en utilisant le vocabulaire mathématique précis.
\item À la fin de cette leçon, je sais résoudre un problème de synthèse mobilisant l'algèbre et la géométrie (intersection de droites, optimisation simple).
\end{itemize}
\end{objectifs}

\begin{prerequis}
\begin{itemize}
\item Notion de proportionnalité et d'application linéaire (classe de 4ème).
\item Résolution d'équations du premier degré à une inconnue (ax + b = c).
\item Repérage dans le plan : coordonnées d'un point, placement de points.
\item Calcul littéral : développement, factorisation simple, ordre de priorité des opérations.
\item Notion de fonction (ensemble de départ, ensemble d'arrivée, image, antécédent) vue en début de 3ème.
\end{itemize}
\end{prerequis}

\newpage

\coursec{Définitions et vocabulaire : Linéaire vs Affine}

\courssub{Situation d'accroche : Deux courses, deux mathématiques}

Imagine la scène : nous sommes un lundi matin à Yaoundé, au carrefour \emph{Warda}. Monsieur Kamga hèle un taxi pour aller au travail. Le compteur affiche déjà \textbf{500 FCFA} avant même que la voiture ne démarre : c'est la \emph{prise en charge}. Ensuite, chaque kilomètre parcouru ajoute \textbf{100 FCFA}. S'il fait 12 km, il paiera $500 + 100 \times 12 = 1\,700$ FCFA.

Pendant ce temps, sa fille Assana achète du crédit téléphonique chez un revendeur \emph{MTN} du quartier. Elle donne \textbf{1\,000 FCFA} et reçoit exactement \textbf{500 Mo} d'internet. Pas de frais de dossier, pas de seuil : 2\,000 FCFA donnent 1\,000 Mo, 500 FCFA donnent 250 Mo. Le rapport est constant.

\begin{attention}[Problématique de la leçon]
Deux situations de la vie courante, deux règles de calcul différentes :
\begin{itemize}
    \item Taxi : $f(x) = 100x + 500$ (il y a un \textbf{départ non nul}).
    \item Crédit téléphone : $g(x) = 0,5x$ (le rapport est \textbf{constant}, passage par l'origine).
\end{itemize}
Comment les distinguer mathématiquement ? Quels sont leurs noms, leurs propriétés, leurs représentations graphiques ? C'est tout l'objet de cette première section.
\end{attention}

\courssub{Rappel : l'application linéaire (la proportionnalité)}

En classe de 4ème, vous avez étudié les situations de \textbf{proportionnalité}. Deux grandeurs $x$ et $y$ sont proportionnelles s'il existe un nombre $a$ tel que $y = a \times x$ pour toutes les valeurs de $x$.

\begin{definition}[Application linéaire]
Une application $f$ de $\mathbb{R}$ dans $\mathbb{R}$ est dite \textbf{linéaire} s'il existe un réel $a$ tel que, pour tout $x \in \mathbb{R}$ :
\[
f(x) = ax
\]
Le réel $a$ s'appelle le \textbf{coefficient de proportionnalité} (ou \textbf{coefficient directeur}). L'ensemble de définition est $D_f = \mathbb{R}$.
\end{definition}

\begin{propriete}[Caractérisation d'une application linéaire]
Soit $f$ une application de $\mathbb{R}$ dans $\mathbb{R}$. Dire que $f$ est linéaire signifie :
\begin{enumerate}
    \item $f(0) = 0$ : la représentation graphique de $f$ est une droite qui \textbf{passe par l'origine} du repère.
    \item Le quotient $\dfrac{f(x)}{x}$ est \textbf{constant} (égal à $a$) pour tout $x \neq 0$ : c'est le critère de proportionnalité.
\end{enumerate}
\emph{En 3ème, on retient ces deux tests : le test numérique $f(0)=0$ et le test du quotient constant.}
\end{propriete}

\begin{exemplebox}[Achat au kilogramme : le modèle linéaire pur]
Au marché \emph{Mfoundi}, les tomates coûtent \textbf{800 FCFA le kg}.
\begin{itemize}
    \item $x$ = masse achetée (en kg), $y$ = prix à payer (en FCFA).
    \item Relation : $y = 800x$.
    \item Tableau de valeurs :
\end{itemize}
\begin{center}
\begin{tabularx}{\linewidth}{|c|*{5}{>{\centering\arraybackslash}X|}}\hline
$x$ (kg) & 0 & 0,5 & 1 & 2 & 3 \\ \hline
$y = 800x$ (FCFA) & 0 & 400 & 800 & 1\,600 & 2\,400 \\ \hline
\end{tabularx}
\end{center}
\begin{itemize}
    \item Vérification : $\dfrac{y}{x} = 800$ est constant pour $x \neq 0$, et $f(0)=0$. C'est une application \textbf{linéaire}.
\end{itemize}
\end{exemplebox}

\begin{popfigure}
\begin{tikzpicture}[scale=0.9, >=Stealth]
% Axes
\draw[->] (-0.5,0) -- (4,0) node[right] {$x$ (kg)};
\draw[->] (0,-0.5) -- (0,3) node[above] {$y$ (centaines de FCFA)};
% Graduations
\foreach \x in {1,2,3} \draw (\x,0.05) -- (\x,-0.05) node[below] {\x};
\foreach \y/\lab in {0.8/8, 1.6/16, 2.4/24} \draw (0.05,\y) -- (-0.05,\y) node[left] {\lab};
% Points
\foreach \x/\y in {0/0, 0.5/0.4, 1/0.8, 2/1.6, 3/2.4} {
    \fill[popblue] (\x,\y) circle (2.5pt);
}
% Droite
\draw[thick, popblue] (0,0) -- (3.5, 2.8);
\node[popblue, anchor=south west] at (2.6, 2.3) {$y = 800x$};
% Origine
\node[below left, popmuted] at (0,0) {$O$};
\end{tikzpicture}
\legende{Représentation graphique d'une application linéaire : la droite passe par l'origine $O(0\,;0)$.}
\end{popfigure}

\courssub{Définition : l'application affine}

Dans la vie réelle, il y a souvent un \textbf{coût de départ}, un \textbf{forfait}, une \emph{prise en charge} : une somme à payer \emph{avant} même de consommer la première unité. C'est le cas du taxi de Monsieur Kamga.

\begin{definition}[Application affine]
Une application $f$ de $\mathbb{R}$ dans $\mathbb{R}$ est dite \textbf{affine} s'il existe deux réels $a$ et $b$ tels que, pour tout $x \in \mathbb{R}$ :
\[
f(x) = ax + b
\]
\begin{itemize}
    \item $a$ est le \textbf{coefficient directeur} (taux de variation, prix unitaire, pente).
    \item $b$ est l'\textbf{ordonnée à l'origine} (valeur initiale, coût fixe, forfait, prise en charge).
    \item L'ensemble de définition est $D_f = \mathbb{R}$.
\end{itemize}
Si $b = 0$, l'application affine redevient une application \textbf{linéaire}. On dit que la linéaire est un \emph{cas particulier} de l'affine.
\end{definition}

\begin{propriete}[Lien entre affine et linéaire : translation verticale]
Soit $f(x) = ax + b$ une application affine et $g(x) = ax$ l'application linéaire associée (même coefficient $a$).
Alors, pour tout $x \in \mathbb{R}$ :
\[
f(x) = g(x) + b
\]
\textbf{Interprétation graphique :} la courbe de $f$ s'obtient en \textbf{translatant} la courbe de $g$ de $b$ unités verticalement (vers le haut si $b>0$, vers le bas si $b<0$). Les deux droites ont le même coefficient directeur : elles sont \textbf{parallèles}.
\end{propriete}

\begin{exemplebox}[Le taxi de Monsieur Kamga : modèle affine]
Reprenons la situation d'accroche :
\begin{itemize}
    \item $x$ = distance parcourue (en km) ; $y$ = prix de la course (en FCFA).
    \item Prise en charge (fixe) = 500 FCFA $\rightarrow b = 500$.
    \item Prix au km (variable) = 100 FCFA/km $\rightarrow a = 100$.
    \item Fonction : $f(x) = 100x + 500$.
\end{itemize}
\begin{center}
\begin{tabularx}{\linewidth}{|c|*{5}{>{\centering\arraybackslash}X|}}\hline
$x$ (km) & 0 & 5 & 10 & 12 & 20 \\ \hline
$f(x) = 100x+500$ (FCFA) & 500 & 1\,000 & 1\,500 & 1\,700 & 2\,500 \\ \hline
\end{tabularx}
\end{center}
\begin{itemize}
    \item $f(0) = 500 \neq 0$ : \textbf{pas de proportionnalité}, c'est une application \textbf{affine}.
    \item Le quotient $\dfrac{f(x)}{x}$ n'est pas constant : $\dfrac{1\,000}{5}=200$, $\dfrac{1\,500}{10}=150$, $\dfrac{1\,700}{12}\approx 141,7$.
\end{itemize}
\end{exemplebox}

\begin{popfigure}
\begin{tikzpicture}[scale=0.62, >=Stealth]
% Axes
\draw[->] (-0.5,0) -- (13.5,0) node[right] {$x$ (km)};
\draw[->] (0,-0.5) -- (0,3) node[above] {$y$ (centaines de FCFA)};
% Graduations
\foreach \x in {2,4,6,8,10,12} \draw (\x,0.05) -- (\x,-0.05) node[below] {\x};
\foreach \y/\lab in {0.5/5, 1/10, 1.5/15, 2/20, 2.5/25} \draw (0.05,\y) -- (-0.05,\y) node[left] {\lab};
% Linéaire g(x)=100x (échelle /100)
\draw[popgreen!70!popdark, dashed, thick] (0,0) -- (12,1.2) node[below right, popgreen!70!popdark] {$g(x)=100x$};
% Affine f(x)=100x+500
\draw[poporange, thick] (0,0.5) -- (12,1.7) node[above right, poporange] {$f(x)=100x+500$};
% Points sur l'affine
\foreach \x/\y in {0/0.5, 5/1, 10/1.5, 12/1.7} {
    \fill[poporange] (\x,\y) circle (2.5pt);
}
% Points sur la linéaire
\foreach \x/\y in {5/0.5, 10/1, 12/1.2} {
    \fill[popgreen!70!popdark] (\x,\y) circle (2pt);
}
% Flèche de translation
\draw[<->, poppurple, thick] (0,0) -- (0,0.5) node[midway, left, poppurple] {$b=500$};
\end{tikzpicture}
\legende{Comparaison graphique : $g(x)=100x$ (linéaire, pointillés verts) et $f(x)=100x+500$ (affine, orange). Les deux droites sont parallèles ; la seconde est translatée de 500 FCFA vers le haut.}
\end{popfigure}

\courssub{Vocabulaire précis et cas particuliers}

Il est indispensable de maîtriser le vocabulaire officiel pour réussir les exercices du BEPC et rédiger correctement.

\begin{aretenir}[Vocabulaire de l'application affine $f(x) = ax + b$]
\begin{itemize}
    \item \cle{Coefficient directeur} : le nombre $a$. Il indique de combien $f(x)$ varie quand $x$ augmente de 1. C'est la \emph{pente} de la droite.
    \item \cle{Ordonnée à l'origine} : le nombre $b$. C'est la valeur de $f(0)$, l'image de 0. Graphiquement, c'est l'ordonnée du point d'intersection de la droite avec l'axe des ordonnées : le point $I(0\,;b)$.
    \item \cle{Ensemble de définition} $D_f$ : ensemble des valeurs que $x$ peut prendre. Pour les applications affines étudiées en 3ème, $D_f = \mathbb{R}$ (sauf mention contraire de l'énoncé).
    \item \cle{Image} : pour une valeur $x_0$ donnée, son image est $f(x_0) = ax_0 + b$.
    \item \cle{Antécédent} : pour une valeur $y_0$ donnée, un antécédent de $y_0$ est une solution de l'équation $ax + b = y_0$.
\end{itemize}
\end{aretenir}

\begin{definition}[Cas particuliers usuels]
\begin{enumerate}
    \item \textbf{Application constante} : $a = 0$. Alors $f(x) = b$ pour tout $x$. La courbe est une droite \emph{horizontale} (parallèle à l'axe des abscisses). Exemple : un forfait illimité à 5\,000 FCFA/mois, quelle que soit la consommation.
    \item \textbf{Application linéaire} : $b = 0$. Alors $f(x) = ax$. La courbe passe par l'origine.
    \item \textbf{Application identité} : $a = 1$ et $b = 0$. Alors $f(x) = x$ : l'image de tout nombre est ce nombre lui-même. Sa courbe est la droite qui fait un angle de 45° avec les axes (première bissectrice).
\end{enumerate}
\end{definition}

\begin{methode}[Identifier le type d'application : linéaire ou affine ?]
Devant une application $f$ donnée par une formule, un tableau ou un énoncé :
\begin{enumerate}
    \item \textbf{Calculer $f(0)$} (ou lire la valeur pour $x=0$ dans le tableau).
    \item \textbf{Si $f(0) = 0$} : vérifier que le quotient $\dfrac{f(x)}{x}$ est constant. Si oui, c'est une application \textbf{linéaire} de coefficient $a = \dfrac{f(x)}{x}$ pour tout $x \neq 0$.
    \item \textbf{Si $f(0) \neq 0$} : c'est une application \textbf{affine} d'ordonnée à l'origine $b = f(0)$. Le coefficient directeur se calcule par $a = \dfrac{f(x_1) - f(x_2)}{x_1 - x_2}$ pour deux valeurs distinctes $x_1$ et $x_2$.
\end{enumerate}
\end{methode}

\begin{attention}[Piège classique : « proportionnel » ou « affine » ?]
\begin{itemize}
    \item \textbf{Erreur fréquente :} « Le prix est proportionnel à la distance car plus on roule, plus on paie. »
    \item \textbf{Correction :} la proportionnalité exige \textbf{le passage par l'origine} ($f(0)=0$) et un \textbf{quotient constant}. Dès qu'il y a un forfait, une prise en charge, un abonnement ($b \neq 0$), ce n'est \textbf{plus} de la proportionnalité : c'est une relation \textbf{affine}.
    \item \textbf{Astuce :} pose-toi toujours la question : « Si je consomme 0 (0 km, 0 kg, 0 Mo), est-ce que je paie 0 ? » Si la réponse est \emph{non}, l'application est affine (non linéaire) ; si la réponse est \emph{oui}, elle peut être linéaire (vérifier alors le quotient).
\end{itemize}
\end{attention}

\courssub{Exemple chiffré complet : identification et paramètres}

\begin{exoresolu}[Analyse de $f(x) = 3x - 5$]
\textbf{Énoncé :} On considère l'application $f$ définie sur $\mathbb{R}$ par $f(x) = 3x - 5$.
\begin{enumerate}
    \item Préciser la nature de $f$ (linéaire ou affine). Justifier.
    \item Identifier le coefficient directeur $a$ et l'ordonnée à l'origine $b$.
    \item Calculer $f(0)$, $f(2)$ et déterminer l'antécédent de $4$ par $f$.
\end{enumerate}
\tcblower
\textbf{Solution détaillée :}
\begin{enumerate}
    \item \textbf{Nature :} calculons $f(0) = 3 \times 0 - 5 = -5$. Comme $f(0) \neq 0$, $f$ n'est pas linéaire. C'est une application \textbf{affine}.
    \item \textbf{Paramètres :} la forme générale est $f(x) = ax + b$. Par identification :
        \begin{itemize}
            \item $a = 3$ (coefficient directeur) : si $x$ augmente de 1, alors $f(x)$ augmente de 3 ;
            \item $b = -5$ (ordonnée à l'origine) : le point $I(0\,;-5)$ appartient à la courbe.
        \end{itemize}
    \item \textbf{Calculs :}
        \begin{itemize}
            \item $f(0) = -5$ (déjà calculé) ;
            \item $f(2) = 3 \times 2 - 5 = 6 - 5 = 1$. L'image de 2 est 1 ;
            \item Antécédent de 4 : on résout $f(x) = 4$ :
            \[
            3x - 5 = 4 \iff 3x = 4 + 5 \iff 3x = 9 \iff x = 3.
            \]
            L'antécédent de 4 est 3. Vérification : $f(3) = 3 \times 3 - 5 = 9 - 5 = 4$. \checkmark
        \end{itemize}
\end{enumerate}
\end{exoresolu}

\courssub{Schéma fléché : ensemble de départ et ensemble d'arrivée}

Pour visualiser la correspondance entre les nombres, on utilise un \textbf{schéma fléché}. C'est un outil très utile pour comprendre les notions d'image et d'antécédent : chaque nombre $x$ de l'ensemble de départ est relié par une flèche à son image $f(x)$ dans l'ensemble d'arrivée.

\begin{popfigure}
\begin{tikzpicture}[>=Stealth, scale=1]
% Ensemble de départ
\draw[popblue, thick, fill=popblueL] (0,0) ellipse (1.1 and 2.4);
\node[popblue, above] at (0,2.4) {Départ : $D_f=\mathbb{R}$};
\foreach \y/\lab in {1.5/3, 0.75/2, 0/0, -0.75/-1, -1.5/-2} {
    \fill[popdark] (0,\y) circle (1.5pt) node[left=2pt, popdark] {$\lab$};
}
% Ensemble d'arrivée
\draw[poporange, thick, fill=poporangeL] (6,0) ellipse (1.1 and 2.4);
\node[poporange, above] at (6,2.4) {Arrivée : $\mathbb{R}$};
\foreach \y/\lab in {1.5/4, 0.75/1, 0/-5, -0.75/-8, -1.5/-11} {
    \fill[popdark] (6,\y) circle (1.5pt) node[right=2pt, popdark] {$\lab$};
}
% Flèches de correspondance
\foreach \y in {1.5, 0.75, 0, -0.75, -1.5} {
    \draw[->, poppurple, thick] (0.15,\y) -- (5.85,\y);
}
% Nom de l'application
\node[poppurple, above] at (3,1.7) {$f(x)=3x-5$};
\end{tikzpicture}
\legende{Schéma fléché pour $f(x)=3x-5$. Chaque antécédent $x$ a une unique image : $3 \mapsto 4$, $2 \mapsto 1$, $0 \mapsto -5$, $-1 \mapsto -8$, $-2 \mapsto -11$.}
\end{popfigure}

\begin{attention}[Image et antécédent : ne pas confondre les rôles]
\begin{itemize}
    \item Chercher l'\textbf{image} de $x_0$, c'est \emph{calculer} : on remplace $x$ par $x_0$ dans la formule. Exemple : l'image de $2$ par $f(x)=3x-5$ est $f(2)=1$.
    \item Chercher l'\textbf{antécédent} de $y_0$, c'est \emph{résoudre une équation} : on cherche $x$ tel que $f(x)=y_0$. Exemple : l'antécédent de $4$ est la solution de $3x-5=4$, soit $x=3$.
    \item Sens de lecture de la phrase « $f(2)=1$ » : \textbf{1 est l'image de 2} et \textbf{2 est un antécédent de 1}.
\end{itemize}
\end{attention}

\courssub{Tableau de valeurs comparé : linéaire vs affine}

Le tableau de valeurs est l'outil privilégié en 3ème pour détecter la nature d'une application à partir de données numériques.

\begin{exemplebox}[Comparaison systématique : $f(x)=2x$ et $g(x)=2x+3$]
On compare la linéaire $f(x)=2x$ et l'affine $g(x)=2x+3$ (même coefficient directeur $a=2$).
\begin{center}
\begin{tabularx}{\linewidth}{|c|*{6}{>{\centering\arraybackslash}X|}}\hline
$x$ & $-2$ & $-1$ & $0$ & $1$ & $2$ & $3$ \\ \hline
$f(x) = 2x$ & $-4$ & $-2$ & \textbf{0} & $2$ & $4$ & $6$ \\ \hline
$g(x) = 2x+3$ & $-1$ & $1$ & \textbf{3} & $5$ & $7$ & $9$ \\ \hline
$g(x) - f(x)$ & $3$ & $3$ & \textbf{3} & $3$ & $3$ & $3$ \\ \hline
$\dfrac{f(x)}{x}$ ($x\neq0$) & $2$ & $2$ & -- & $2$ & $2$ & $2$ \\ \hline
$\dfrac{g(x)}{x}$ ($x\neq0$) & $0,5$ & $-1$ & -- & $5$ & $3,5$ & $3$ \\ \hline
\end{tabularx}
\end{center}
\textbf{Observations clés :}
\begin{itemize}
    \item La différence $g(x) - f(x)$ est \textbf{constante} (égale à $b=3$) : c'est la translation verticale.
    \item Le quotient $\dfrac{f(x)}{x}$ est constant (égal à 2) : proportionnalité pour $f$.
    \item Le quotient $\dfrac{g(x)}{x}$ \textbf{n'est pas constant} : pas de proportionnalité pour $g$.
    \item $f(0)=0$ tandis que $g(0)=3=b$.
\end{itemize}
\end{exemplebox}

\begin{popfigure}
\begin{tikzpicture}
\begin{axis}[
    width=\linewidth, height=8cm,
    axis lines=middle,
    xlabel=$x$, ylabel=$y$,
    xmin=-3, xmax=3.5,
    ymin=-5, ymax=10,
    xtick={-2,-1,0,1,2,3},
    ytick={-4,-2,0,2,4,6,8},
    grid=major, grid style={popline},
    legend style={at={(0.02,0.98)}, anchor=north west, fill=popcream, draw=popmuted},
    clip=false
]
% Linéaire f(x)=2x
\addplot[popcurve, popblue, thick, domain=-2.4:3, samples=50] {2*x};
\addlegendentry{$f(x)=2x$ (linéaire)};
% Affine g(x)=2x+3
\addplot[popcurve, poporange, thick, domain=-3:3, samples=50] {2*x+3};
\addlegendentry{$g(x)=2x+3$ (affine)};
% Points sur la linéaire
\addplot[only marks, mark=*, popblue] coordinates {(-2,-4) (-1,-2) (0,0) (1,2) (2,4) (3,6)};
% Points sur l'affine
\addplot[only marks, mark=square*, poporange] coordinates {(-2,-1) (-1,1) (0,3) (1,5) (2,7) (3,9)};
% Flèches de translation verticale
\draw[<->, poppurple, dashed, thick] (axis cs:-2,-4) -- (axis cs:-2,-1);
\draw[<->, poppurple, dashed, thick] (axis cs:-1,-2) -- (axis cs:-1,1);
\draw[<->, poppurple, dashed, thick] (axis cs:0,0) -- (axis cs:0,3);
\draw[<->, poppurple, dashed, thick] (axis cs:1,2) -- (axis cs:1,5);
\draw[<->, poppurple, dashed, thick] (axis cs:2,4) -- (axis cs:2,7);
\node[poppurple, anchor=west] at (axis cs:0.15,1.5) {$b=3$};
\end{axis}
\end{tikzpicture}
\legende{Superposition des courbes $y=2x$ (bleu) et $y=2x+3$ (orange). Même pente ($a=2$), décalage vertical de 3 unités : les droites sont parallèles. Les flèches violettes matérialisent la translation.}
\end{popfigure}

\courssub{Le saviez-vous ? Lecture économique des paramètres}

\begin{savaistu}[Coût fixe et coût unitaire : le langage des entreprises]
En économie, en gestion et dans les offres des opérateurs télécoms (\emph{MTN}, \emph{Orange}, \emph{Camtel}), on retrouve exactement nos paramètres $a$ et $b$ sous d'autres noms :
\begin{itemize}
    \item \textbf{$b$ = coût fixe, frais fixes, abonnement, prise en charge.} C'est ce qu'il faut payer \textbf{avant de consommer la première unité}. Exemples : location du local, abonnement internet mensuel, prise en charge du taxi.
    \item \textbf{$a$ = coût variable unitaire, prix unitaire.} C'est le coût \textbf{supplémentaire} pour consommer ou produire \textbf{une unité de plus}. Exemples : prix du kilomètre en taxi, prix du Mo de données, prix du kilogramme de tomates.
\end{itemize}
Le coût total pour $x$ unités est donc une application affine : $C(x) = ax + b$.
\begin{itemize}
    \item Si $b=0$ (pas de frais fixes), le coût est \textbf{proportionnel} à la quantité (linéaire) : c'est le cas de l'achat-revente simple, comme le revendeur de crédit téléphone.
    \item Le \textbf{seuil de rentabilité} d'un petit commerce correspond à l'antécédent de 0 pour la fonction \emph{Bénéfice} = \emph{Recette} $-$ \emph{Coût}. On le trouve en résolvant une équation du type $px - (ax+b) = 0$.
\end{itemize}
\textbf{En résumé :} maîtriser l'application affine $f(x)=ax+b$, c'est déjà savoir lire un budget, une facture ou un business plan !
\end{savaistu}

\begin{exercice}[Entraînement immédiat : identification]
Pour chacune des situations ci-dessous, dire si l'application qui associe $y$ à $x$ est linéaire ou affine. Préciser $a$ et $b$ dans chaque cas.
\begin{enumerate}
    \item Un abonnement \emph{Camtel} Fibre coûte 15\,000 FCFA/mois, plus 500 FCFA par Go supplémentaire au-delà du forfait. $x$ = nombre de Go supplémentaires, $y$ = facture totale (FCFA).
    \item Un maçon facture 3\,000 FCFA par m$^2$ de carrelage posé, fournitures comprises. $x$ = surface (m$^2$), $y$ = facture (FCFA).
    \item La température d'une pièce augmente de 2 °C par heure à partir d'une température initiale de 18 °C. $x$ = temps (h), $y$ = température (°C).
    \item $f(x) = -4x$.
    \item $g(x) = 7$.
\end{enumerate}
\end{exercice}

\begin{corrige}[Exercice n°1]
\begin{enumerate}
    \item \textbf{Affine.} $y = 500x + 15\,000$, avec $a=500$ (prix du Go supplémentaire) et $b=15\,000$ (abonnement fixe). On a $f(0)=15\,000 \neq 0$.
    \item \textbf{Linéaire.} $y = 3\,000x$, avec $a=3\,000$ et $b=0$. On a $f(0)=0$ : c'est une situation de proportionnalité.
    \item \textbf{Affine.} $y = 2x + 18$, avec $a=2$ (vitesse de chauffe en °C/h) et $b=18$ (température initiale). On a $f(0)=18 \neq 0$.
    \item \textbf{Linéaire.} $f(x) = -4x$, avec $a=-4$ et $b=0$. On a $f(0)=0$. (Le coefficient directeur est négatif : les images diminuent quand $x$ augmente.)
    \item \textbf{Affine constante.} $g(x) = 0 \times x + 7$, avec $a=0$ et $b=7$. On a $g(0)=7 \neq 0$ ; la représentation graphique est une droite horizontale.
\end{enumerate}
\end{corrige}

\courssub{Synthèse de la section : ce qu'il faut retenir absolument}

\begin{aretenir}[Bilan : linéaire vs affine]
\begin{center}
\begin{tabularx}{\linewidth}{|>{\bfseries}l|X|X|}\hline
\textbf{Critère} & \textbf{Application linéaire} & \textbf{Application affine} \\ \hline
Formule & $f(x) = ax$ & $f(x) = ax + b$ \\ \hline
Paramètres & $a \in \mathbb{R}$ (coefficient directeur) & $a \in \mathbb{R}$ (coefficient directeur), $b \in \mathbb{R}$ (ordonnée à l'origine) \\ \hline
Condition & $b = 0$ & $b$ quelconque \\ \hline
Test numérique & $f(0) = 0$ & $f(0) = b$ \\ \hline
Quotient $\dfrac{f(x)}{x}$ ($x\neq0$) & Constant (égal à $a$) & Non constant (sauf si $b=0$) \\ \hline
Graphique & Droite passant par l'origine $O(0\,;0)$ & Droite coupant l'axe des ordonnées au point $(0\,;b)$ \\ \hline
Lien entre les deux & Cas particulier de l'affine ($b=0$) & Cas général : translation verticale de la linéaire $g(x)=ax$ \\ \hline
Exemples concrets & Achat au kg, conversion d'unités, crédit téléphone & Taxi (prise en charge), forfait + dépassement, chauffage, coût total d'une entreprise \\ \hline
\end{tabularx}
\end{center}
\end{aretenir}

\begin{methode}[Réflexe « début d'exercice BEPC »]
Devant un énoncé parlant d'une relation entre deux grandeurs $x$ et $y$ :
\begin{enumerate}
    \item \textbf{Identifier les variables :} $x$ = grandeur indépendante (souvent temps, quantité, distance) ; $y$ = grandeur dépendante (prix, température, volume).
    \item \textbf{Chercher le « départ » :} quelle est la valeur de $y$ quand $x=0$ ? C'est $b$.
    \item \textbf{Chercher le « rythme » :} de combien $y$ varie-t-il quand $x$ augmente de 1 ? C'est $a$.
    \item \textbf{Écrire la formule :} $f(x) = ax + b$.
    \item \textbf{Conclure par une phrase :} « C'est une application affine (linéaire si $b=0$) de coefficient directeur $a$ et d'ordonnée à l'origine $b$. »
\end{enumerate}
\end{methode}

Cette première section pose les fondations. Dans la section suivante, nous apprendrons à \textbf{calculer des images et des antécédents} (résoudre $f(x)=k$) et à exploiter ces notions pour résoudre des problèmes concrets.

\coursec{Images, antécédents et résolution d'équations}

Dans la section précédente, nous avons défini les applications affines : une application $f$ est affine lorsqu'elle s'écrit $f(x) = ax + b$, où $a$ et $b$ sont deux réels fixés. Nous avons aussi distingué le cas particulier des applications linéaires ($b = 0$). Mais une fois qu'on a défini une telle machine mathématique, deux questions naturelles se posent :

\begin{itemize}
\item \textbf{Question directe :} si je connais un nombre $x$, quel est son image $f(x)$ ?
\item \textbf{Question inverse :} si je connais une image $k$, quel(s) nombre(s) $x$ a (ont) pour image $k$ ?
\end{itemize}

La première question est un simple calcul. La seconde, plus subtile, nous ramène à résoudre une \cle{équation du premier degré} — et c'est là tout l'intérêt de cette section : tu vas découvrir que résoudre $ax + b = k$ et chercher un antécédent, c'est \emph{exactement la même activité}, vue sous deux angles différents.

\courssub{Calcul d'une image : la question directe}

\begin{definition}[Image d'un nombre par une application affine]
Soit $f : x \mapsto ax + b$ une application affine et $x_0$ un nombre réel. L'\cle{image} de $x_0$ par $f$ est le nombre obtenu en remplaçant $x$ par $x_0$ dans l'expression de $f$ :
\[ f(x_0) = a \times x_0 + b. \]
On note aussi $f(x_0)$ et on lit « $f$ de $x_0$ ».
\end{definition}

L'intuition est celle d'une machine : on introduit un nombre $x$, la machine le multiplie par $a$, ajoute $b$, et sort le résultat $f(x)$. Le calcul d'image est donc un calcul \emph{direct}, sans aucune équation à résoudre. La seule vigilance requise concerne les \cle{priorités opératoires} : la multiplication par $a$ se fait \textbf{avant} l'addition de $b$, et il faut être très attentif aux signes lorsque $x_0$ ou $a$ est négatif.

\begin{exemplebox}[Calculs d'images avec $f(x) = -2x + 6$]
Soit $f$ l'application affine définie sur $\mathbb{R}$ par $f(x) = -2x + 6$.
\begin{itemize}
\item Image de $4$ : $f(4) = -2 \times 4 + 6 = -8 + 6 = -2$. L'image de $4$ est $-2$.
\item Image de $-3$ : $f(-3) = -2 \times (-3) + 6 = 6 + 6 = 12$. L'image de $-3$ est $12$.
\item Image de $0$ : $f(0) = -2 \times 0 + 6 = 6$. L'image de $0$ est $6$ (c'est l'ordonnée à l'origine $b$).
\end{itemize}
\end{exemplebox}

\begin{attention}[Pièges du calcul d'image]
\begin{itemize}
\item \textbf{Signe du produit :} pour $f(x) = -2x + 6$, on a $f(-3) = -2 \times (-3) + 6 = +6 + 6 = 12$, et non $-6 + 6 = 0$. Le produit de deux nombres négatifs est positif !
\item \textbf{Priorité :} on calcule d'abord $a \times x$, \emph{puis} on ajoute $b$. Ne jamais calculer $(x + b)$ d'abord.
\item \textbf{Parenthèses :} quand $x_0$ est négatif, écris-le entre parenthèses : $f(-3) = -2 \times (-3) + 6$. Cela évite 90 \% des erreurs.
\end{itemize}
\end{attention}

\courssub{Recherche d'antécédent : la question inverse}

\begin{definition}[Antécédent d'un nombre]
Soit $f : x \mapsto ax + b$ une application affine et $k$ un nombre réel. Un \cle{antécédent} de $k$ par $f$ est un nombre $x$ dont l'image est $k$, c'est-à-dire un nombre $x$ tel que :
\[ f(x) = k. \]
L'ensemble de tous les antécédents de $k$ se note $f^{-1}(\{k\})$.
\end{definition}

La différence avec l'image est fondamentale : pour une image, on part de $x$ et on calcule ; pour un antécédent, on part du \emph{résultat} $k$ et on cherche à \emph{remonter} vers $x$. C'est exactement ce que fait un chauffeur de taxi-moto qui connaît le prix de la course et veut retrouver la distance parcourue : il « inverse » le tarif.

Or, dire $f(x) = k$, c'est dire $ax + b = k$ : chercher un antécédent revient donc à \textbf{résoudre une équation du premier degré}.

\begin{methode}[Calcul d'un antécédent]
Pour trouver le ou les antécédents d'un nombre $k$ par $f(x) = ax + b$ :
\begin{enumerate}
\item \textbf{Écrire l'équation} $f(x) = k$, c'est-à-dire $ax + b = k$.
\item \textbf{Isoler $x$} : on retranche $b$ des deux membres, puis on divise par $a$ (si $a \neq 0$) :
\[ ax + b = k \iff ax = k - b \iff x = \frac{k - b}{a}. \]
\item \textbf{Conclure par une phrase} : « L'antécédent de $k$ par $f$ est $\dfrac{k-b}{a}$ » ou, le cas échéant, « $k$ n'a pas d'antécédent ».
\end{enumerate}
\end{methode}

L'arbre suivant résume visuellement la démarche de résolution :

\begin{popfigure}
\begin{tikzpicture}[
  node distance=1.1cm,
  etape/.style={rectangle, rounded corners=4pt, draw=popblue, fill=popblueL, thick, minimum width=4.2cm, minimum height=0.9cm, font=\small},
  action/.style={font=\small\itshape, text=popdark},
  fleche/.style={-{Stealth[length=3mm]}, thick, poporange}
]
\node[etape] (e1) {$ax + b = k$};
\node[etape, below=of e1] (e2) {$ax = k - b$};
\node[etape, below=of e2] (e3) {$x = \dfrac{k-b}{a}$ \quad (si $a \neq 0$)};
\draw[fleche] (e1) -- (e2) node[midway, right, action]{on retranche $b$ aux deux membres};
\draw[fleche] (e2) -- (e3) node[midway, right, action]{on divise les deux membres par $a$};
\end{tikzpicture}
\legende{Arbre de résolution de l'équation $ax + b = k$ : deux étapes suffisent pour isoler $x$.}
\end{popfigure}

\begin{exemplebox}[Recherche d'antécédents avec $f(x) = -2x + 6$]
Soit $f$ définie par $f(x) = -2x + 6$.
\begin{itemize}
\item \textbf{Antécédent de $0$ :} on résout $-2x + 6 = 0 \iff -2x = -6 \iff x = \dfrac{-6}{-2} = 3$. L'antécédent de $0$ est $3$. Vérification : $f(3) = -6 + 6 = 0$. \checkmark
\item \textbf{Antécédent de $10$ :} on résout $-2x + 6 = 10 \iff -2x = 4 \iff x = \dfrac{4}{-2} = -2$. L'antécédent de $10$ est $-2$. Vérification : $f(-2) = 4 + 6 = 10$. \checkmark
\end{itemize}
Remarque : on peut appliquer directement la formule $x = \dfrac{k-b}{a}$. Pour $k = 10$ : $x = \dfrac{10 - 6}{-2} = \dfrac{4}{-2} = -2$.
\end{exemplebox}

\begin{attention}[Image ou antécédent : ne pas confondre les questions]
« Calculer l'image de $4$ » et « calculer l'antécédent de $4$ » sont deux questions \textbf{différentes} qui donnent en général des réponses différentes. Avec $f(x) = -2x + 6$ :
\begin{itemize}
\item image de $4$ : $f(4) = -2$ (on remplace $x$ par $4$) ;
\item antécédent de $4$ : $-2x + 6 = 4 \iff x = 1$ (on résout une équation).
\end{itemize}
Lis toujours l'énoncé deux fois avant de te lancer !
\end{attention}

\courssub{Le cas particulier $a = 0$ : fonction constante}

La formule $x = \dfrac{k-b}{a}$ contient une division par $a$ : elle n'a de sens que si $a \neq 0$. Que se passe-t-il si $a = 0$ ? Alors $f(x) = b$ pour \emph{tout} réel $x$ : c'est une \cle{fonction constante}. Tous les nombres ont la même image $b$. Deux cas se présentent :

\begin{itemize}
\item Si $k = b$ : l'équation $f(x) = k$ devient $b = b$, toujours vraie. \textbf{Tout réel $x$ est antécédent} : $f^{-1}(\{k\}) = \mathbb{R}$.
\item Si $k \neq b$ : l'équation devient $b = k$, toujours fausse. \textbf{Aucun réel n'est antécédent} : $f^{-1}(\{k\}) = \varnothing$ (ensemble vide).
\end{itemize}

\begin{exemplebox}[Fonction constante $g(x) = 5$]
Soit $g$ définie sur $\mathbb{R}$ par $g(x) = 5$.
\begin{itemize}
\item Antécédents de $5$ : tout nombre $x$ vérifie $g(x) = 5$. Donc $g^{-1}(\{5\}) = \mathbb{R}$.
\item Antécédents de $3$ : aucun nombre n'a pour image $3$. Donc $g^{-1}(\{3\}) = \varnothing$.
\end{itemize}
\end{exemplebox}

\begin{attention}[Toujours discuter selon la valeur de $a$]
Devant l'équation $ax + b = k$, réflexe obligatoire : écris « \textbf{Si $a \neq 0$}... \textbf{Si $a = 0$}... ». Oublier le cas $a = 0$ et diviser par $a$ sans précaution est une erreur classique, sanctionnée au BEPC. La division par zéro est interdite : toute formule contenant $\dfrac{\ldots}{a}$ exige l'hypothèse $a \neq 0$.
\end{attention}

\courssub{Unicité de l'antécédent : une propriété fondamentale}

Pour une application affine non constante, la question inverse admet toujours une réponse, et une seule. C'est une propriété remarquable.

\begin{propriete}[Existence et unicité de l'antécédent]
Soit $f : x \mapsto ax + b$ une application affine avec $a \neq 0$. Alors tout nombre réel $k$ admet un \textbf{unique} antécédent par $f$, donné par :
\[ f^{-1}(\{k\}) = \left\{ \frac{k - b}{a} \right\}. \]
On dit que $f$ réalise une correspondance biunivoque (on dit aussi que $f$ est \cle{bijective}) entre $\mathbb{R}$ et $\mathbb{R}$.
\end{propriete}

\begin{experience}[Démonstration guidée de la propriété]
Fixons un réel $k$ quelconque et cherchons tous les $x$ tels que $f(x) = k$.

\textbf{Étape 1 — Existence (tout réel a au moins un antécédent).} Posons $x_0 = \dfrac{k-b}{a}$ (ce nombre existe car $a \neq 0$). Vérifions :
\[ f(x_0) = a \times \frac{k-b}{a} + b = (k - b) + b = k. \]
Donc $x_0$ est bien un antécédent de $k$ : il en existe au moins un.

\textbf{Étape 2 — Unicité (il n'y en a pas deux).} Supposons que deux nombres $x_1$ et $x_2$ soient tous deux antécédents de $k$ : $f(x_1) = k$ et $f(x_2) = k$. Alors :
\[ ax_1 + b = ax_2 + b \iff ax_1 = ax_2 \iff x_1 = x_2 \quad (\text{car } a \neq 0). \]
Les deux antécédents sont forcément égaux : il n'y en a donc qu'un seul.

\textbf{Conclusion.} Pour tout réel $k$, l'équation $f(x) = k$ admet exactement une solution : $x = \dfrac{k-b}{a}$. $\blacksquare$
\end{experience}

\begin{savaistu}[La notation $f^{-1}(k)$ : un abus de langage toléré]
Quand $a \neq 0$, l'antécédent de $k$ est unique : on écrit alors souvent $f^{-1}(k)$ \emph{sans accolades}, par exemple $f^{-1}(10) = -2$, pour désigner directement ce nombre unique. C'est un abus de langage commode et toléré, mais garde en tête que l'écriture rigoureuse est $f^{-1}(\{k\}) = \left\{\dfrac{k-b}{a}\right\}$ : $f^{-1}(\{k\})$ désigne un \emph{ensemble} de nombres, qui ici ne contient qu'un seul élément. Attention aussi : $f^{-1}$ ne signifie \textbf{pas} $\dfrac{1}{f}$ ! Ce n'est pas un inverse au sens de la division.
\end{savaistu}

\courssub{Lecture graphique : l'antécédent comme intersection}

Géométriquement, chercher l'antécédent de $k$ revient à chercher le point de la droite $y = ax + b$ dont l'ordonnée vaut $k$. C'est le point d'intersection de la droite représentative de $f$ avec la droite horizontale d'équation $y = k$.

\begin{popfigure}
\begin{tikzpicture}
\begin{axis}[
  width=0.85\linewidth, height=7cm,
  axis lines=middle,
  xlabel={$x$}, ylabel={$y$},
  xlabel style={anchor=west}, ylabel style={anchor=south},
  xmin=-4, xmax=6, ymin=-3, ymax=12,
  xtick={-3,-2,-1,1,2,3,4,5}, ytick={-2,2,4,6,8,10},
  grid=both, grid style={popline!40, very thin},
  tick label style={font=\scriptsize}
]
\addplot[popcurve, domain=-3.5:5.5, samples=50]{-2*x+6};
\addplot[poporange, thick, domain=-4:6, samples=2]{10};
\addplot[popgreen, thick, domain=-4:6, samples=2]{0};
\addplot[only marks, mark=*, mark size=2pt, popink] coordinates {(-2,10) (3,0)};
\node[popblue, anchor=west, font=\small] at (axis cs:3.4,7) {$y = -2x + 6$};
\node[poporange, anchor=south west, font=\small] at (axis cs:-3.9,10.2) {$y = 10$};
\node[popgreen, anchor=north west, font=\small] at (axis cs:-3.9,-0.4) {$y = 0$};
\node[popink, anchor=south east, font=\small] at (axis cs:-2,10) {$(-2\,;\,10)$};
\node[popink, anchor=south west, font=\small] at (axis cs:3,0) {$(3\,;\,0)$};
\draw[dashed, popmuted] (axis cs:-2,0) -- (axis cs:-2,10);
\draw[dashed, popmuted] (axis cs:0,10) -- (axis cs:-2,10);
\end{axis}
\end{tikzpicture}
\legende{Lecture graphique des antécédents pour $f(x) = -2x + 6$ : l'antécédent de $k$ est l'abscisse du point d'intersection de la droite avec la droite horizontale $y = k$. Ici, l'antécédent de $10$ est $-2$ et celui de $0$ est $3$.}
\end{popfigure}

Cette lecture confirme visuellement la propriété démontrée plus haut : une droite non horizontale ($a \neq 0$) coupe \emph{toujours} une droite horizontale en \emph{exactement un} point. En revanche, une fonction constante ($a = 0$) a pour représentation une droite horizontale : soit elle coïncide avec $y = k$ (infinité d'antécédents), soit elle ne la rencontre jamais (aucun antécédent).

\courssub{Tableau de valeurs : conjecturer avant de calculer}

Un tableau de valeurs permet de conjecturer rapidement images et antécédents, et de vérifier ensuite par le calcul. Voici le tableau de $f(x) = -2x + 6$ :

\begin{center}
\begin{tabular}{|c|c|c|c|c|c|c|}
\hline
\rowcolor{popblueL} $x$ & $-2$ & $-1$ & $0$ & $1$ & $3$ & $4$ \\
\hline
$f(x) = -2x + 6$ & $10$ & $8$ & $6$ & $4$ & $0$ & $-2$ \\
\hline
\end{tabular}
\end{center}

Ce tableau se lit dans les deux sens :
\begin{itemize}
\item \textbf{Sens « image »} (de haut en bas) : on part de la ligne des $x$ et on descend. Par exemple, l'image de $4$ est $-2$ : $4 \longrightarrow -2$.
\item \textbf{Sens « antécédent »} (de bas en haut) : on part de la ligne des $f(x)$ et on remonte. Par exemple, l'antécédent de $10$ est $-2$ : $-2 \longleftarrow 10$.
\end{itemize}

\begin{attention}[Limites du tableau de valeurs]
Un tableau ne donne que quelques valeurs isolées. Si l'antécédent cherché n'est pas un nombre entier (par exemple l'antécédent de $5$ est $x = \dfrac{5-6}{-2} = 0,5$), il n'apparaîtra pas dans un tableau à pas entier. Le tableau sert à \emph{conjecturer} ou \emph{vérifier}, jamais à prouver : seule la résolution de l'équation $ax + b = k$ donne la réponse exacte et complète.
\end{attention}

\courssub{Lien avec l'équation du premier degré : les zéros de la fonction}

Un cas particulier mérite une attention spéciale : la recherche de l'antécédent de $0$.

\begin{definition}[Zéro d'une application affine]
Soit $f : x \mapsto ax + b$ avec $a \neq 0$. Le \cle{zéro} (ou \cle{racine}) de $f$ est l'unique antécédent de $0$, c'est-à-dire la solution de l'équation :
\[ ax + b = 0 \iff x = -\frac{b}{a}. \]
Graphiquement, c'est l'abscisse du point où la droite représentative de $f$ coupe l'axe des abscisses.
\end{definition}

Résoudre l'équation du premier degré $ax + b = 0$ (notion vue au chapitre des équations) et chercher le zéro de l'application affine $f(x) = ax + b$ sont donc \textbf{deux formulations d'un même problème}. Cette unité des mathématiques est précieuse : chaque équation du premier degré peut être « vue » comme une question sur une fonction, et réciproquement.

\begin{exoresolu}[Images, antécédents et zéro d'une application affine]
On considère l'application affine $g$ définie sur $\mathbb{R}$ par $g(x) = 3x - 7$.
\begin{enumerate}
\item Calculer l'image de $5$ par $g$.
\item Déterminer l'antécédent de $8$ par $g$.
\item Déterminer le zéro de $g$. Interpréter graphiquement le résultat.
\item Le nombre $-10$ admet-il un antécédent par $g$ ? Si oui, lequel ?
\end{enumerate}
\tcblower
\textbf{Solution détaillée.}

\textbf{1.} Image de $5$ : on remplace $x$ par $5$ :
\[ g(5) = 3 \times 5 - 7 = 15 - 7 = 8. \]
L'image de $5$ par $g$ est $8$.

\textbf{2.} Antécédent de $8$ : on résout $g(x) = 8$ :
\[ 3x - 7 = 8 \iff 3x = 8 + 7 \iff 3x = 15 \iff x = \frac{15}{3} = 5. \]
L'antécédent de $8$ par $g$ est $5$. (Remarque : c'est cohérent avec la question 1 — image et antécédent sont deux notions réciproques.)

\textbf{3.} Zéro de $g$ : on résout $g(x) = 0$ :
\[ 3x - 7 = 0 \iff 3x = 7 \iff x = \frac{7}{3}. \]
Le zéro de $g$ est $\dfrac{7}{3}$. Interprétation graphique : la droite d'équation $y = 3x - 7$ coupe l'axe des abscisses au point $\left(\dfrac{7}{3}\,;\,0\right)$.

\textbf{4.} Comme $a = 3 \neq 0$, l'application $g$ est bijective : \emph{tout} réel admet un antécédent. On résout $g(x) = -10$ :
\[ 3x - 7 = -10 \iff 3x = -3 \iff x = -1. \]
L'antécédent de $-10$ par $g$ est $-1$. Vérification : $g(-1) = -3 - 7 = -10$. \checkmark
\end{exoresolu}

\begin{exercice}[À toi de jouer]
Soit $h$ l'application affine définie par $h(x) = -4x + 12$.
\begin{enumerate}
\item Calcule les images de $0$, de $2$ et de $-1$.
\item Détermine les antécédents de $12$, de $0$ et de $20$.
\item Quel est le zéro de $h$ ?
\item Construis un tableau de valeurs de $h$ pour $x \in \{-1\,;\,0\,;\,1\,;\,2\,;\,3\,;\,5\}$ et retrouve graphiquement tes réponses.
\end{enumerate}
\end{exercice}

\begin{corrige}[Exercice « À toi de jouer »]
\textbf{1.} $h(0) = 12$ ; $h(2) = -8 + 12 = 4$ ; $h(-1) = 4 + 12 = 16$.

\textbf{2.} Antécédent de $12$ : $-4x + 12 = 12 \iff -4x = 0 \iff x = 0$. Antécédent de $0$ : $-4x + 12 = 0 \iff x = 3$. Antécédent de $20$ : $-4x + 12 = 20 \iff -4x = 8 \iff x = -2$.

\textbf{3.} Le zéro de $h$ est $3$ (antécédent de $0$).

\textbf{4.} Tableau :
\begin{center}
\begin{tabular}{|c|c|c|c|c|c|c|}
\hline
\rowcolor{popblueL} $x$ & $-1$ & $0$ & $1$ & $2$ & $3$ & $5$ \\
\hline
$h(x)$ & $16$ & $12$ & $8$ & $4$ & $0$ & $-8$ \\
\hline
\end{tabular}
\end{center}
On y lit bien : image de $2$ égale à $4$, antécédent de $12$ égal à $0$, antécédent de $0$ égal à $3$.
\end{corrige}

\begin{aretenir}[L'essentiel de la section]
\begin{itemize}
\item \textbf{Image} de $x_0$ : on calcule $f(x_0) = ax_0 + b$ (calcul direct, attention aux signes et priorités).
\item \textbf{Antécédent} de $k$ : on résout $ax + b = k$. Si $a \neq 0$, unique solution $x = \dfrac{k-b}{a}$.
\item Si $a = 0$ (fonction constante $f(x) = b$) : $k = b \Rightarrow$ tout réel est antécédent ($\mathbb{R}$) ; $k \neq b \Rightarrow$ aucun antécédent ($\varnothing$).
\item Si $a \neq 0$, $f$ est \textbf{bijective} : tout réel a un unique antécédent, noté $f^{-1}(\{k\})$.
\item Le \textbf{zéro} de $f$ est l'antécédent de $0$ : $x = -\dfrac{b}{a}$ ; c'est l'abscisse du point où la droite coupe l'axe des abscisses.
\item Graphiquement, l'antécédent de $k$ est l'abscisse de l'intersection de la droite $y = ax + b$ avec la droite horizontale $y = k$.
\end{itemize}
\end{aretenir}

Maintenant que tu maîtrises le calcul des images et des antécédents, une question se pose naturellement : comment les images évoluent-elles lorsque $x$ augmente ? C'est l'objet de la section suivante, consacrée au \emph{sens de variation} et au rôle décisif du coefficient $a$.

\coursec{Sens de variation et coefficient directeur}

Après avoir maîtrisé le calcul d'images et d'antécédents (section 2), nous nous intéressons maintenant à la \emph{manière} dont une application affine évolue lorsque $x$ parcourt $\mathbb{R}$. Cette notion de variation est au cœur de la modélisation : elle permet de prévoir si une quantité augmente, diminue ou reste stable, et de comparer deux situations (coûts, revenus, distances, etc.).

\courssub{Rappel : notion de variation d'une fonction}

\begin{definition}[Fonction croissante, décroissante, constante]
Soit $f$ une fonction définie sur un intervalle $I$.
\begin{itemize}
\item $f$ est \textbf{croissante} sur $I$ si pour tout $x_1, x_2 \in I$, $x_1 < x_2 \implies f(x_1) \le f(x_2)$.
\item $f$ est \textbf{strictement croissante} sur $I$ si pour tout $x_1, x_2 \in I$, $x_1 < x_2 \implies f(x_1) < f(x_2)$.
\item $f$ est \textbf{décroissante} (resp. \textbf{strictement décroissante}) si l'inégalité est inversée.
\item $f$ est \textbf{constante} sur $I$ si pour tout $x_1, x_2 \in I$, $f(x_1) = f(x_2)$.
\end{itemize}
\end{definition}

\courssub{Théorème fondamental : variation d'une application affine}

\begin{propriete}[Variation d'une application affine]
Soit $f : \mathbb{R} \to \mathbb{R}$ définie par $f(x) = ax + b$ avec $a,b \in \mathbb{R}$.
\begin{itemize}
\item Si $a > 0$, alors $f$ est \textbf{strictement croissante} sur $\mathbb{R}$.
\item Si $a < 0$, alors $f$ est \textbf{strictement décroissante} sur $\mathbb{R}$.
\item Si $a = 0$, alors $f$ est \textbf{constante} sur $\mathbb{R}$ (valeur $b$).
\end{itemize}
\emph{Cette démonstration est exigible au BEPC.}
\end{propriete}

\begin{experience}[Démonstration guidée]
Prenons $x_1, x_2 \in \mathbb{R}$ tels que $x_1 < x_2$.
\begin{enumerate}
\item Multiplions l'inégalité par $a$ :
  \begin{itemize}
  \item Si $a > 0$, le sens est conservé : $ax_1 < ax_2$.
  \item Si $a < 0$, le sens est inversé : $ax_1 > ax_2$.
  \item Si $a = 0$, on a $ax_1 = ax_2 = 0$.
  \end{itemize}
\item Ajoutons $b$ des deux côtés (cela ne change pas l'ordre) :
  \begin{itemize}
  \item $a>0 \implies ax_1+b < ax_2+b \implies f(x_1) < f(x_2)$.
  \item $a<0 \implies ax_1+b > ax_2+b \implies f(x_1) > f(x_2)$.
  \item $a=0 \implies f(x_1)=b=f(x_2)$.
  \end{itemize}
\item Conclusion : le signe de $a$ détermine seul le sens de variation.
\end{enumerate}
\end{experience}

\courssub{Tableau de variation : construction et lecture}

Le tableau de variation synthétise l'évolution de $f$ sur $\mathbb{R}$. Pour une application affine, il est très simple car il n'y a pas de valeur remarquable intérieure (pas d'extremum).

\begin{methode}[Construire le tableau de variation de $f(x)=ax+b$]
\begin{enumerate}
\item Écrire la ligne $x$ avec les bornes $-\infty$ et $+\infty$ (et une colonne intermédiaire pour la flèche).
\item Sur la ligne $f'(x)$ (ou « signe de $a$ »), noter $+$ si $a>0$, $-$ si $a<0$, $0$ si $a=0$.
\item Sur la ligne $f$ :
  \begin{itemize}
  \item Si $a>0$ : $f(-\infty)=-\infty$, flèche $\nearrow$, $f(+\infty)=+\infty$.
  \item Si $a<0$ : $f(-\infty)=+\infty$, flèche $\searrow$, $f(+\infty)=-\infty$.
  \item Si $a=0$ : $f(x)=b$ constant, trait horizontal (ou flèche $\rightarrow$).
  \end{itemize}
\end{enumerate}
\end{methode}

\begin{exemplebox}[Exemple chiffré : $f(x) = -0,5x + 10$]
Ici $a = -0,5 < 0$, donc $f$ est strictement décroissante sur $\mathbb{R}$.
Le tableau de variation est :
\[
\begin{tabular}{|c|ccc|}\hline
$x$ & $-\infty$ & & $+\infty$ \\ \hline
$f'(x)$ & & $-$ & \\ \hline
$f$ & $+\infty$ & $\searrow$ & $-\infty$ \\ \hline
\end{tabular}
\]
\end{exemplebox}

\begin{attention}[Piège classique : ne pas confondre $a$ et $b$]
L'ordonnée à l'origine $b$ détermine \emph{où} la droite coupe l'axe des ordonnées, mais \textbf{n'influence pas le sens de variation}. Une fonction $f(x)=2x-100$ est croissante ($a=2>0$) même si elle prend des valeurs négatives pour de petits $x$. À l'inverse, $g(x)=-0,01x+1000$ est décroissante ($a=-0,01<0$) malgré son ordonnée à l'origine très grande.
\end{attention}

\begin{popfigure}
\centering
\begin{minipage}{0.45\linewidth}
\centering
\begin{tabular}{|c|ccc|}\hline
$x$ & $-\infty$ & & $+\infty$ \\ \hline
$f'(x)$ & & $+$ & \\ \hline
$f$ & $-\infty$ & $\nearrow$ & $+\infty$ \\ \hline
\end{tabular}
\par\medskip
\textbf{(a) $a>0$ : croissante}
\end{minipage}
\hfill
\begin{minipage}{0.45\linewidth}
\centering
\begin{tabular}{|c|ccc|}\hline
$x$ & $-\infty$ & & $+\infty$ \\ \hline
$f'(x)$ & & $-$ & \\ \hline
$f$ & $+\infty$ & $\searrow$ & $-\infty$ \\ \hline
\end{tabular}
\par\medskip
\textbf{(b) $a<0$ : décroissante}
\end{minipage}
\legende{Tableaux de variation types pour une application affine sur $\mathbb{R}$.}
\end{popfigure}

\courssub{Lien géométrique : coefficient directeur et pente}

Dans un repère orthonormé, l'application affine $f(x)=ax+b$ est représentée par une droite $\mathcal{D}$ d'équation $y=ax+b$. Le coefficient directeur $a$ est exactement la \textbf{pente} de cette droite.

\begin{popfigure}
\begin{tikzpicture}[scale=1.2, >=Stealth]
% Axes
\draw[->] (-0.5,0) -- (5,0) node[right] {$x$};
\draw[->] (0,-0.5) -- (0,4) node[above] {$y$};
% Droite y = 0.8x + 1
\draw[popblue, thick, domain=-0.5:3.5] plot (\x, {0.8*\x+1});
\node[popblue, right] at (3.5, 0.8*3.5+1) {$y=ax+b$};
% Points
\coordinate (A) at (1, 0.8*1+1);
\coordinate (B) at (3, 0.8*3+1);
\fill[poporange] (A) circle (2pt) node[below left] {$A(x_1, f(x_1))$};
\fill[poporange] (B) circle (2pt) node[above right] {$B(x_2, f(x_2))$};
% Triangle rectangle
\draw[popgreen, dashed] (A) -- (3, 0.8*1+1) node[midway, below] {$\Delta x = x_2-x_1$};
\draw[popgreen, dashed] (3, 0.8*1+1) -- (B) node[midway, right] {$\Delta y = f(x_2)-f(x_1)$};
\draw[popgreen, thick] (A) -- (B);
% Pente
\node[popgreen, align=center] at (2.2, 2.5) {$a = \dfrac{\Delta y}{\Delta x}$};
% Graduations
\foreach \x in {1,2,3} \draw (\x, -0.1) -- (\x, 0.1) node[below] {\x};
\foreach \y in {1,2,3} \draw (-0.1, \y) -- (0.1, \y) node[left] {\y};
\end{tikzpicture}
\legende{Visualisation de la pente $a = \frac{\Delta y}{\Delta x}$ : triangle rectangle sous la droite.}
\end{popfigure}

\emph{Intuition :} si $a>0$, quand $x$ augmente ($\Delta x>0$), $y$ augmente ($\Delta y>0$) : la droite « monte » en allant vers la droite. Si $a<0$, $\Delta y$ est de signe opposé à $\Delta x$ : la droite « descend ».

\courssub{Comparaison de deux applications affines}

Soient $f(x)=ax+b$ et $g(x)=a'x+b'$ deux applications affines. Résoudre $f(x) < g(x)$ (ou $f(x)=g(x)$) revient à comparer les deux droites.

\begin{propriete}[Intersection et ordre relatif]
L'équation $f(x)=g(x)$ s'écrit $(a-a')x = b'-b$.
\begin{itemize}
\item Si $a \neq a'$ (droites non parallèles), il existe une \textbf{unique solution} $x_0 = \dfrac{b'-b}{a-a'}$.
  \begin{itemize}
  \item Pour $x < x_0$, l'ordre entre $f(x)$ et $g(x)$ dépend du signe de $a-a'$.
  \item Pour $x > x_0$, l'ordre s'inverse.
  \end{itemize}
\item Si $a = a'$ et $b \neq b'$ (droites parallèles distinctes), pas de solution : l'une est toujours au-dessus de l'autre.
\item Si $a = a'$ et $b = b'$ (droites confondues), $f(x)=g(x)$ pour tout $x$.
\end{itemize}
\end{propriete}

\begin{exoresolu}[Comparer $f(x)=2x-3$ et $g(x)=-x+6$]
\textbf{Énoncé :} Résoudre $f(x) < g(x)$ et représenter graphiquement.
\textbf{Solution :}
\begin{align*}
f(x) < g(x) &\iff 2x-3 < -x+6 \\
&\iff 3x < 9 \\
&\iff x < 3.
\end{align*}
L'intersection a lieu pour $x_0=3$ : $f(3)=g(3)=3$.
Pour $x<3$, $f(x) < g(x)$ ; pour $x>3$, $f(x) > g(x)$.
\end{exoresolu}

\begin{popfigure}
\begin{tikzpicture}
\begin{axis}[
    width=\linewidth,
    height=7cm,
    axis lines=middle,
    xlabel=$x$, ylabel=$y$,
    xmin=-1, xmax=5,
    ymin=-2, ymax=8,
    xtick={-1,0,1,2,3,4,5},
    ytick={-2,0,2,4,6,8},
    grid=major,
    grid style={popline},
    legend style={at={(0.98,0.02)}, anchor=south east, draw=popdark, fill=popcream},
    clip=false
]
\addplot[popblue, thick, domain=-1:5, samples=100] {2*x-3};
\addlegendentry{$f(x)=2x-3$}
\addplot[poporange, thick, domain=-1:5, samples=100] {-x+6};
\addlegendentry{$g(x)=-x+6$}
% Intersection point
\addplot[only marks, mark=*, mark size=2.5, popgreen] coordinates {(3,3)};
\node[popgreen, above left] at (axis cs:3,3) {$I(3,3)$};
% Annotation for inequality region
\node[popblue, align=center, font=\small] at (axis cs:1,5) {$f(x) < g(x)$\\pour $x < 3$};
\node[poporange, align=center, font=\small] at (axis cs:4,5) {$f(x) > g(x)$\\pour $x > 3$};
\end{axis}
\end{tikzpicture}
\legende{Graphiques de $f$ et $g$ : intersection en $I(3,3)$ ; pour $x<3$, la courbe de $f$ (bleue) est sous celle de $g$ (orange).}
\end{popfigure}

\courssub{Ouverture : applications affines par intervalles (aperçu)}

Certaines situations réelles ne sont pas modélisées par une seule formule affine sur tout $\mathbb{R}$, mais par \emph{plusieurs} formules selon l'intervalle où se trouve $x$ (tarifs progressifs, tranches d'imposition, forfaits téléphoniques MTN/Orange avec seuil). On parle alors de \textbf{fonctions affines par morceaux} ou \textbf{applications affines par intervalles}. Leur étude (variation, représentation graphique, résolution d'équations) repose sur les mêmes outils, appliqués \emph{sur chaque intervalle}. Vous approfondirez ce thème l'an prochain.

\begin{savaistu}[Le saviez-vous ?]
Le concept de « pente » ($a$) est utilisé par les ingénieurs de la SONARA (raffinerie de Limbé) pour dimensionner les pipelines : la pente détermine la vitesse d'écoulement du pétrole. En topographie, le coefficient directeur d'une droite de niveau sur une carte IGN du Cameroun donne la déclivité du terrain.
\end{savaistu}

\coursec{Représentation graphique dans un repère orthonormé}

Dans la section précédente, nous avons appris à lire le \cle{sens de variation} d'une application affine $f(x)=ax+b$ directement sur son écriture algébrique : si $a>0$ elle est croissante, si $a<0$ elle est décroissante. Mais les mathématiciens aiment aussi \emph{voir} les choses. Comment représenter une application affine par un dessin ? Comment, à partir d'un dessin, retrouver des images, des antécédents, voire résoudre des inéquations sans aucun calcul ? C'est tout l'objet de cette section : le passage du \textbf{calcul} au \textbf{graphique}, et retour.

\courssub{Le graphe d'une application affine est une droite}

\begin{definition}[Courbe représentative]
Soit $f$ une application affine définie sur $\mathbb{R}$ par $f(x)=ax+b$. Dans un repère $(O\,;I,J)$ du plan, la \cle{courbe représentative} de $f$, notée $\mathcal{C}_f$, est l'ensemble des points $M(x\,;y)$ du plan tels que :
\[ y = f(x) \quad \text{c'est-à-dire} \quad y = ax+b. \]
\end{definition}

\begin{propriete}[Forme graphique]
La courbe représentative d'une application affine $f(x)=ax+b$ est une \cle{droite} $(D)$, appelée droite d'\cle{équation} $y=ax+b$.
\begin{itemize}
\item Le nombre $a$ est le \cle{coefficient directeur} de la droite.
\item Le nombre $b$ est l'\cle{ordonnée à l'origine} : la droite coupe l'axe des ordonnées au point $A(0\,;b)$.
\end{itemize}
\end{propriete}

\begin{attention}[Deux cas particuliers à reconnaître]
\begin{itemize}
\item Si $a=0$ : la droite a pour équation $y=b$. Elle est \textbf{parallèle à l'axe des abscisses} (application constante).
\item Si $b=0$ : la droite a pour équation $y=ax$. Elle \textbf{passe par l'origine} $O$ du repère : c'est le cas des applications \emph{linéaires}, qui traduisent les situations de proportionnalité.
\end{itemize}
\end{attention}

\courssub{Méthode de tracé : deux points suffisent}

Puisque le graphe est une droite, il suffit de connaître \textbf{deux points} pour la tracer entièrement à la règle. On choisit des abscisses simples pour limiter les erreurs de calcul.

\begin{methode}[Tracer la courbe $\mathcal{C}_f$ de $f(x)=ax+b$]
\begin{enumerate}
\item \textbf{Choisir le repère} : un repère orthonormé si possible (même unité sur les deux axes).
\item \textbf{Placer le point $A(0\,;b)$} : c'est l'ordonnée à l'origine, il se place directement sur l'axe des ordonnées.
\item \textbf{Placer un second point $B(x_0\,;\,ax_0+b)$} avec une abscisse $x_0$ simple : $x_0=1$ donne $B(1\,;a+b)$, mais $-1$ ou $2$ conviennent aussi.
\item \textbf{Tracer la droite $(AB)$} à la règle, en la prolongeant de part et d'autre.
\item \textbf{Légender} la droite : écrire son équation $y=ax+b$ à côté du tracé.
\end{enumerate}
\end{methode}

\begin{exemplebox}[Tracé de la droite d'équation $y=2x-1$]
Soit $f(x)=2x-1$. On a $a=2$ et $b=-1$.
\begin{itemize}
\item Premier point : $A(0\,;b)$, donc $A(0\,;-1)$.
\item Second point : pour $x_0=1$, $f(1)=2\times 1-1=1$, donc $B(1\,;1)$.
\end{itemize}
On place $A$ et $B$, puis on trace la droite $(AB)$ : c'est $\mathcal{C}_f$.
\end{exemplebox}

\begin{popfigure}
\begin{center}
\begin{tikzpicture}[scale=1.1]
% grille
\draw[popline!40, very thin] (-2.4,-2.4) grid (3.4,3.4);
% axes
\draw[->, thick, popdark] (-2.4,0) -- (3.4,0) node[below right] {\small $x$};
\draw[->, thick, popdark] (0,-2.4) -- (0,3.4) node[above left] {\small $y$};
\draw[popdark] (0,0) node[below left] {\small $O$};
% unités
\draw[thick, popdark] (1,0.06) -- (1,-0.06) node[below] {\small $1$};
\draw[thick, popdark] (0.06,1) -- (-0.06,1) node[left] {\small $1$};
% droite y=2x-1
\draw[popblue, very thick, domain=-0.6:2.1, samples=50] plot (\x,{2*\x-1}) node[right, popblue] {\small $y=2x-1$};
% points A et B
\fill[poporange] (0,-1) circle (2.2pt) node[below left, poporange] {\small $A(0\,;-1)$};
\fill[poporange] (1,1) circle (2.2pt) node[above right, poporange] {\small $B(1\,;1)$};
% Lecture image de 0.5 : f(0.5)=0
\draw[dashed, popgreen, thick] (0.5,0) -- (0.5,0) node[below, popgreen] {\small $0{,}5$}; % point sur axe x
\draw[dashed, popgreen, thick, ->] (0.5,0) -- (0,0) node[left, popgreen] {\small $0$}; % horizontal vers y
% Antécédent de 3 : f(x)=3 => x=2
\draw[dashed, poppurple, thick] (0,3) node[left, poppurple] {\small $3$} -- (2,3); % horizontal vers droite
\draw[dashed, poppurple, thick, ->] (2,3) -- (2,0) node[below, poppurple] {\small $2$}; % vertical vers x
% Points d'intersection
\fill[popgreen] (0.5,0) circle (2pt);
\fill[poppurple] (2,3) circle (2pt);
\end{tikzpicture}
\end{center}
\legende{Tracé de $y=2x-1$ par les points $A(0\,;-1)$ et $B(1\,;1)$. En vert : lecture de l'image de $0{,}5$ (qui vaut $0$). En violet : lecture de l'antécédent de $3$ (qui vaut $2$).}
\end{popfigure}

\courssub{Lecture graphique : images et antécédents}

Le grand intérêt du graphique est de lire des résultats \emph{sans calcul}. Mais il faut respecter le sens de lecture, sinon tout est inversé.

\begin{methode}[Lire une image et un antécédent sur le graphique]
\begin{itemize}
\item \textbf{Image de $x_0$} (on cherche $f(x_0)$) : on part de $x_0$ sur l'\textbf{axe des abscisses} (horizontal), on monte ou descend \textbf{verticalement} jusqu'à la droite, puis on rejoint \textbf{horizontalement} l'\textbf{axe des ordonnées}. La valeur lue est $f(x_0)$.
\item \textbf{Antécédent de $y_0$} (on cherche $x$ tel que $f(x)=y_0$) : on part de $y_0$ sur l'\textbf{axe des ordonnées} (vertical), on rejoint \textbf{horizontalement} la droite, puis on redescend \textbf{verticalement} vers l'axe des abscisses. La valeur lue est l'antécédent.
\end{itemize}
\end{methode}

\begin{attention}[Erreur fréquente : inverser les axes]
Beaucoup d'élèves lisent l'antécédent en partant de l'axe des abscisses : c'est l'inverse ! Retiens bien :
\begin{itemize}
\item \textbf{Abscisse} $\to$ axe des $x$ $\to$ axe \textbf{horizontal} ;
\item \textbf{Ordonnée} $\to$ axe des $y$ $\to$ axe \textbf{vertical}.
\end{itemize}
Astuce : dans « image », il y a le départ sur l'axe des $x$ ; dans « antécédent », on part de l'axe des $y$. Une lecture graphique donne une valeur \emph{approchée} : pour une valeur exacte, on résout l'équation $ax+b=y_0$.
\end{attention}

\begin{exoresolu}[Lectures graphiques avec $f(x)=3x-2$]
\textbf{Énoncé.} Soit $f(x)=3x-2$. (1) Tracer $\mathcal{C}_f$ dans un repère orthonormé. (2) Lire graphiquement l'image de $0{,}5$. (3) Lire graphiquement l'antécédent de $4$.
\tcblower
\textbf{Solution.}
(1) On a $a=3$ et $b=-2$. Premier point : $A(0\,;-2)$. Second point : $f(1)=3\times 1-2=1$, donc $B(1\,;1)$. On trace la droite $(AB)$.

(2) Image de $0{,}5$ : on part de $0{,}5$ sur l'axe des abscisses, on descend verticalement jusqu'à la droite, puis on rejoint l'axe des ordonnées : on lit environ $-0{,}5$. Vérification par le calcul : $f(0{,}5)=3\times 0{,}5-2=1{,}5-2=-0{,}5$. La lecture est exacte.

(3) Antécédent de $4$ : on part de $4$ sur l'axe des ordonnées, on rejoint horizontalement la droite, puis on redescend verticalement : on lit $x=2$. Vérification : $3x-2=4 \iff 3x=6 \iff x=2$. Exact.
\end{exoresolu}

\courssub{Coefficient directeur et inclinaison de la droite}

Dans un repère \textbf{orthonormé} (même unité sur les deux axes), le coefficient directeur $a$ a une interprétation géométrique précise : quand on avance de $1$ unité en abscisse, on monte (ou descend) de $a$ unités en ordonnée. C'est pourquoi on l'appelle aussi la \emph{pente} de la droite.

\begin{propriete}[Coefficient directeur et trigonométrie]
Dans un repère orthonormé, le coefficient directeur $a$ d'une droite est la tangente de l'angle que fait la droite avec l'axe des abscisses :
\[ a = \tan\widehat{(Ox\,;\,D)}. \]
Ce résultat se démontre en 3ème à l'aide de la trigonométrie dans le triangle rectangle : pour un déplacement horizontal de $1$, le déplacement vertical est $a$, donc $\tan = \frac{\text{opposé}}{\text{adjacent}} = \frac{a}{1} = a$.
Entre deux points $M(x_M\,;y_M)$ et $N(x_N\,;y_N)$ d'une droite, on a toujours, dans \emph{tout} repère (même non orthonormé) :
\[ a = \dfrac{\Delta y}{\Delta x} = \dfrac{y_N - y_M}{x_N - x_M}. \]
\end{propriete}

\begin{attention}[Repère non orthonormé : la pente visuelle trompe !]
Si les unités sur les deux axes sont différentes (par exemple $1$ cm pour $10$ unités en abscisse et $1$ cm pour $1000$ FCFA en ordonnée, comme dans un problème de facturation MTN), l'angle visuel de la droite ne correspond \textbf{plus} à $a$. Une droite peut paraître presque horizontale alors que $a=50$ ! Dans ce cas, ne jamais mesurer l'angle au rapporteur : \textbf{calculer} $a=\dfrac{\Delta y}{\Delta x}$ avec les \emph{unités réelles} des deux axes.
\end{attention}

\begin{popfigure}
\begin{center}
\begin{tikzpicture}[xscale=0.5, yscale=1.2]
% même droite y=2x-1 mais échelles différentes
\draw[->, thick, popdark] (-1,0) -- (7,0) node[below right] {\small $x$};
\draw[->, thick, popdark] (0,-2.5) -- (0,3.5) node[above left] {\small $y$};
\draw[popdark] (0,0) node[below left] {\small $O$};
\draw[thick, popdark] (1,0.08) -- (1,-0.08) node[below] {\small $1$};
\draw[thick, popdark] (0.08,1) -- (-0.08,1) node[left] {\small $1$};
% Domaine restreint pour rester dans l'axe y visible (y max 3.5 => 2x-1 <= 3.5 => x <= 2.25)
\draw[poporange, very thick, domain=-0.4:2.25, samples=50] plot (\x,{2*\x-1}) node[right, poporange] {\small $y=2x-1$};
\fill[popblue] (0,-1) circle (2.2pt) node[below left, popblue] {\small $A(0\,;-1)$};
\fill[popblue] (1,1) circle (2.2pt) node[above right, popblue] {\small $B(1\,;1)$};
\end{tikzpicture}
\end{center}
\legende{La même droite $y=2x-1$ tracée dans un repère non orthonormé (abscisses comprimées). Elle paraît beaucoup plus « raide » : l'angle visuel ne vaut plus $\arctan(2)$. Seul le calcul $a=\frac{\Delta y}{\Delta x}=\frac{1-(-1)}{1-0}=2$ donne la vraie pente.}
\end{popfigure}

\courssub{Résolution graphique d'une inéquation}

Voici une application très puissante du graphique : résoudre une inéquation \emph{sans calcul algébrique}, simplement en observant la position de la droite par rapport à l'axe des abscisses.

\begin{methode}[Résoudre graphiquement $f(x)>0$ ou $f(x)<0$]
\begin{enumerate}
\item Tracer la droite $\mathcal{C}_f$ d'équation $y=ax+b$.
\item Repérer le point où elle coupe l'axe des abscisses (c'est l'antécédent de $0$, solution de $ax+b=0$).
\item $f(x)>0$ correspond aux points de la droite situés \textbf{au-dessus} de l'axe des abscisses ; $f(x)<0$ à ceux situés \textbf{en dessous}.
\item Lire les abscisses correspondantes : ce sont les solutions.
\end{enumerate}
\end{methode}

\begin{exoresolu}[Résolution graphique de $-2x+4>0$]
\textbf{Énoncé.} Soit $g(x)=-2x+4$. Résoudre graphiquement l'inéquation $-2x+4>0$, puis vérifier par le calcul.
\tcblower
\textbf{Solution.} La droite d'équation $y=-2x+4$ passe par $A(0\,;4)$ et $B(2\,;0)$. Elle coupe l'axe des abscisses en $x=2$. Comme $a=-2<0$, la droite « descend » : elle est \textbf{au-dessus} de l'axe des abscisses pour toutes les abscisses strictement inférieures à $2$.
\[ \mathcal{S} = \left]-\infty\,;\,2\right[ \quad \text{c'est-à-dire } x<2. \]
Vérification algébrique : $-2x+4>0 \iff -2x>-4 \iff x<2$ (on divise par $-2<0$, donc on change le sens de l'inégalité). Le graphique et le calcul concordent.
\end{exoresolu}

\begin{popfigure}
\begin{center}
\begin{tikzpicture}[scale=0.95]
\draw[popline!40, very thin] (-1.4,-2.4) grid (4.4,4.4);
\draw[->, thick, popdark] (-1.4,0) -- (4.4,0) node[below right] {\small $x$};
\draw[->, thick, popdark] (0,-2.4) -- (0,4.4) node[above left] {\small $y$};
\draw[popdark] (0,0) node[below left] {\small $O$};
\draw[thick, popdark] (1,0.06) -- (1,-0.06) node[below] {\small $1$};
\draw[thick, popdark] (0.06,1) -- (-0.06,1) node[left] {\small $1$};
% partie au-dessus de l'axe : x<2 en vert épais
\draw[popgreen, line width=2.5pt, domain=-1.2:1.95, samples=50] plot (\x,{-2*\x+4});
% partie en dessous : x>2 en gris
\draw[popmuted, very thick, domain=2.05:3.1, samples=50] plot (\x,{-2*\x+4});
\fill[poporange] (2,0) circle (2.5pt) node[above right, poporange] {\small $B(2\,;0)$};
\fill[popblue] (0,4) circle (2.2pt) node[above left, popblue] {\small $A(0\,;4)$};
% zone solution sur l'axe
\draw[popgreen, line width=2.5pt] (-1.3,0) -- (2,0);
\draw[popgreen] (2,0) circle (3pt);
\fill[white] (2,0) circle (2.2pt);
\draw[popgreen] (-0.4,-0.55) node {\small solutions : $x<2$};
\draw[popgreen, ->] (-0.4,-0.4) -- (-0.4,-0.08);
\end{tikzpicture}
\end{center}
\legende{Résolution graphique de $-2x+4>0$ : la partie de la droite située au-dessus de l'axe des abscisses (en vert) correspond à $x<2$. Le point $B(2\,;0)$ est exclu (inégalité stricte).}
\end{popfigure}

\begin{aretenir}[L'essentiel de la section]
\begin{itemize}
\item Le graphe de $f(x)=ax+b$ est la \textbf{droite} d'équation $y=ax+b$ ; deux points suffisent pour la tracer, par exemple $A(0\,;b)$ et $B(1\,;a+b)$.
\item Lecture d'une \textbf{image} : axe des $x$ $\to$ verticale $\to$ droite $\to$ horizontale $\to$ axe des $y$. Lecture d'un \textbf{antécédent} : le trajet inverse.
\item Si $a=0$, la droite est parallèle à l'axe des abscisses ; si $b=0$, elle passe par l'origine.
\item En repère orthonormé, $a$ mesure l'inclinaison réelle ($a=\tan(\text{angle})$) ; en repère non orthonormé, seul le calcul $a=\frac{\Delta y}{\Delta x}$ est fiable.
\item $f(x)>0$ se résout en repérant les abscisses où la droite est \textbf{au-dessus} de l'axe des abscisses.
\end{itemize}
\end{aretenir}

\begin{savaistu}[D'où vient le mot « affine » ?]
Le terme \emph{affine} vient du latin \emph{affinis}, qui signifie « parent », « allié », « voisin ». En géométrie, une transformation affine est une transformation qui conserve le parallélisme : deux droites parallèles restent parallèles après transformation. C'est exactement ce qui se passe avec l'ombre d'un objet projetée par le soleil sur un mur ou sur le sol : les rails parallèles d'une voie ferrée gardent des ombres parallèles ! Les applications affines $x\mapsto ax+b$ sont les plus simples de ces transformations, et tu les retrouveras partout : tarifs de transport, factures d'électricité, forfaits téléphoniques...
\end{savaistu}

\begin{exercice}[À toi de jouer]
Soit $h(x)=-x+3$.
\begin{enumerate}
\item Tracer $\mathcal{C}_h$ dans un repère orthonormé en précisant les deux points utilisés.
\item Lire graphiquement l'image de $-1$ et l'antécédent de $5$, puis vérifier par le calcul.
\item Résoudre graphiquement l'inéquation $-x+3\geqslant 0$.
\end{enumerate}
\end{exercice}

\begin{corrige}[Exercice : À toi de jouer]
\begin{enumerate}
\item $a=-1$, $b=3$. Points : $A(0\,;3)$ et $B(1\,;2)$ (car $h(1)=-1+3=2$). On trace la droite $(AB)$.
\item Image de $-1$ : lecture sur le graphique, on trouve $4$. Vérification : $h(-1)=-(-1)+3=1+3=4$. Antécédent de $5$ : on lit $x=-2$. Vérification : $-x+3=5 \iff -x=2 \iff x=-2$.
\item La droite coupe l'axe des abscisses en $x=3$ et est au-dessus (ou sur) l'axe pour $x\leqslant 3$. Donc $\mathcal{S}=\left]-\infty\,;\,3\right]$. Vérification : $-x+3\geqslant 0 \iff -x\geqslant -3 \iff x\leqslant 3$.
\end{enumerate}
\end{corrige}

\coursec{Applications affines par intervalles (Fonctions par morceaux)}

Dans la section précédente, nous avons appris à représenter graphiquement une application affine $f(x)=ax+b$ : sa représentation est une \cle{droite}, tracée à partir de deux points. Mais dans la vie courante, une seule formule ne suffit pas toujours ! Quand tu paies ta facture d'électricité ENEO, le prix du kilowattheure n'est pas le même si tu consommes peu ou beaucoup : il change \emph{par tranches}. Le prix d'une course de taxi n'est pas calculé pareil le jour et la nuit. Comment les mathématiciens décrivent-ils ces situations ? Avec des \cle{applications affines par intervalles}, aussi appelées \cle{fonctions affines par morceaux}. C'est l'objet de cette section, qui te servira énormément au BEPC et dans la vie de tous les jours.

\courssub{Une situation pour comprendre : la facture d'électricité}

Imaginons une compagnie d'électricité qui facture ainsi :
\begin{itemize}
\item si tu consommes entre $0$ et $100$ kWh, tu paies $50$ F par kWh ;
\item si tu consommes entre $100$ et $500$ kWh, le tarif change ;
\item au-delà de $500$ kWh, le tarif change encore.
\end{itemize}

Il est impossible d'écrire une seule formule $f(x)=ax+b$ valable pour toutes les consommations $x$. On écrit donc \emph{plusieurs} formules, chacune valable sur son propre intervalle. C'est exactement l'idée d'une fonction définie par morceaux : \textbf{on découpe l'ensemble de définition en morceaux (intervalles), et sur chaque morceau on applique sa propre formule affine}.

\begin{definition}[Fonction affine par intervalles (ou par morceaux)]
Une fonction $f$ est dite \cle{affine par intervalles} (ou \cle{affine par morceaux}) si son ensemble de définition $D_f$ est réunion d'intervalles $I_1, I_2, \ldots, I_n$ \textbf{disjoints deux à deux} (qui ne se chevauchent pas), et si sur chaque intervalle $I_i$, $f$ coïncide avec une application affine :
\[ f(x) = a_i x + b_i \quad \text{pour } x \in I_i. \]
Les coefficients $a_i$ et $b_i$ peuvent être différents d'un intervalle à l'autre.
\end{definition}

\begin{aretenir}[Les trois conditions à retenir]
\begin{enumerate}
\item Chaque morceau est une formule \textbf{affine} : $a_i x + b_i$.
\item Les intervalles sont \textbf{disjoints} : un nombre $x$ ne peut appartenir qu'à un seul intervalle, donc une seule formule s'applique à lui.
\item La réunion des intervalles recouvre tout l'ensemble de définition : $\bigcup I_i = D_f$.
\end{enumerate}
\end{aretenir}

\begin{exemplebox}[Tarif d'électricité par tranches]
Une compagnie facture la consommation $x$ (en kWh) selon la fonction $f$ (en francs CFA) :
\[
f(x) = \begin{cases}
50x & \text{si } 0 \leq x \leq 100 \\
30x + 2000 & \text{si } 100 < x \leq 500 \\
20x + 7000 & \text{si } x > 500
\end{cases}
\]
Les trois intervalles $[0\,;\,100]$, $]100\,;\,500]$ et $]500\,;\,+\infty[$ sont disjoints et leur réunion est $[0\,;\,+\infty[$ : c'est bien une fonction affine par morceaux.
\end{exemplebox}

\begin{savaistu}[Les tranches dans la vie camerounaise]
La tarification par tranches est partout autour de toi : l'électricité (ENEO), l'eau (CAMWATER), les impôts sur les revenus (tranches marginales), les forfaits de données mobiles MTN ou Orange, le prix d'une course de taxi-moto (\og benskin\fg{}) qui double la nuit à Yaoundé ou Douala. Comprendre les fonctions par morceaux, c'est comprendre ta facture !
\end{savaistu}

\courssub{Étude algébrique : images et antécédents}

La règle d'or est simple : \textbf{avant tout calcul, demande-toi dans quel intervalle se trouve $x$}. Ensuite seulement, applique la formule correspondante.

\begin{methode}[Étudier une fonction affine par intervalles]
\begin{enumerate}
\item \textbf{Identifier l'intervalle} : repérer à quel intervalle $I_i$ appartient la valeur de $x$ (attention aux crochets : ouvert ou fermé !).
\item \textbf{Appliquer la bonne formule} : calculer $f(x) = a_i x + b_i$ avec la formule de cet intervalle, et uniquement celle-là.
\item \textbf{Pour les variations} : étudier le signe de chaque coefficient $a_i$ sur son intervalle, puis faire un tableau de variations par morceaux.
\item \textbf{Pour le graphique} : tracer chaque morceau (un segment ou une demi-droite) sur son intervalle, en gérant soigneusement les extrémités : point \textbf{plein} si la borne est incluse, point \textbf{vide} si elle est exclue.
\end{enumerate}
\end{methode}

\begin{exoresolu}[Images avec le tarif d'électricité]
On reprend la fonction $f$ de l'exemple précédent :
\[
f(x) = \begin{cases}
50x & \text{si } 0 \leq x \leq 100 \\
30x + 2000 & \text{si } 100 < x \leq 500 \\
20x + 7000 & \text{si } x > 500
\end{cases}
\]
Calculer $f(80)$, $f(300)$ et $f(600)$.
\tcblower
\textbf{Solution détaillée.}

\textbf{Image de 80.} Comme $0 \leq 80 \leq 100$, on utilise la première formule :
\[ f(80) = 50 \times 80 = 4000. \]
Une consommation de $80$ kWh coûte $4000$ F.

\textbf{Image de 300.} Comme $100 < 300 \leq 500$, on utilise la deuxième formule :
\[ f(300) = 30 \times 300 + 2000 = 9000 + 2000 = 11000. \]
Une consommation de $300$ kWh coûte $11000$ F.

\textbf{Image de 600.} Comme $600 > 500$, on utilise la troisième formule :
\[ f(600) = 20 \times 600 + 7000 = 12000 + 7000 = 19000. \]
Une consommation de $600$ kWh coûte $19000$ F.
\end{exoresolu}

\begin{exoresolu}[Recherche d'antécédent par morceaux]
Toujours avec la même fonction $f$, déterminer le (ou les) antécédent(s) de $8000$.
\tcblower
\textbf{Solution détaillée.} On résout l'équation $f(x) = 8000$ \textbf{sur chaque intervalle}, puis on vérifie que la solution trouvée appartient bien à cet intervalle.

\textbf{Sur $[0\,;\,100]$ :} $50x = 8000$ donne $x = 160$. Or $160 \notin [0\,;\,100]$ : \textbf{rejeté}.

\textbf{Sur $]100\,;\,500]$ :} $30x + 2000 = 8000$ donne $30x = 6000$, soit $x = 200$. Or $100 < 200 \leq 500$ : \textbf{accepté}.

\textbf{Sur $]500\,;\,+\infty[$ :} $20x + 7000 = 8000$ donne $20x = 1000$, soit $x = 50$. Or $50 \notin \;]500\,;\,+\infty[$ : \textbf{rejeté}.

\textbf{Conclusion :} $8000$ a un unique antécédent : $x = 200$. Une facture de $8000$ F correspond à une consommation de $200$ kWh.
\end{exoresolu}

\begin{attention}[Les deux pièges mortels]
\begin{itemize}
\item \textbf{Ne mélange jamais les formules !} Pour calculer $f(80)$, utiliser $30x+2000$ parce que \og le résultat serait plus grand \fg{} est une erreur grave. Pose-toi toujours d'abord la question : \emph{« $x$ appartient à quel intervalle ? »}.
\item \textbf{Une solution d'équation n'est valable que si elle tombe dans son intervalle.} Dans la recherche d'antécédents, résoudre $50x = 8000$ donne $x = 160$, mais cette solution est \emph{hors de l'intervalle} $[0\,;\,100]$ : il faut la rejeter. Oublier cette vérification, c'est perdre des points au BEPC.
\end{itemize}
\end{attention}

\courssub{Sens de variation morceau par morceau}

Sur chaque intervalle, la fonction est affine : son sens de variation est donc donné par le signe du coefficient $a_i$, comme on l'a vu dans la section 3 :
\begin{itemize}
\item si $a_i > 0$, $f$ est \textbf{croissante} sur $I_i$ ;
\item si $a_i < 0$, $f$ est \textbf{décroissante} sur $I_i$ ;
\item si $a_i = 0$, $f$ est \textbf{constante} sur $I_i$.
\end{itemize}

Pour notre tarif d'électricité, les trois coefficients sont $50$, $30$ et $20$, tous strictement positifs : $f$ est croissante sur chacun des trois intervalles. On résume cela dans un tableau de variations par morceaux :

\begin{center}
\begin{tabular}{|c|ccccccc|}\hline
$x$ & $0$ & & $100$ & & $500$ & & $+\infty$ \\ \hline
$a_i$ & & $50$ & & $30$ & & $20$ & \\ \hline
$f$ & $0$ & $\nearrow$ & $5000$ & $\nearrow$ & $17000$ & $\nearrow$ & $+\infty$ \\ \hline
\end{tabular}
\end{center}

Vérifions les valeurs aux bornes : $f(100) = 50 \times 100 = 5000$ ; $f(500) = 30 \times 500 + 2000 = 17000$. Remarque que la pente diminue ($50$, puis $30$, puis $20$) : plus tu consommes, plus le prix par kWh supplémentaire est faible, mais la facture totale augmente toujours.

\courssub{Représentation graphique : points pleins et points vides}

Pour tracer la représentation graphique d'une fonction affine par morceaux, on trace \textbf{chaque morceau séparément}, sur son intervalle uniquement. La question délicate est : que faire aux bornes des intervalles ?

\begin{aretenir}[Point plein, point vide]
\begin{itemize}
\item Si la borne est \textbf{incluse} dans l'intervalle (crochet fermé $[$ ou $]$), on marque l'extrémité du segment par un \cle{point plein} $\bullet$ : le point appartient à la courbe.
\item Si la borne est \textbf{exclue} (crochet ouvert $]$ ou $[$), on marque l'extrémité par un \cle{point vide} $\circ$ : le point n'appartient pas à la courbe (il est \og laissé \fg{} à l'autre morceau).
\end{itemize}
Oublier les points pleins et vides au graphique est une erreur classique qui coûte des points !
\end{aretenir}

\begin{popfigure}
\begin{center}
\begin{tikzpicture}[scale=1]
\begin{axis}[
    width=0.92\linewidth, height=7cm,
    axis lines=middle,
    xlabel={$x$ (kWh)}, ylabel={$f(x)$ (F)},
    xmin=-20, xmax=650, ymin=-1000, ymax=21000,
    xtick={0,100,200,300,400,500,600},
    ytick={5000,10000,15000,20000},
    grid=both, grid style={popline!40},
    tick label style={font=\small},
    label style={font=\small},
    scaled y ticks=false,
    yticklabel style={/pgf/number format/fixed}
]
% Morceau 1 : [0,100] -> 50x
\addplot[popcurve, domain=0:100, samples=2] {50*x};
% Morceau 2 : ]100,500] -> 30x+2000
\addplot[popcurve, poporange, domain=100:500, samples=2] {30*x+2000};
% Morceau 3 : ]500, +inf[ -> 20x+7000
\addplot[popcurve, popgreen, domain=500:640, samples=2] {20*x+7000};
% Points pleins et vides
\addplot[only marks, mark=*, mark size=2pt, popblue] coordinates {(0,0) (100,5000) (500,17000)};
\addplot[only marks, mark=o, mark size=2.2pt, thick, poporange] coordinates {(100,5000)};
\addplot[only marks, mark=o, mark size=2.2pt, thick, popgreen] coordinates {(500,17000)};
\end{axis}
\end{tikzpicture}
\end{center}
\legende{Représentation graphique du tarif d'électricité : trois morceaux affines. Les points pleins indiquent les bornes incluses, les points vides les bornes exclues. Ici, les morceaux se raccordent : la fonction est continue.}
\end{popfigure}

Dans cet exemple, observe bien : en $x = 100$, la première formule donne $f(100) = 5000$ (point plein bleu) et la deuxième formule donnerait $30 \times 100 + 2000 = 5000$ (point vide orange) : les deux morceaux \textbf{se touchent}. De même en $x = 500$ : $f(500) = 17000$ et $20 \times 500 + 7000 = 17000$. La courbe ne présente aucun \og trou \fg{} ni \og saut \fg{} : on dit que la fonction est \cle{continue}.

\courssub{Continuité et discontinuité aux bornes}

\begin{definition}[Fonction continue, fonction discontinue]
Une fonction affine par morceaux est dite \cle{continue} en une borne commune $c$ de deux intervalles si les deux morceaux \textbf{se raccordent} en $c$, c'est-à-dire si les deux formules donnent la même valeur en $c$. Sinon, la courbe présente un \cle{saut} (ou \cle{rupture}) en $c$ : on dit que la fonction est \cle{discontinue} en $c$.
\end{definition}

\begin{propriete}[Condition de raccordement (admise)]
Une fonction affine par morceaux est continue sur $\mathbb{R}$ si et seulement si, à chaque borne commune $c$ de deux intervalles consécutifs, les deux formules coïncident :
\[ a_1 c + b_1 = a_2 c + b_2. \]
(Ce résultat est admis en classe de 3ème ; il sera très utile au lycée.)
\end{propriete}

\begin{exemplebox}[Le taxi de jour et de nuit : une fonction discontinue]
Un taxi de Douala facture la course de $x$ kilomètres ainsi :
\[
f(x) = \begin{cases}
2x + 500 & \text{si } 0 \leq x \leq 10 \quad \text{(tarif de jour)} \\
3x + 500 & \text{si } x > 10 \quad \text{(tarif de nuit)}
\end{cases}
\]
\textbf{Image de 5 :} $0 \leq 5 \leq 10$, donc $f(5) = 2 \times 5 + 500 = 510$ F.

\textbf{Image de 15 :} $15 > 10$, donc $f(15) = 3 \times 15 + 500 = 545$ F... attendez, vérifions la continuité en $10$ !

\textbf{Étude en $x = 10$ :} avec le tarif de jour, $f(10) = 2 \times 10 + 500 = 520$ F (point plein). Mais dès que $x$ dépasse $10$, la formule de nuit s'applique et donnerait $3 \times 10 + 500 = 800$ F (point vide). Comme $520 \neq 800$, les deux morceaux \textbf{ne se raccordent pas} : il y a un \textbf{saut} de $800 - 520 = 280$ F en $x = 10$. La fonction est discontinue en $10$.
\end{exemplebox}

\begin{popfigure}
\begin{center}
\begin{tikzpicture}[scale=1]
\begin{axis}[
    width=0.9\linewidth, height=6.5cm,
    axis lines=middle,
    xlabel={$x$ (km)}, ylabel={$f(x)$ (F)},
    xmin=-1, xmax=17, ymin=400, ymax=1100,
    xtick={0,5,10,15},
    ytick={500,520,800,950},
    grid=both, grid style={popline!40},
    tick label style={font=\small},
    label style={font=\small}
]
% Jour : [0,10] -> 2x+500
\addplot[popcurve, domain=0:10, samples=2] {2*x+500};
% Nuit : ]10,16] -> 3x+500
\addplot[popcurve, poppurple, domain=10:16, samples=2] {3*x+500};
% Points
\addplot[only marks, mark=*, mark size=2pt, popblue] coordinates {(0,500) (10,520)};
\addplot[only marks, mark=o, mark size=2.2pt, thick, poppurple] coordinates {(10,800)};
% Annotation du saut
\draw[<->, thick, poppink] (axis cs:10.6,520) -- (axis cs:10.6,800) node[midway, right, font=\small]{saut};
\end{axis}
\end{tikzpicture}
\end{center}
\legende{Tarif du taxi : en $x = 10$, le point plein (jour, $520$ F) et le point vide (nuit, $800$ F) ne coïncident pas : la courbe présente un saut, la fonction est discontinue en $10$.}
\end{popfigure}

\begin{attention}[Rédaction au BEPC]
Quand on te demande si une fonction par morceaux est continue en une borne $c$, il faut \textbf{calculer les deux valeurs} (celle de la formule de gauche et celle de la formule de droite) et \textbf{comparer} : si elles sont égales, conclure \og continue \fg{} ; si elles diffèrent, conclure \og discontinue : il y a un saut de ... \fg{}. Une simple réponse \og oui \fg{} ou \og non \fg{} sans calcul ne rapporte aucun point.
\end{attention}

\begin{exoresolu}[Synthèse : un forfait internet]
Un cybercafé de Bafoussam facture la connexion selon la durée $x$ (en heures) :
\[
g(x) = \begin{cases}
200x & \text{si } 0 \leq x \leq 2 \\
100x + 200 & \text{si } x > 2
\end{cases}
\]
\begin{enumerate}
\item Calculer $g(1)$, $g(2)$ et $g(5)$.
\item La fonction $g$ est-elle continue en $2$ ?
\item Combien d'heures de connexion peut-on avoir avec $700$ F ?
\end{enumerate}
\tcblower
\textbf{Solution détaillée.}

\textbf{1. Images.} $1 \in [0\,;\,2]$ : $g(1) = 200 \times 1 = 200$ F. $2 \in [0\,;\,2]$ : $g(2) = 200 \times 2 = 400$ F. $5 > 2$ : $g(5) = 100 \times 5 + 200 = 700$ F.

\textbf{2. Continuité en 2.} Formule de gauche : $200 \times 2 = 400$. Formule de droite : $100 \times 2 + 200 = 400$. Les deux valeurs sont égales : les morceaux se raccordent, $g$ est \textbf{continue} en $2$.

\textbf{3. Antécédent de 700.} Sur $[0\,;\,2]$ : $200x = 700$ donne $x = 3{,}5$, rejeté car $3{,}5 \notin [0\,;\,2]$. Sur $]2\,;\,+\infty[$ : $100x + 200 = 700$ donne $100x = 500$, soit $x = 5$, accepté car $5 > 2$. Avec $700$ F, on peut se connecter \textbf{5 heures}.
\end{exoresolu}

\begin{exercice}[À toi de jouer]
Une société de transport facture le transport de marchandises entre Yaoundé et Bafoussam selon la masse $x$ (en tonnes) :
\[
h(x) = \begin{cases}
10000x & \text{si } 0 < x \leq 5 \\
8000x + 10000 & \text{si } x > 5
\end{cases}
\]
\begin{enumerate}
\item Calculer $h(3)$ et $h(8)$.
\item La fonction $h$ est-elle continue en $5$ ? Justifier.
\item Un commerçant paie $74000$ F. Quelle masse a-t-il transportée ?
\item Représenter graphiquement $h$ (points pleins et vides exigés).
\end{enumerate}
\end{exercice}

\begin{corrige}[Exercice : transport de marchandises]
\textbf{1.} $3 \in ]0\,;\,5]$ : $h(3) = 10000 \times 3 = 30000$ F. $8 > 5$ : $h(8) = 8000 \times 8 + 10000 = 74000$ F.

\textbf{2.} Formule de gauche en $5$ : $10000 \times 5 = 50000$. Formule de droite en $5$ : $8000 \times 5 + 10000 = 50000$. Les deux valeurs sont égales : $h$ est continue en $5$.

\textbf{3.} Sur $]0\,;\,5]$ : $10000x = 74000$ donne $x = 7{,}4$, rejeté ($7{,}4 > 5$). Sur $]5\,;\,+\infty[$ : $8000x + 10000 = 74000$ donne $8000x = 64000$, soit $x = 8$, accepté. Le commerçant a transporté $8$ tonnes.

\textbf{4.} Premier morceau : segment de $(0\,;\,0)$ [point vide] à $(5\,;\,50000)$ [point plein]. Deuxième morceau : demi-droite partant de $(5\,;\,50000)$ (les deux morceaux se raccordent, le point est plein) et passant par $(8\,;\,74000)$.
\end{corrige}

\begin{aretenir}[L'essentiel de la section]
\begin{itemize}
\item Une fonction \cle{affine par intervalles} applique une formule $a_i x + b_i$ différente sur chaque intervalle d'une partition de $D_f$.
\item Avant tout calcul : \textbf{identifier l'intervalle de $x$}, puis appliquer la bonne formule.
\item Pour un antécédent : résoudre sur chaque intervalle et \textbf{rejeter} les solutions hors intervalle.
\item Au graphique : point \textbf{plein} = borne incluse, point \textbf{vide} = borne exclue.
\item \cle{Continuité} en une borne $c$ : les deux formules donnent la même valeur en $c$ ; sinon il y a un \cle{saut} (discontinuité).
\end{itemize}
\end{aretenir}

Tu sais désormais manipuler ces fonctions \og à plusieurs visages \fg{} qui modélisent tant de situations concrètes. Dans la prochaine et dernière section, nous rassemblerons tout ce que tu as appris dans des problèmes de synthèse et des activités d'intégration, en vue du BEPC.

\coursec{Modélisation, problèmes de synthèse et ouverture}

Dans la section précédente, nous avons découvert les applications affines par intervalles : ces fonctions qui changent de formule selon la valeur de $x$, comme les tarifs de taxi ou les tranches d'imposition. Nous voici maintenant à l'aboutissement de toute la leçon : à quoi servent vraiment les applications linéaires et affines ? Réponse : à \cle{modéliser} des situations de la vie courante (tarifs, coûts, abonnements), à les comparer et à prendre des décisions économiques. C'est exactement l'esprit des problèmes du BEPC. Nous terminerons par une petite ouverture vers les classes supérieures, pour montrer que tout ce que tu viens d'apprendre est le fondement de notions que tu rencontreras plus tard.

\courssub{La démarche de modélisation : de la réalité aux mathématiques}

\textbf{Intuition.} Un opérateur téléphonique propose un forfait : 5\,000 FCFA fixes par mois, plus 50 FCFA par minute d'appel. Combien paieras-tu si tu appelles 120 minutes ? Et si tu as un budget de 9\,000 FCFA, combien de minutes peux-tu te permettre ? Pour répondre, on \emph{traduit} la situation en langage mathématique : c'est la modélisation. Elle suit toujours les mêmes étapes.

\begin{methode}[Les cinq étapes de la modélisation]
\begin{enumerate}
\item \textbf{Choisir les variables.} Identifier la grandeur que l'on fait varier (variable \cle{indépendante} $x$ : nombre de minutes, de kilomètres, de sacs de ciment...) et la grandeur qui en dépend (variable \cle{dépendante} $y$ : le prix, le coût total...).
\item \textbf{Poser les hypothèses.} La situation est-elle \emph{linéaire} (proportionnalité : $y = ax$), \emph{affine} (part fixe + part proportionnelle : $y = ax + b$), ou \emph{affine par morceaux} (la formule change selon les intervalles) ?
\item \textbf{Écrire le modèle mathématique.} Traduire en écriture $f(x) = \ldots$, en précisant le domaine (souvent $x \geq 0$).
\item \textbf{Résoudre.} Calculer une image, un antécédent, ou résoudre une équation du type $f_A(x) = f_B(x)$.
\item \textbf{Interpréter et valider.} Traduire le résultat en français avec l'\cle{unité}, et vérifier la cohérence : l'ordre de grandeur est-il raisonnable ? Le résultat appartient-il au domaine ?
\end{enumerate}
\end{methode}

\begin{popfigure}
\begin{tikzpicture}[
  node distance=0.6cm,
  box/.style={draw=popblue, thick, rounded corners, fill=popblueL, text width=4.2cm, align=center, minimum height=1cm, font=\small},
  fleche/.style={-{Stealth[length=3mm]}, thick, popdark}]
\node[box, align=center] (sit){\textbf{Situation concrète}\\ (tarifs, coûts, consommation)};
\node[box, below=of sit, align=center] (choix){\textbf{Choix du modèle}\\ linéaire ? affine ? par morceaux ?};
\node[box, below=of choix, align=center] (ecrit){\textbf{Écriture mathématique}\\ $f(x)=ax+b$, domaine};
\node[box, below=of ecrit, align=center] (reso){\textbf{Résolution}\\ équation, image, antécédent};
\node[box, below=of reso, fill=popgreenL, draw=popgreen, align=center] (concl){\textbf{Conclusion interprétée}\\ avec unité + validation};
\draw[fleche] (sit) -- (choix);
\draw[fleche] (choix) -- (ecrit);
\draw[fleche] (ecrit) -- (reso);
\draw[fleche] (reso) -- (concl);
\draw[fleche, dashed, poporange] (concl.west) to[bend right=40] node[left, font=\small\itshape, text=poporange]{si incohérent : on revient au modèle} (choix.west);
\end{tikzpicture}
\legende{L'arbre de décision de la modélisation : si la conclusion est incohérente (prix négatif, ordre de grandeur absurde), on remet en cause le modèle choisi.}
\end{popfigure}

\begin{exemplebox}[Forfait téléphonique]
Un forfait coûte 5\,000 FCFA fixes plus 50 FCFA la minute. \textbf{Étape 1 :} $x$ = nombre de minutes (indépendante), $y$ = coût total en FCFA (dépendante). \textbf{Étape 2 :} part fixe + part proportionnelle $\Rightarrow$ modèle affine. \textbf{Étape 3 :} $f(x) = 50x + 5000$, pour $x \geq 0$. \textbf{Étape 4 :} pour 120 minutes : $f(120) = 50 \times 120 + 5000 = 11\,000$. Avec un budget de 9\,000 FCFA : $50x + 5000 = 9000 \Leftrightarrow 50x = 4000 \Leftrightarrow x = 80$. \textbf{Étape 5 :} 120 minutes coûtent 11\,000 FCFA ; avec 9\,000 FCFA on peut appeler 80 minutes. Cohérent : 80 minutes à 50 FCFA font bien 4\,000 FCFA, plus la part fixe.
\end{exemplebox}

\courssub{Problèmes types : comparaison de tarifs, point d'équilibre, seuil de rentabilité}

\textbf{Intuition.} Deux agences de location de voitures à Douala proposent des tarifs différents : l'une a une part fixe faible mais un prix au kilomètre élevé, l'autre l'inverse. Laquelle choisir ? Tout dépend de la distance parcourue ! Il existe une distance « charnière » où les deux tarifs sont égaux : c'est le \cle{point d'équilibre}.

\begin{definition}[Point d'équilibre de deux tarifs]
Soient $f_A$ et $f_B$ deux applications affines modélisant deux tarifs. Le \cle{point d'équilibre} (ou seuil d'indifférence) est la valeur $x_0$ solution de l'équation $f_A(x) = f_B(x)$. Pour cette quantité, les deux tarifs sont identiques.
\end{definition}

\begin{methode}[Résoudre un problème de comparaison de tarifs]
\begin{enumerate}
\item \textbf{Définir $x$} : la quantité qui varie (km, minutes, mois...).
\item \textbf{Écrire} les deux applications affines $f_A(x)$ et $f_B(x)$.
\item \textbf{Résoudre} l'équation $f_A(x) = f_B(x)$ : on obtient le seuil $x_0$.
\item \textbf{Étudier le signe} de la différence $f_A(x) - f_B(x)$ (ou comparer les coefficients directeurs) pour savoir qui est moins cher de chaque côté du seuil.
\item \textbf{Conclure par une phrase} complète : « Pour $x < x_0$, le tarif A est moins cher ; pour $x > x_0$, c'est le tarif B. »
\end{enumerate}
\end{methode}

\begin{exoresolu}[Location de voiture : lequel choisir ?]
\textbf{Énoncé.} Agence A : 5\,000 FCFA par jour plus 200 FCFA par km. Agence B : 8\,000 FCFA par jour plus 100 FCFA par km. Pour une journée, déterminer pour quelles distances l'agence A est la moins chère.
\tcblower
\textbf{Solution.} \textbf{Étape 1.} Soit $x$ la distance parcourue en km ($x \geq 0$).
\textbf{Étape 2.} $f_A(x) = 200x + 5000$ et $f_B(x) = 100x + 8000$.
\textbf{Étape 3.} Point d'équilibre :
\begin{align*}
f_A(x) &= f_B(x)\\
200x + 5000 &= 100x + 8000\\
200x - 100x &= 8000 - 5000\\
100x &= 3000\\
x &= 30
\end{align*}
\textbf{Étape 4.} Différence : $f_A(x) - f_B(x) = 100x - 3000$. C'est une application affine de coefficient directeur $100 > 0$, donc croissante, qui s'annule en $x = 30$. Ainsi $f_A(x) - f_B(x) < 0$ pour $x < 30$ et $f_A(x) - f_B(x) > 0$ pour $x > 30$.
\textbf{Étape 5.} Conclusion : pour un trajet de \textbf{moins de 30 km}, l'agence A est moins chère ; pour \textbf{plus de 30 km}, l'agence B est moins chère ; pour exactement 30 km, les deux coûtent 11\,000 FCFA. Vérification : $f_A(30) = 200 \times 30 + 5000 = 11\,000$ et $f_B(30) = 100 \times 30 + 8000 = 11\,000$. \checkmark
\end{exoresolu}

\begin{popfigure}
\begin{tikzpicture}
\begin{axis}[
  width=\linewidth, height=7.5cm,
  axis lines=middle,
  xlabel={Distance $x$ (km)}, ylabel={Coût (FCFA)},
  xmin=0, xmax=60, ymin=0, ymax=18000,
  xtick={0,10,20,30,40,50,60},
  ytick={0,5000,8000,11000,15000},
  legend style={at={(0.02,0.98)}, anchor=north west, font=\small},
  clip=false]
\addplot[popcurve, popblue, very thick, domain=0:60, samples=2]{200*x+5000};
\addplot[popcurve, poporange, very thick, domain=0:60, samples=2]{100*x+8000};
\legend{$C_A(x)=200x+5000$, $C_B(x)=100x+8000$}
\addplot[mark=*, mark size=2.5pt, popdark] coordinates {(30,11000)};
\draw[dashed, popdark] (axis cs:30,0) -- (axis cs:30,11000) -- (axis cs:0,11000);
\node[font=\small, popdark, anchor=west] at (axis cs:31,11500) {Point d'équilibre $(30\,;\,11\,000)$};
\node[font=\small, popblue, align=center] at (axis cs:12,15500) {\textbf{Avantage A}\\ ($x<30$ km)};
\node[font=\small, poporange, align=center] at (axis cs:50,6000) {\textbf{Avantage B}\\ ($x>30$ km)};
\end{axis}
\end{tikzpicture}
\legende{Graphiquement, le point d'équilibre est l'intersection des deux droites de coûts. En dessous de 30 km, la droite bleue (A) est sous la droite orange (B) : A est moins chère. Au-delà, c'est l'inverse.}
\end{popfigure}

\textbf{Intersection de deux droites : le lien algèbre-géométrie.} Ce que nous venons de faire algébriquement a un sens géométrique précis : résoudre $f_A(x) = f_B(x)$, c'est chercher l'abscisse du point d'intersection des deux droites représentatives. Résoudre le système
\[
\begin{cases} y = 200x + 5000 \\ y = 100x + 8000 \end{cases}
\]
c'est trouver les coordonnées $(30\,;\,11\,000)$ du point commun aux deux droites. La \cle{résolution algébrique} donne la valeur exacte ; la \cle{résolution graphique} permet de visualiser et de vérifier, mais ne donne souvent qu'une valeur approchée. Au BEPC, on utilise les deux : le graphique pour conjecturer, le calcul pour prouver.

\begin{attention}[Les deux pièges classiques de la conclusion]
\begin{itemize}
\item \textbf{L'unité oubliée.} Une réponse « $x = 30$ » seule ne vaut rien : il faut écrire « 30 km ». Le seuil est une \emph{distance}, le coût est en \emph{FCFA}, la durée en \emph{mois}... Précise toujours l'unité dans la phrase de conclusion.
\item \textbf{Le seuil hors domaine.} Si tu trouves $x_0 = -5$ alors que $x$ désigne un nombre de kilomètres ($x \geq 0$), le point d'équilibre n'a aucun sens concret : l'un des tarifs est alors \emph{toujours} le moins cher sur le domaine. Vérifie systématiquement que le seuil appartient au domaine de définition avant de conclure.
\end{itemize}
\end{attention}

\begin{exemplebox}[Seuil de rentabilité d'un petit commerce]
Mama Ngo fabrique des beignets à Yaoundé. Ses coûts : 10\,000 FCFA de frais fixes (charbon, location de l'emplacement) plus 30 FCFA par beignet (farine, huile, sucre). Elle vend chaque beignet 50 FCFA. Le \cle{seuil de rentabilité} est le nombre de beignets à partir duquel elle commence à gagner de l'argent. Coût : $C(x) = 30x + 10\,000$ ; recette : $R(x) = 50x$. Équilibre : $50x = 30x + 10\,000 \Leftrightarrow 20x = 10\,000 \Leftrightarrow x = 500$. Elle doit vendre \textbf{au moins 500 beignets} pour couvrir ses frais ; à partir du 501\textsuperscript{e}, elle réalise un bénéfice.
\end{exemplebox}

\courssub{Ouverture vers la classe de Première : le nombre dérivé}

\textbf{Intuition.} Reprenons $f(x) = 50x + 5000$. Si tu appelles une minute de plus, ton coût augmente de 50 FCFA, \emph{quelle que soit} ta consommation actuelle. Cette « vitesse de variation » constante est une idée immense des mathématiques : elle deviendra, en classe de Première, la notion de \cle{nombre dérivé}.

\begin{propriete}[Taux d'accroissement d'une fonction affine — ouverture Première]
Soit $f$ une application affine définie par $f(x) = ax + b$. Pour tout réel $x$ et tout réel $h \neq 0$, le \cle{taux d'accroissement} de $f$ entre $x$ et $x + h$ est constant :
\[
\frac{f(x+h) - f(x)}{h} = \frac{a(x+h) + b - (ax + b)}{h} = \frac{ah}{h} = a.
\]
En classe de Première, on appellera ce nombre le \cle{nombre dérivé} de $f$ en $x$, noté $f'(x)$. Ainsi, pour une fonction affine : $f'(x) = a$ pour tout $x$. Le coefficient directeur, le taux d'accroissement et la dérivée sont \emph{trois visages du même nombre}.
\end{propriete}

Ce résultat confirme ce que tu sais déjà : le coefficient directeur $a$ mesure la variation de $f$ quand $x$ augmente de 1. Dire « $f'(x) = a$ » revient à dire « la fonction affine varie partout à la même vitesse ».

\courssub{Ouverture : l'approximation affine, ou pourquoi les tangentes sont des droites}

\textbf{Intuition.} Zoome sur une courbe, n'importe laquelle, en un point : plus tu zoomes, plus la courbe ressemble à une droite. Cette droite « limite » est la \cle{tangente} à la courbe en ce point, et son équation est... affine ! C'est le principe de l'\cle{approximation affine} : près d'un point, on peut remplacer une fonction compliquée par une fonction affine qui lui ressemble.

\begin{popfigure}
\begin{tikzpicture}
\begin{axis}[
  width=\linewidth, height=7.5cm,
  axis lines=middle,
  xlabel={$x$}, ylabel={$y$},
  xmin=-0.5, xmax=2.5, ymin=-1.5, ymax=5,
  xtick={1}, ytick={1},
  legend style={at={(0.02,0.98)}, anchor=north west, font=\small}]
\addplot[popcurve, popblue, very thick, domain=-0.4:2.3, samples=100]{x^2};
\addplot[popcurve, poporange, very thick, domain=-0.2:2.4, samples=2]{2*x-1};
\legend{$y=x^2$ (courbe), $y=2x-1$ (tangente en $x=1$)}
\addplot[mark=*, mark size=2.5pt, popdark] coordinates {(1,1)};
\node[font=\small, popdark, anchor=south west] at (axis cs:1,1) {$(1\,;\,1)$};
\draw[popgreen, thick, dashed] (axis cs:0.6,0.36) circle (0.28);
\node[font=\small\itshape, popgreen, align=center] at (axis cs:1.75,4.4) {Près du point, courbe\\ et tangente se confondent};
\end{axis}
\end{tikzpicture}
\legende{La parabole $y = x^2$ et sa tangente au point $(1\,;\,1)$, d'équation $y = 2x - 1$. Au voisinage de ce point, la courbe « ressemble » à sa tangente : c'est l'approximation affine, que tu étudieras en Terminale.}
\end{popfigure}

\begin{exemplebox}[L'approximation affine en action]
Pour la fonction $f(x) = x^2$, au voisinage de $x = 1$, on utilise la tangente $y = 2x - 1$. Calculons $f(1,01) = 1,01^2 = 1,0201$. L'approximation affine donne $2 \times 1,01 - 1 = 1,02$. L'erreur est seulement de $0,0001$ ! Voilà pourquoi les ingénieurs adorent les fonctions affines : elles sont simples, et elles approchent \emph{toutes} les fonctions « lisses » localement. Tout ce que tu as appris dans cette leçon (coefficient directeur, équation de droite, sens de variation) est la brique de base de cette idée.
\end{exemplebox}

\begin{savaistu}[Ta fiche de paie est une fonction affine par morceaux]
L'impôt sur le revenu au Cameroun (l'IRPP, Impôt sur le Revenu des Personnes Physiques) est calculé par \emph{tranches} : une première partie du revenu est taxée à un certain taux, la suivante à un taux supérieur, et ainsi de suite. Mathématiquement, l'impôt à payer est une \textbf{application affine par intervalles} — et même \emph{continue} : à la frontière entre deux tranches, les deux formules donnent le même montant, ce qui garantit qu'on ne « saute » pas brutalement d'un impôt à un autre. Les droites se raccordent parfaitement, comme dans la section 5. Comprendre cette leçon, c'est donc déjà comprendre le mécanisme de ta future fiche de paie, les tarifs d'électricité d'ENEO ou les grilles de prix de transport !
\end{savaistu}

\courssub{Synthèse de la leçon : linéaire, affine, affine par morceaux}

\begin{center}
\begin{tabularx}{\linewidth}{|l|X|X|X|}
\hline
\rowcolor{popblueL} & \textbf{Linéaire} & \textbf{Affine} & \textbf{Affine par morceaux} \\
\hline
\textbf{Définition} & $f(x) = ax$ & $f(x) = ax + b$ & une formule affine différente sur chaque intervalle \\
\hline
\textbf{Graphe} & droite passant par l'origine $O$ & droite coupant l'axe des ordonnées en $(0\,;\,b)$ & segments de droites raccordés (parfois avec des sauts) \\
\hline
\textbf{Variation} & croissante si $a>0$, décroissante si $a<0$ & croissante si $a>0$, décroissante si $a<0$ & variation étudiée \emph{intervalle par intervalle} \\
\hline
\textbf{Modélisation} & proportionnalité pure (prix au kg, $d = vt$) & part fixe + part variable (forfait, location, taxi) & tarifs par tranches (impôts, électricité, transport) \\
\hline
\end{tabularx}
\end{center}

\begin{aretenir}[L'essentiel de la section]
\begin{itemize}
\item Modéliser, c'est suivre 5 étapes : variables $\rightarrow$ hypothèses $\rightarrow$ écriture $f(x)$ $\rightarrow$ résolution $\rightarrow$ interprétation avec unité et validation.
\item Le \cle{point d'équilibre} de deux tarifs est la solution de $f_A(x) = f_B(x)$ ; graphiquement, c'est l'intersection des deux droites. On conclut toujours par une phrase avec les unités, après avoir vérifié que le seuil est dans le domaine.
\item Ouverture : pour une fonction affine, le taux d'accroissement est constant et égal à $a$ ; en Première, ce nombre s'appellera la dérivée $f'(x) = a$. Près d'un point, toute courbe « lisse » ressemble à sa tangente, qui est une fonction affine : c'est l'approximation affine.
\end{itemize}
\end{aretenir}

\begin{exercice}[Abonnements internet]
Deux fournisseurs proposent : Fournisseur 1 : 10\,000 FCFA d'installation puis 15\,000 FCFA par mois. Fournisseur 2 : 25\,000 FCFA d'installation puis 12\,000 FCFA par mois. À partir de combien de mois le fournisseur 2 devient-il plus avantageux ?
\end{exercice}

\begin{corrige}[Exercice 1]
Soit $x$ le nombre de mois ($x \geq 0$). $f_1(x) = 15\,000x + 10\,000$ et $f_2(x) = 12\,000x + 25\,000$. Point d'équilibre : $15\,000x + 10\,000 = 12\,000x + 25\,000 \Leftrightarrow 3\,000x = 15\,000 \Leftrightarrow x = 5$. Différence : $f_1(x) - f_2(x) = 3\,000x - 15\,000$, négative pour $x < 5$, positive pour $x > 5$. Conclusion : avant 5 mois, le fournisseur 1 est moins cher ; à 5 mois, les deux coûtent 85\,000 FCFA ; \textbf{à partir du 6\textsuperscript{e} mois, le fournisseur 2 est plus avantageux}. Le seuil $x = 5$ est bien dans le domaine $x \geq 0$. \checkmark
\end{corrige}

\newpage

\coursec{Activité d'intégration}

\coursec{Activité d'intégration et de synthèse : Choix d'un transporteur pour une sortie pédagogique}

\courssub{Objectifs de l'activité}
Cette activité d'intégration vise à mobiliser l'ensemble des savoirs et savoir-faire travaillés dans la leçon sur les applications affines. Elle permet à l'élève de :
\begin{itemize}
    \item Modéliser une situation concrète par des fonctions affines, linéaires ou constantes.
    \item Calculer des images et des antécédents, résoudre des équations du premier degré pour trouver des seuils de décision.
    \item Étudier le sens de variation de chaque fonction et l'interpréter dans le contexte.
    \item Représenter graphiquement les trois fonctions dans un même repère orthonormé avec une échelle adaptée.
    \item Lire graphiquement les intersections et vérifier les résultats algébriques.
    \item Rédiger une argumentation structurée (note de décision) en s'appuyant sur le calcul et le graphique.
    \item Construire la fonction « coût minimal » par intervalles (application affine par morceaux).
\end{itemize}
L'activité se déroule en groupes de trois élèves (45 minutes de travail + 15 minutes de mise en commun) et s'inscrit dans la compétence « Résoudre des problèmes » du référentiel MINESEC.

\courssub{Situation}
L'association des élèves de la classe de 3ème du Lycée Bilingue de Nkolbisson organise une sortie pédagogique au Parc National de la Bénoué (région du Nord). Pour le transport, l'association doit louer un bus auprès d'un transporteur. Après consultation, trois transporteurs basés à Yaoundé proposent les offres suivantes :

\begin{center}
\begin{tabularx}{\linewidth}{|l|X|X|}\hline
\textbf{Transporteur} & \textbf{Description de l'offre} & \textbf{Observations} \\ \hline
\textbf{A -- Voyages \emph{Express}} & 50\,000 FCFA de frais fixes (mise à disposition du bus et chauffeur) + 200 FCFA par kilomètre parcouru. & Offre dégressive au km pour les longs trajets. \\ \hline
\textbf{B -- \emph{Transports Rapides}} & 30\,000 FCFA de frais fixes + 300 FCFA par kilomètre parcouru. & Moins cher au départ, mais coût au km plus élevé. \\ \hline
\textbf{C -- \emph{Confort Plus}} & Forfait unique de 100\,000 FCFA (kilométrage illimité, assurance et péages inclus). & Coût constant, quel que soit le nombre de km. \\ \hline
\end{tabularx}
\end{center}

Le trajet aller-retour Yaoundé -- Parc de la Bénoué -- Yaoundé fait environ 1\,200 km. Cependant, l'association hésite encore sur le programme exact : visite partielle (environ 400 km au total), visite complète (environ 800 km) ou extension vers le Lac Tchad (environ 1\,500 km). Le président de l'association demande à la commission « Logistique » (votre groupe) d'analyser les trois offres et de rédiger une note de décision claire pour le choix final.

\courssub{Consignes}
Travaillez en groupe et répondez aux questions dans l'ordre. Chaque étape doit être justifiée.

\begin{enumerate}
    \item \textbf{Modélisation.} Soit $x$ la distance parcourue en kilomètres ($x \ge 0$) et $y$ le coût total en FCFA.
    \begin{enumerate}
        \item Écrire les expressions littérales des fonctions $f_A$, $f_B$, $f_C$ qui associent à $x$ le coût total pour chaque transporteur.
        \item Préciser la nature de chaque fonction (linéaire, affine, constante) et donner son coefficient directeur et son ordonnée à l'origine.
    \end{enumerate}

    \item \textbf{Calculs numériques et tableau de valeurs.}
    Calculer le coût total pour chaque transporteur aux distances $x = 100$ km, $x = 200$ km, $x = 300$ km, $x = 500$ km et $x = 1\,000$ km. Présenter les résultats dans un tableau comparatif.

    \item \textbf{Recherche algébrique des seuils (intersections).}
    Déterminer exactement (par le calcul) les distances pour lesquelles deux offres coûtent le même prix :
    \begin{enumerate}
        \item Résoudre $f_A(x) = f_B(x)$.
        \item Résoudre $f_A(x) = f_C(x)$.
        \item Résoudre $f_B(x) = f_C(x)$.
    \end{enumerate}
    Interpréter chaque résultat : à partir de quel kilométrage une offre devient-elle plus avantageuse qu'une autre ?

    \item \textbf{Représentation graphique.}
    Dans un repère orthonormé $(O; \vec{\imath}, \vec{\jmath})$, tracer les droites représentants $f_A$, $f_B$ et $f_C$.
    \begin{itemize}
        \item Choisir une échelle adaptée sur l'axe des abscisses (distance en km) et sur l'axe des ordonnées (coût en FCFA). Indiquer les échelles choisies.
        \item Tracer les trois droites sur le même graphique en utilisant des couleurs ou des styles différents. Repérer les points d'intersection trouvés à la question 3.
        \item Légender le graphique (titres des axes, noms des droites, unités).
    \end{itemize}

    \item \textbf{Lecture graphique et vérification.}
    À l'aide du graphique, lire approximativement les abscisses des points d'intersection. Vérifier qu'elles correspondent aux valeurs exactes calculées à la question 3. Observer l'ordre relatif des droites sur chaque intervalle.

    \item \textbf{Fonction « coût minimal » par intervalles.}
    Définir la fonction $f_{\min}$ qui associe à toute distance $x \ge 0$ le coût le plus bas parmi les trois offres. L'exprimer sous forme d'une fonction définie par intervalles (fonction affine par morceaux). Tracer sa courbe représentative en traits épais sur le graphique de la question 4.

    \item \textbf{Note de décision finale.}
    Rédiger une note adressée au président de l'association, structurée ainsi :
    \begin{quote}
        \textbf{Objet : Choix du transporteur pour la sortie au Parc de la Bénoué.}
        
        \textbf{Analyse :} [Rappeler les trois fonctions et les seuils trouvés.]
        
        \textbf{Recommandation :}
        \begin{itemize}
            \item Si le kilométrage prévu est \textbf{inférieur à [valeur]} km : choisir le transporteur \dots\ (justification).
            \item Si le kilométrage prévu est \textbf{compris entre [valeur] et [valeur]} km : choisir le transporteur \dots\ (justification).
            \item Si le kilométrage prévu est \textbf{supérieur à [valeur]} km : choisir le transporteur \dots\ (justification).
        \end{itemize}
        \textbf{Conclusion :} En fonction du programme retenu (400 km, 800 km ou 1\,500 km), indiquer le choix optimal et son coût.
    \end{quote}
    La note doit être signée par les membres du groupe.
\end{enumerate}

\courssub{Ressources à mobiliser}
\begin{aretenir}[Rappels de cours -- Applications affines]
\textbf{Fonction affine :} $f(x) = ax + b$ avec $a,b \in \mathbb{R}$.
\begin{itemize}
    \item Coefficient directeur : $a$ (taux d'accroissement). Ordonnée à l'origine : $b = f(0)$.
    \item Si $a > 0$, $f$ est strictement croissante sur $\mathbb{R}$.
    \item Si $a < 0$, $f$ est strictement décroissante sur $\mathbb{R}$.
    \item Si $a = 0$, $f$ est constante (fonction constante $f(x) = b$).
\end{itemize}
\textbf{Fonction linéaire :} $f(x) = ax$ (cas particulier $b=0$). La droite passe par l'origine $O(0,0)$.
\textbf{Image :} $y = f(x)$ se calcule en remplaçant $x$ par sa valeur.
\textbf{Antécédent :} Résoudre $f(x) = k \iff ax + b = k \iff x = \frac{k-b}{a}$ (si $a \ne 0$).
\textbf{Intersection de deux droites $y = a_1x+b_1$ et $y = a_2x+b_2$ ($a_1 \ne a_2$) :} résoudre $a_1x+b_1 = a_2x+b_2 \iff x = \frac{b_2-b_1}{a_1-a_2}$.
\textbf{Graphique :} Dans un repère orthonormé, une fonction affine se représente par une droite. Deux points suffisent pour la tracer (ex. $x=0$ et $x$ une valeur commode).
\end{aretenir}

\begin{propriete}[Fonction définie par intervalles (affine par morceaux)]
Une fonction définie par des expressions affines différentes sur des intervalles disjoints qui couvrent le domaine de définition est appelée \textbf{fonction affine par morceaux} (ou fonction par intervalles). Sa courbe représentative est constituée de segments de droites (ou demi-droites) joints ou non aux bornes des intervalles.
\end{propriete}

\courssub{Démarche et corrigé commenté}
\begin{exoresolu}[Corrigé détaillé de l'activité d'intégration]
\textbf{1. Modélisation}
Soit $x \in [0; +\infty[$ la distance en km, $y$ le coût en FCFA.
\begin{itemize}
    \item \textbf{Transporteur A (Voyages Express) :} Frais fixes 50\,000 + 200 FCFA/km. \quad $f_A(x) = 200x + 50\,000$.
    \item \textbf{Transporteur B (Transports Rapides) :} Frais fixes 30\,000 + 300 FCFA/km. \quad $f_B(x) = 300x + 30\,000$.
    \item \textbf{Transporteur C (Confort Plus) :} Forfait 100\,000 FCFA. \quad $f_C(x) = 100\,000$.
\end{itemize}
\underline{Nature :}
\begin{itemize}
    \item $f_A$ et $f_B$ sont des fonctions \textbf{affines} (coefficient directeur $a \ne 0$, ordonnée à l'origine $b \ne 0$).
    \item $f_C$ est une fonction \textbf{constante} (affine de coefficient directeur $a=0$).
\end{itemize}
\underline{Coefficients directeurs et ordonnées à l'origine :}
\begin{center}
\begin{tabular}{|c|c|c|c|}\hline
Fonction & Expression & Coeff. directeur $a$ & Ord. à l'origine $b$ \\ \hline
$f_A$ & $200x + 50\,000$ & $200$ & $50\,000$ \\ \hline
$f_B$ & $300x + 30\,000$ & $300$ & $30\,000$ \\ \hline
$f_C$ & $100\,000$ & $0$ & $100\,000$ \\ \hline
\end{tabular}
\end{center}
\emph{Interprétation :} $f_B$ augmente plus vite que $f_A$ ($300 > 200$). $f_C$ ne varie pas.

\textbf{2. Calculs numériques et tableau}
On calcule $f_A(x) = 200x+50\,000$, $f_B(x) = 300x+30\,000$, $f_C(x) = 100\,000$.

\begin{center}
\begin{tabular}{|c|c|c|c|}\hline
$x$ (km) & $f_A(x)$ (FCFA) & $f_B(x)$ (FCFA) & $f_C(x)$ (FCFA) \\ \hline
100 & $200\times100 + 50\,000 = 70\,000$ & $300\times100 + 30\,000 = 60\,000$ & 100\,000 \\ \hline
200 & $200\times200 + 50\,000 = 90\,000$ & $300\times200 + 30\,000 = 90\,000$ & 100\,000 \\ \hline
300 & $200\times300 + 50\,000 = 110\,000$ & $300\times300 + 30\,000 = 120\,000$ & 100\,000 \\ \hline
500 & $200\times500 + 50\,000 = 150\,000$ & $300\times500 + 30\,000 = 180\,000$ & 100\,000 \\ \hline
1\,000 & $200\times1\,000 + 50\,000 = 250\,000$ & $300\times1\,000 + 30\,000 = 330\,000$ & 100\,000 \\ \hline
\end{tabular}
\end{center}
\emph{Observation :} Pour 100 km, B est le moins cher (60\,000). Pour 200 km, A et B sont à égalité (90\,000). Pour 300 km, A devient moins cher que B, mais C (100\,000) est le moins cher. Pour 500 km et 1\,000 km, C reste le moins cher.

\textbf{3. Recherche algébrique des seuils}
On résout les équations $f_i(x) = f_j(x)$.

\begin{itemize}
    \item \textbf{Seuil $f_A = f_B$ :}
    \begin{align*}
        200x + 50\,000 &= 300x + 30\,000 \\
        50\,000 - 30\,000 &= 300x - 200x \\
        20\,000 &= 100x \\
        x &= 200
    \end{align*}
    Pour $x = 200$ km, les deux offres coûtent 90\,000 FCFA.
    \begin{itemize}
        \item Si $x < 200$ : $f_B(x) < f_A(x)$ (B moins cher que A).
        \item Si $x > 200$ : $f_A(x) < f_B(x)$ (A moins cher que B).
    \end{itemize}

    \item \textbf{Seuil $f_A = f_C$ :}
    \begin{align*}
        200x + 50\,000 &= 100\,000 \\
        200x &= 50\,000 \\
        x &= 250
    \end{align*}
    Pour $x = 250$ km, l'offre A coûte 100\,000 FCFA (comme C).
    \begin{itemize}
        \item Si $x < 250$ : $f_A(x) < 100\,000$ (A moins cher que C).
        \item Si $x > 250$ : $f_A(x) > 100\,000$ (C moins cher que A).
    \end{itemize}

    \item \textbf{Seuil $f_B = f_C$ :}
    \begin{align*}
        300x + 30\,000 &= 100\,000 \\
        300x &= 70\,000 \\
        x &= \frac{700}{3} \approx 233,33 \text{ km}
    \end{align*}
    Pour $x = \frac{700}{3}$ km, l'offre B coûte 100\,000 FCFA.
    \begin{itemize}
        \item Si $x < \frac{700}{3}$ : $f_B(x) < 100\,000$ (B moins cher que C).
        \item Si $x > \frac{700}{3}$ : $f_B(x) > 100\,000$ (C moins cher que B).
    \end{itemize}
\end{itemize}

\textbf{4. Représentation graphique}
\begin{popfigure}
\begin{tikzpicture}[scale=0.85]
% Échelle : 1 unité TikZ = 100 km en abscisse, 10 000 FCFA en ordonnée.
% Axe x : 0 à 5.5 (0 à 550 km). Axe y : 0 à 16 (0 à 160 000 FCFA).
% Axes
\draw[->, thick] (0,0) -- (5.8,0) node[right] {$x$ (km $\times 100$)}; % 5.8 pour marge
\draw[->, thick] (0,0) -- (0,16.5) node[above] {$y$ (FCFA $\times 10^4$)};

% Graduations axe x
\foreach \x in {0,1,2,3,4,5} {
    \draw (\x, -0.2) -- (\x, 0.2) node[below] {\small $\x00$};
}
% Graduations axe y
\foreach \y in {0,5,10,15} {
    \draw (-0.15, \y) -- (0.15, \y) node[left] {\small $\y$};
}
\node[below left] at (0,0) {0};

% Droite f_A : y = 200x + 50000 -> y/10000 = 0.02x + 5. En unités TikZ (x_tikz = x_km/100) : y_tikz = 2*x_tikz + 5
\draw[popblue, thick, domain=0:5.5] plot (\x, {2*\x + 5});
\node[popblue, above left] at (5.5, 16) {$f_A: y=200x+50\,000$};

% Droite f_B : y = 300x + 30000 -> y_tikz = 3*x_tikz + 3
\draw[poporange, thick, domain=0:5.5] plot (\x, {3*\x + 3});
\node[poporange, above right] at (4, 15) {$f_B: y=300x+30\,000$};

% Droite f_C : y = 100000 -> y_tikz = 10
\draw[popgreen, thick, dashed, domain=0:5.5] plot (\x, {10});
\node[popgreen, right] at (5.5, 10) {$f_C: y=100\,000$};

% Points d'intersection
% I_AB : x=200 -> x_tikz=2, y=9
\fill[poppink] (2, 9) circle (2.5pt);
\node[above left, poppink, font=\small] at (2, 9) {$I_{AB}(200; 90\,000)$};

% I_AC : x=250 -> x_tikz=2.5, y=10
\fill[poppink] (2.5, 10) circle (2.5pt);
\node[above right, poppink, font=\small] at (2.5, 10) {$I_{AC}(250; 100\,000)$};

% I_BC : x=700/3 ~ 233.33 -> x_tikz=2.333, y=10
\fill[poppink] (2.333, 10) circle (2.5pt);
\node[below right, poppink, font=\small] at (2.333, 10) {$I_{BC}(233,3; 100\,000)$};

% Courbe f_min (coût minimal) en traits épais noirs
% [0, 700/3] : f_B
\draw[black, ultra thick, domain=0:2.333] plot (\x, {3*\x + 3});
% [700/3, 250] : f_A
\draw[black, ultra thick, domain=2.333:2.5] plot (\x, {2*\x + 5});
% [250, 5.5] : f_C
\draw[black, ultra thick, domain=2.5:5.5] plot (\x, {10});
\node[black, above] at (1.5, 7.5) {\textbf{$f_{\min}$}};

\end{tikzpicture}
\legende{Représentation graphique des coûts des trois transporteurs. Échelle : 1 cm pour 100 km en abscisses, 1 cm pour 10\,000 FCFA en ordonnées. Les points d'intersection $I_{AB}$, $I_{BC}$, $I_{AC}$ sont repérés. La courbe épaisse noire représente la fonction « coût minimal » $f_{\min}$.}
\end{popfigure}

\textbf{5. Lecture graphique et vérification}
Sur le graphique, on lit :
\begin{itemize}
    \item $I_{AB}$ : abscisse 2 (soit 200 km), ordonnée 9 (soit 90\,000 FCFA). \textbf{Confirmé.}
    \item $I_{AC}$ : abscisse 2,5 (soit 250 km), ordonnée 10 (soit 100\,000 FCFA). \textbf{Confirmé.}
    \item $I_{BC}$ : abscisse $\approx 2,33$ (soit $\approx 233$ km), ordonnée 10 (soit 100\,000 FCFA). \textbf{Confirmé (valeur approchée de $700/3$).}
\end{itemize}
On observe l'ordre des droites (du plus bas au plus haut = du moins cher au plus cher) :
\begin{itemize}
    \item $[0 ; \frac{700}{3}[$ : $f_B$ (bas) $< f_A < f_C$ (haut). \textbf{B est optimal.}
    \item $]\frac{700}{3} ; 250[$ : $f_A$ (bas) $< f_C < f_B$ (haut). \textbf{A est optimal.}
    \item $]250 ; +\infty[$ : $f_C$ (bas) $< f_A < f_B$ (haut). \textbf{C est optimal.}
\end{itemize}

\textbf{6. Fonction « coût minimal » par intervalles}
La fonction $f_{\min}$ qui associe à $x \ge 0$ le coût le plus avantageux est une fonction affine par morceaux :
\[
f_{\min}(x) = 
\begin{cases}
300x + 30\,000 & \text{si } 0 \le x \le \frac{700}{3} \\[4pt]
200x + 50\,000 & \text{si } \frac{700}{3} \le x \le 250 \\[4pt]
100\,000 & \text{si } x \ge 250
\end{cases}
\]
Sa courbe (tracée en noir épais sur le graphique) est continue et convexe. Elle est constituée de trois tronçons : un segment de la droite $f_B$, un segment de la droite $f_A$, et un demi-droite horizontale (tronçon de $f_C$).

\textbf{7. Note de décision finale}

\begin{center}
\fbox{\begin{minipage}{0.92\linewidth}
\textbf{LYCÉE BILINGUE DE NKOLBISSON -- ASSOCIATION DES ÉLÈVES DE 3ÈME}\\[0.5em]
\textbf{Objet : Choix du transporteur pour la sortie au Parc National de la Bénoué.}\\[1em]
\textbf{Monsieur le Président,}\\[0.5em]
La commission Logistique a analysé les trois offres de transport reçues. Nous avons modélisé le coût total $y$ (en FCFA) en fonction de la distance $x$ (en km) par trois fonctions :
\begin{itemize}
    \item Transporteur A (Voyages Express) : $f_A(x) = 200x + 50\,000$ (affine, croissante).
    \item Transporteur B (Transports Rapides) : $f_B(x) = 300x + 30\,000$ (affine, croissante).
    \item Transporteur C (Confort Plus) : $f_C(x) = 100\,000$ (constante).
\end{itemize}
L'étude algébrique et graphique des intersections fait apparaître trois seuils décisifs :
\begin{itemize}
    \item $f_B = f_A$ pour $x = 200$ km (coût 90\,000 FCFA).
    \item $f_B = f_C$ pour $x = \frac{700}{3} \approx 233,3$ km (coût 100\,000 FCFA).
    \item $f_A = f_C$ pour $x = 250$ km (coût 100\,000 FCFA).
\end{itemize}

\textbf{Recommandation formelle :}
\begin{itemize}
    \item \textbf{Si le kilométrage prévu est inférieur à 233 km :} choisir le \textbf{Transporteur B (Transports Rapides)}. \\
    Justification : Sur $[0 ; \frac{700}{3}[$, la droite $f_B$ est la plus basse (coût le plus faible). Exemple : pour 100 km, B coûte 60\,000 FCFA contre 70\,000 pour A et 100\,000 pour C.
    
    \item \textbf{Si le kilométrage prévu est compris entre 233 km et 250 km :} choisir le \textbf{Transporteur A (Voyages Express)}. \\
    Justification : Sur $]\frac{700}{3} ; 250[$, $f_A$ passe sous $f_C$ et reste sous $f_B$. Exemple : pour 240 km, A coûte 98\,000 FCFA, B coûte 102\,000 FCFA, C coûte 100\,000 FCFA.
    
    \item \textbf{Si le kilométrage prévu est supérieur à 250 km :} choisir le \textbf{Transporteur C (Confort Plus)}. \\
    Justification : Au-delà de 250 km, le forfait illimité de C (100\,000 FCFA) devient imbattable. Exemple : pour 500 km, C coûte 100\,000 FCFA contre 150\,000 pour A et 180\,000 pour B.
\end{itemize}

\textbf{Application au programme prévu :}
\begin{itemize}
    \item Visite partielle ($\approx 400$ km) $> 250$ km $\rightarrow$ \textbf{Choix C} (Coût : 100\,000 FCFA).
    \item Visite complète ($\approx 800$ km) $> 250$ km $\rightarrow$ \textbf{Choix C} (Coût : 100\,000 FCFA).
    \item Extension Lac Tchad ($\approx 1\,500$ km) $> 250$ km $\rightarrow$ \textbf{Choix C} (Coût : 100\,000 FCFA).
\end{itemize}

\textbf{Conclusion :} Quel que soit le scénario retenu (400, 800 ou 1\,500 km), le \textbf{Transporteur C (Confort Plus)} est le choix optimal et le plus sûr (budget maîtrisé, assurances et péages inclus). Nous recommandons de signer le contrat avec Confort Plus dès cette semaine.

\vspace{1em}
\hfill Fait à Yaoundé, le \dots\ \\
\hfill \textbf{La Commission Logistique} \\
\hfill \textit{(Signatures des 3 membres du groupe)}
\end{minipage}}
\end{center}
\end{exoresolu}

\courssub{Bilan de l'activité}
Cette activité a permis de mettre en évidence les points suivants, essentiels pour le BEPC et la poursuite d'études :

\begin{itemize}
    \item \textbf{Modélisation :} Traduire un énoncé réel (tarifs transporteurs) en langage mathématique (fonctions affines/constantes) est la première étape indispensable. L'identification correcte du coefficient directeur (coût au km) et de l'ordonnée à l'origine (frais fixes) conditionne tout le reste.
    \item \textbf{Calcul algébrique vs Lecture graphique :} La résolution d'équations du premier degré ($f_i = f_j$) donne les seuils \textbf{exacts} (200 ; $700/3$ ; 250). Le graphique permet une \textbf{visualisation globale} immédiate de l'ordre des coûts et une vérification cohérente. Les deux approches sont complémentaires.
    \item \textbf{Sens de variation et décision :} Comparer des fonctions affines croissantes de coefficients directeurs différents (200 vs 300) montre que celle qui démarre plus haut (A : 50\,000) devient plus avantageuse après intersection car elle augmente moins vite. La fonction constante (C) agit comme un « plafond » : toute fonction croissante la coupe une fois et la dépasse ensuite.
    \item \textbf{Fonction affine par morceaux :} La recherche du coût minimal conduit naturellement à définir une fonction $f_{\min}$ par intervalles. C'est une application directe du point « Applications affines par intervalles » du programme.
    \item \textbf{Rédaction argumentée :} La note de décision impose de structurer sa pensée : \emph{constats chiffrés $\rightarrow$ seuils $\rightarrow$ intervalles $\rightarrow$ recommandation contextuelle}. C'est l'exercice type de l'épreuve de résolution de problèmes au BEPC.
    \item \textbf{Compétences transverses :} Travail en groupe, gestion du temps (45 min), communication orale (mise en commun), utilisation des unités (km, FCFA), choix d'échelles graphiques adaptées.
\end{itemize}

\begin{attention}[Pièges évités grâce à l'activité]
\begin{itemize}
    \item \textbf{Confondre $x$ et $y$ :} $x$ est la distance (variable indépendante), $y$ le coût (variable dépendante). Ne pas inverser les axes.
    \item \textbf{Oublier la condition $x \ge 0$ :} Une distance négative n'a pas de sens. Les droites ne sont tracées que pour $x \ge 0$.
    \item \textbf{Échelle graphique inadaptée :} Si on prend 1 cm pour 10 km en abscisse et 1 cm pour 1000 FCFA en ordonnée, le graphique tient sur 1 m de large et 25 cm de haut... mais les droites sont quasi confondues. Le choix 1 cm / 100 km et 1 cm / 10\,000 FCFA est lisible sur une feuille A4.
    \item \textbf{Ne pas vérifier l'ordre des droites \emph{entre} les intersections :} Il ne suffit pas de trouver $I_{AB}$, $I_{BC}$, $I_{AC}$ ; il faut tester un point dans chaque intervalle (ex. $x=100$, $x=240$, $x=300$) pour savoir qui est le minimum.
    \item \textbf{Réponse « C est le moins cher » sans nuance :} C n'est le moins cher qu'au-delà de 250 km. Pour un court trajet (ex. 50 km), C est le plus cher (100\,000 vs 60\,000 pour B). La réponse doit être \textbf{conditionnée par l'intervalle de $x$}.
    \item \textbf{Précision sur les bornes :} Aux valeurs exactes des intersections (200, $700/3$, 250), deux offres sont à égalité. La recommandation doit préciser si on inclut la borne dans un intervalle ou l'autre (ici, continuité de $f_{\min}$ assure l'égalité).
\end{itemize}
\end{attention}

\begin{savaistu}[Contexte camerounais -- Transport interurbain]
Au Cameroun, le transport interurbain (agences de voyage comme \emph{Voyages Express}, \emph{Transports Rapides}, \emph{Confort Plus}, \emph{General Voyage}, etc.) fonctionne souvent sur ce modèle : un prix de base + un coût au kilomètre, ou un forfait « tout compris ». Les associations scolaires, les groupes de paroissiens, les entreprises (séminaires) et les familles (deuils, mariages) font régulièrement ce type d'arbitrage. La maîtrise des applications affines permet d'économiser des centaines de milliers de FCFA sur l'année. C'est une compétence citoyenne et économique concrète.
\end{savaistu}

\newpage

\coursec{Méthodes et exercices résolus}

\coursec{Applications linéaires et affines}

\courssub{Définitions et vocabulaire : Linéaire vs Affine}

\begin{definition}[Application linéaire]
    Une application $f : \mathbb{R} \to \mathbb{R}$ est dite \textbf{linéaire} s'il existe un nombre réel $a$ tel que pour tout $x \in \mathbb{R}$, $f(x) = ax$.
    Le nombre $a$ est appelé \cle{coefficient directeur}. La représentation graphique est une droite passant par l'origine $O(0;0)$.
\end{definition}

\begin{definition}[Application affine]
    Une application $f : \mathbb{R} \to \mathbb{R}$ est dite \textbf{affine} s'il existe deux nombres réels $a$ et $b$ tels que pour tout $x \in \mathbb{R}$, $f(x) = ax + b$.
    Le nombre $a$ est le \cle{coefficient directeur} (ou taux d'accroissement), $b$ est l'\cle{ordonnée à l'origine}. La représentation graphique est une droite coupant l'axe des ordonnées au point $A(0;b)$.
    \begin{itemize}
        \item Toute application linéaire est affine (cas $b=0$).
        \item Une application affine avec $b \neq 0$ n'est \textbf{pas} linéaire.
        \item L'application constante $f(x) = b$ (cas $a=0$) est affine.
    \end{itemize}
\end{definition}

\begin{propriete}[Forme canonique]
    Toute application affine s'écrit de manière unique sous la forme $f(x) = ax + b$ avec $a, b \in \mathbb{R}$.
    \begin{itemize}
        \item $a = \dfrac{f(x_2) - f(x_1)}{x_2 - x_1}$ pour $x_1 \neq x_2$ (taux d'accroissement constant).
        \item $b = f(0)$.
    \end{itemize}
\end{propriete}

\begin{methode}[Reconnaître le type d'application]
    \textbf{Objectif :} Déterminer si une application $f$ est linéaire, affine ou ni l'une ni l'autre.
    \begin{enumerate}
        \item \textbf{Écrire l'expression algébrique} de $f(x)$ sous la forme $ax+b$ (développer et réduire si nécessaire).
        \item \textbf{Identifier les coefficients} $a$ (coefficient directeur) et $b$ (ordonnée à l'origine).
        \item \textbf{Conclure} :
        \begin{itemize}
            \item Si $b = 0$ : $f(x) = ax$ est une \cle{application linéaire}.
            \item Si $b \neq 0$ : $f(x) = ax+b$ est une \cle{application affine} (non linéaire).
            \item Si l'expression n'est pas de la forme $ax+b$ (ex: $x^2$, $\frac{1}{x}$, $\sqrt{x}$, $|x|$) : ce n'est \textbf{ni} une application linéaire \textbf{ni} affine.
        \end{itemize}
        \item \textbf{Réflexe « Tableau de valeurs »} : Si tu as un tableau, vérifie si les différences $f(x_2)-f(x_1)$ sont proportionnelles aux différences $x_2-x_1$ (taux d'accroissement constant $\Rightarrow$ affine). Si en plus $f(0)=0$ $\Rightarrow$ linéaire.
    \end{enumerate}
\end{methode}

\begin{exoresolu}[Type d'application : QCM corrigé]
    Parmi les applications suivantes, lesquelles sont affines ? Lesquelles sont linéaires ? Justifie.
    \begin{enumerate}
        \item $f_1(x) = 3x - 5$
        \item $f_2(x) = -2x$
        \item $f_3(x) = x^2 + 1$
        \item $f_4(x) = \dfrac{4x+2}{2}$
        \item $f_5(x) = 7$
    \end{enumerate}
    \tcblower
    \textbf{Solution détaillée}
    \begin{enumerate}
        \item $f_1(x) = 3x - 5$ : Forme $ax+b$ avec $a=3$, $b=-5 \neq 0$. \textbf{Affine (non linéaire)}.
        \item $f_2(x) = -2x$ : Forme $ax+b$ avec $a=-2$, $b=0$. \textbf{Linéaire} (donc aussi affine).
        \item $f_3(x) = x^2 + 1$ : Présence de $x^2$. Pas de la forme $ax+b$. \textbf{Ni linéaire, ni affine}.
        \item $f_4(x) = \dfrac{4x+2}{2} = 2x + 1$ : Après simplification, forme $ax+b$ avec $a=2$, $b=1 \neq 0$. \textbf{Affine (non linéaire)}.
        \item $f_5(x) = 7 = 0x + 7$ : Forme $ax+b$ avec $a=0$, $b=7 \neq 0$. \textbf{Affine (constante, non linéaire)}.
    \end{enumerate}
    \textbf{Conclusion} : Linéaires : $f_2$. Affines : $f_1, f_2, f_4, f_5$. Ni l'une ni l'autre : $f_3$.
\end{exoresolu}

\courssub{Images, antécédents et résolution d'équations}

\begin{definition}[Image et antécédent]
    Soit $f : x \mapsto ax+b$ une application affine.
    \begin{itemize}
        \item Pour $x_0 \in \mathbb{R}$, le nombre $f(x_0) = ax_0 + b$ est l'\cle{image} de $x_0$ par $f$.
        \item Pour $y_0 \in \mathbb{R}$, tout $x_0$ tel que $f(x_0) = y_0$ est un \cle{antécédent} de $y_0$ par $f$.
    \end{itemize}
\end{definition}

\begin{propriete}[Résolution de $f(x) = y_0$ et $f(x) < y_0$]
    Soit $f(x) = ax+b$ avec $a, b \in \mathbb{R}$.
    \begin{itemize}
        \item \textbf{Équation $f(x) = y_0$} :
        \begin{itemize}
            \item Si $a \neq 0$ : \textbf{Un unique antécédent} $x = \dfrac{y_0 - b}{a}$.
            \item Si $a = 0$ (fonction constante $f(x)=b$) :
            \begin{itemize}
                \item Si $y_0 = b$ : \textbf{Tous les réels} sont antécédents ($\mathbb{R}$).
                \item Si $y_0 \neq b$ : \textbf{Aucun antécédent} ($\varnothing$).
            \end{itemize}
        \end{itemize}
        \item \textbf{Inéquation $f(x) < y_0$ (ou $\le, >, \ge$)} : Résoudre $ax+b < y_0$ en isolant $x$.
        \textbf{Règle d'or} : Si on divise par $a < 0$, \textbf{le sens de l'inégalité change}.
    \end{itemize}
\end{propriete}

\begin{methode}[Images, antécédents et équations]
    \textbf{Objectif :} Maîtriser le vocabulaire et les calculs associés à $f(x)=ax+b$.
    \begin{enumerate}
        \item \textbf{Calculer l'image} d'un nombre $x_0$ : Remplacer $x$ par $x_0$ dans l'expression de $f$ et calculer $f(x_0)=ax_0+b$.
        \item \textbf{Chercher l'antécédent(s)} d'un nombre $y_0$ : Résoudre l'équation $f(x)=y_0 \Leftrightarrow ax+b=y_0$ (cf. propriété ci-dessus).
        \item \textbf{Résoudre $f(x) < k$ ou $f(x) \ge k$} : Résoudre $ax+b < k$ en isolant $x$. \textbf{Attention :} Si on divise par $a<0$, le sens de l'inégalité change !
        \item \textbf{Phrase de conclusion obligatoire} : « L'image de ... est ... » ou « L'antécédent de ... est ... » ou « L'ensemble des solutions est $S = \{...\}$ ».
    \end{enumerate}
\end{methode}

\begin{exoresolu}[Images, antécédents et inéquation]
    Soit l'application affine $f : x \mapsto -2x + 5$.
    \begin{enumerate}
        \item Calculer $f(3)$ et $f(-1)$.
        \item Déterminer l'antécédent de $11$.
        \item Résoudre l'inéquation $f(x) \le 1$.
    \end{enumerate}
    \tcblower
    \textbf{Solution détaillée}
    \begin{enumerate}
        \item \textbf{Images} :
        $f(3) = -2 \times 3 + 5 = -6 + 5 = -1$.
        $f(-1) = -2 \times (-1) + 5 = 2 + 5 = 7$.
        \textit{L'image de 3 est -1 ; l'image de -1 est 7.}
        \item \textbf{Antécédent de 11} : On résout $f(x) = 11$.
        \begin{align*}
            -2x + 5 &= 11 \\
            -2x &= 11 - 5 \\
            -2x &= 6 \\
            x &= \frac{6}{-2} = -3
        \end{align*}
        \textit{L'antécédent de 11 est -3.}
        \item \textbf{Inéquation $f(x) \le 1$} :
        \begin{align*}
            -2x + 5 &\le 1 \\
            -2x &\le 1 - 5 \\
            -2x &\le -4
        \end{align*}
        On divise par $-2$ (négatif), \textbf{le sens change} :
        $x \ge \dfrac{-4}{-2} \Leftrightarrow x \ge 2$.
        \textit{L'ensemble des solutions est $S = [2 ; +\infty[$}.
    \end{enumerate}
\end{exoresolu}

\courssub{Sens de variation et coefficient directeur}

\begin{propriete}[Variation d'une application affine]
    Soit $f(x) = ax+b$ définie sur $\mathbb{R}$ (ou un intervalle).
    \begin{itemize}
        \item Si $a > 0$ : $f$ est \textbf{strictement croissante}.
        \item Si $a < 0$ : $f$ est \textbf{strictement décroissante}.
        \item Si $a = 0$ : $f$ est \textbf{constante} (ni croissante, ni décroissante).
    \end{itemize}
    \textbf{Justification} : Pour $x_1 < x_2$, $\dfrac{f(x_2)-f(x_1)}{x_2-x_1} = a$. Le signe de $a$ donne le signe de $f(x_2)-f(x_1)$.
\end{propriete}

\begin{methode}[Variation d'une application affine]
    \textbf{Objectif :} Lier le signe du coefficient directeur $a$ à la monotonie de $f(x)=ax+b$.
    \begin{enumerate}
        \item \textbf{Identifier $a$} dans l'expression réduite $ax+b$.
        \item \textbf{Appliquer la règle} (valable sur $\mathbb{R}$ ou tout intervalle) selon le signe de $a$.
        \item \textbf{Tableau de variation} : Une ligne pour $x$, une pour $f$. Flèches $\nearrow$ (croissante) ou $\searrow$ (décroissante). Bornes $-\infty$ et $+\infty$ si domaine non borné.
    \end{enumerate}
\end{methode}

\begin{exoresolu}[Variation et lecture graphique]
    On donne $f(x) = -0,5x + 4$ et $g(x) = 3$.
    \begin{enumerate}
        \item Étudier la variation de $f$ et de $g$ sur $\mathbb{R}$.
        \item Construire le tableau de variation de $f$.
        \item Si $x$ varie de $-2$ à $6$, entre quelles valeurs varie $f(x)$ ?
    \end{enumerate}
    \tcblower
    \textbf{Solution détaillée}
    \begin{enumerate}
        \item $f(x) = -0,5x + 4 \Rightarrow a = -0,5 < 0$. Donc $f$ est \textbf{strictement décroissante} sur $\mathbb{R}$.
        $g(x) = 3 = 0x + 3 \Rightarrow a = 0$. Donc $g$ est \textbf{constante} sur $\mathbb{R}$.
        \item \textbf{Tableau de variation de $f$} :
        \begin{center}
        \begin{tabular}{|c|ccc|}\hline
        $x$ & $-\infty$ & & $+\infty$ \\ \hline
        $f(x)$ & $+\infty$ & $\searrow$ & $-\infty$ \\ \hline
        \end{tabular}
        \end{center}
        \item $f$ est décroissante. Pour $x_1 = -2$ et $x_2 = 6$ ($x_1 < x_2$), on a $f(x_1) > f(x_2)$.
        $f(-2) = -0,5(-2)+4 = 1+4 = 5$.
        $f(6) = -0,5(6)+4 = -3+4 = 1$.
        Quand $x$ varie de $-2$ à $6$, $f(x)$ varie de \textbf{5 à 1} (c'est-à-dire $f(x) \in [1 ; 5]$).
    \end{enumerate}
\end{exoresolu}

\courssub{Représentation graphique dans un repère orthonormé}

\begin{propriete}[Droite représentative]
    Dans un repère orthonormé $(O; \vec{\imath}, \vec{\jmath})$, la représentation graphique d'une application affine $f(x)=ax+b$ est une \textbf{droite} $\mathcal{C}_f$ :
    \begin{itemize}
        \item qui coupe l'axe des ordonnées au point $A(0; b)$ ;
        \item de coefficient directeur $a = \dfrac{\text{déplacement vertical}}{\text{déplacement horizontal}}$.
    \end{itemize}
    Réciproquement, toute droite non parallèle à l'axe des ordonnées est la représentation graphique d'une unique application affine.
\end{propriete}

\begin{methode}[Tracer et lire la droite représentative]
    \textbf{Objectif :} Passer de l'algébrique au graphique et inversement.
    \begin{enumerate}
        \item \textbf{Tracer la droite $\mathcal{C}_f$ de $f(x)=ax+b$} :
        \begin{itemize}
            \item Repérer le point $A(0 ; b)$ (intersection avec l'axe des ordonnées).
            \item Utiliser le coefficient directeur $a = \frac{\text{déplacement vertical}}{\text{déplacement horizontal}}$.
            \item Exemple : $a = \frac{3}{2}$ $\Rightarrow$ de $A$, avancer 2 à droite, 3 en haut $\to$ point $B$. $a = -1$ $\Rightarrow$ 1 à droite, 1 en bas.
            \item Tracer la droite $(AB)$.
        \end{itemize}
        \item \textbf{Lire une image graphiquement} : Pour $x_0$, monter à la droite, lire l'ordonnée.
        \item \textbf{Lire un antécédent graphiquement} : Pour $y_0$, couper la droite par l'horizontale $y=y_0$, lire l'abscisse.
        \item \textbf{Retrouver l'expression à partir du graphique} :
        \begin{itemize}
            \item Lire $b$ = ordonnée du point d'intersection avec l'axe des $y$ ($x=0$).
            \item Lire $a$ = $\frac{y_B - y_A}{x_B - x_A}$ pour deux points $A,B$ de la droite aux coordonnées entières lisibles.
        \end{itemize}
    \end{enumerate}
\end{methode}

\begin{popfigure}
\begin{tikzpicture}[scale=0.8, >=Stealth]
    % Axes
    \draw[->] (-1,0) -- (5,0) node[right] {$x$};
    \draw[->] (0,-1) -- (0,4) node[above] {$y$};
    % Grid
    \draw[popline, step=1] (-0.5,-0.5) grid (4.5,3.5);
    % Line f(x) = -0.5x + 3
    \draw[popblue, thick, domain=-0.5:5] plot (\x, {-0.5*\x+3});
    % Points
    \fill[popblue] (0,3) circle (2pt) node[left] {$A(0;b)$};
    \fill[popblue] (2,2) circle (2pt) node[above right] {$B$};
    % Slope triangle
    \draw[poporange, dashed] (0,3) -- (2,3) -- (2,2);
    \node[poporange] at (1, 3.2) {$\Delta x = 2$};
    \node[poporange] at (2.2, 2.5) {$\Delta y = -1$};
    % Labels
    \node[popblue] at (4.5, 1.5) {$\mathcal{C}_f : y = ax+b$};
    \node at (0,0) [below left] {$O$};
\end{tikzpicture}
\legende{Construction graphique de $f(x)=ax+b$ : point $A(0;b)$ et triangle de pente $a = \frac{\Delta y}{\Delta x}$.}
\end{popfigure}

\begin{exoresolu}[Graphique et expression]
    Dans un repère orthonormé $(O; \vec{\imath}, \vec{\jmath})$, on a tracé la droite $\mathcal{D}$ passant par $A(0; 2)$ et $B(3; -1)$.
    \begin{enumerate}
        \item Déterminer l'expression de l'application affine $f$ dont $\mathcal{D}$ est la représentation graphique.
        \item Résoudre graphiquement $f(x) = 0$ et $f(x) < 0$.
    \end{enumerate}
    \tcblower
    \textbf{Solution détaillée}
    \begin{enumerate}
        \item La droite coupe l'axe des ordonnées en $A(0;2)$, donc $b = 2$.
        Coefficient directeur $a = \dfrac{y_B - y_A}{x_B - x_A} = \dfrac{-1 - 2}{3 - 0} = \dfrac{-3}{3} = -1$.
        L'application affine est $f(x) = -x + 2$.
        \item \textbf{Graphiquement} :
        \begin{itemize}
            \item $f(x) = 0$ : Intersection de $\mathcal{D}$ avec l'axe des abscisses ($y=0$). Le point est $(2;0)$. Solution : $x=2$.
            \item $f(x) < 0$ : Portion de $\mathcal{D}$ \textbf{sous} l'axe des abscisses ($y<0$). Cela correspond aux abscisses $x > 2$.
        \end{itemize}
        \textbf{Vérification algébrique} : $-x+2 < 0 \Leftrightarrow -x < -2 \Leftrightarrow x > 2$. Cohérent.
    \end{enumerate}
\end{exoresolu}

\courssub{Applications affines par intervalles (Fonctions par morceaux)}

\begin{definition}[Fonction affine par intervalles]
    Une fonction $f$ est dite \textbf{affine par intervalles} (ou \textbf{fonction par morceaux}) si son ensemble de définition $D$ se partitionne en intervalles $I_1, I_2, \dots, I_n$ tels que la restriction de $f$ à chaque $I_k$ est une application affine (forme $a_k x + b_k$).
    On l'écrit souvent sous la forme :
    \[
    f(x) = \begin{cases}
    a_1 x + b_1 & \text{si } x \in I_1 \\
    a_2 x + b_2 & \text{si } x \in I_2 \\
    \vdots & \\
    a_n x + b_n & \text{si } x \in I_n
    \end{cases}
    \]
\end{definition}

\begin{methode}[Étudier une fonction définie par intervalles]
    \textbf{Objectif :} Gérer les fonctions dont la formule change selon $x$.
    \begin{enumerate}
        \item \textbf{Identifier les intervalles de définition} (ex: $x < 0$, $0 \le x \le 10$, $x > 10$). Vérifier la couverture du domaine.
        \item \textbf{Pour chaque intervalle} : Appliquer les méthodes précédentes (expression, variation, graphique) \textbf{uniquement sur cet intervalle}.
        \item \textbf{Points de jonction} : Vérifier la continuité (les formules donnent-elles la même valeur à la frontière ?).
        \item \textbf{Tableau de variation global} : Rassembler les variations de chaque morceau. Attention aux sauts (discontinuités) : double valeur ou trait de suspension.
        \item \textbf{Résolution $f(x)=k$} : Résoudre sur \textbf{chaque intervalle} et ne garder que les solutions appartenant à l'intervalle concerné.
    \end{enumerate}
\end{methode}

\begin{savaistu}[Tarification par tranches au Cameroun]
    De nombreux services au Cameroun utilisent une tarification affine par intervalles : forfaits Internet \textbf{MTN} / \textbf{Orange} (premiers Go à prix fixe, puis prix au Go), facture d'électricité \textbf{ENEO} (tranches de kWh), taxi (prix à la prise en charge + prix au km), location de salle (forfait + prix par invité). Modéliser ces situations permet de choisir l'offre la plus avantageuse.
\end{savaistu}

\begin{exoresolu}[Tarification par tranches (Contexte MTN/Orange)]
    Un forfait « Data » coûte 1000 F pour les 2 premiers Go, puis 300 F par Go supplémentaire.
    Soit $x$ le volume consommé en Go ($x \ge 0$) et $P(x)$ le prix en Francs CFA.
    \begin{enumerate}
        \item Exprimer $P(x)$ en fonction de $x$ (définition par intervalles).
        \item Calculer $P(1)$, $P(2)$, $P(5)$.
        \item Quel volume a été consommé si la facture est 2200 F ?
        \item Tracer la courbe représentative pour $x \in [0; 6]$.
    \end{enumerate}
    \tcblower
    \textbf{Solution détaillée}
    \begin{enumerate}
        \item Pour $0 \le x \le 2$ : Prix fixe 1000 F. $P(x) = 1000$.
        Pour $x > 2$ : 1000 F + 300 F $\times$ (Go au-delà de 2) = $1000 + 300(x-2) = 300x + 400$.
        \[
        P(x) = \begin{cases}
        1000 & \text{si } 0 \le x \le 2 \\
        300x + 400 & \text{si } x > 2
        \end{cases}
        \]
        \item $P(1) = 1000$ F (1er intervalle). $P(2) = 1000$ F (1er intervalle, borne incluse).
        $P(5)$ : $5 > 2$, 2ème intervalle. $P(5) = 300 \times 5 + 400 = 1500 + 400 = 1900$ F.
        \item Facture 2200 F $> 1000$ F $\Rightarrow$ on est dans le 2ème intervalle ($x > 2$).
        On résout $300x + 400 = 2200 \Leftrightarrow 300x = 1800 \Leftrightarrow x = 6$.
        Vérification : $6 > 2$ (OK). \textbf{Volume consommé : 6 Go.}
        \item \textbf{Graphique} :
        \begin{itemize}
            \item Segment horizontal $y=1000$ de $x=0$ à $x=2$ (points pleins aux extrémités).
            \item Demi-droite $y=300x+400$ pour $x \ge 2$. Point de départ $(2; 1000)$ (point plein, continuité assurée car $300\times2+400=1000$). Passe par $(4; 1600)$, $(6; 2200)$.
        \end{itemize}
    \end{enumerate}
\end{exoresolu}

\begin{popfigure}
\begin{tikzpicture}[x=1cm, y=0.0025cm, scale=0.7, >=Stealth]
    % Axes
    \draw[->] (-0.5,0) -- (6.5,0) node[right] {$x$ (Go)};
    \draw[->] (0,-100) -- (0,2400) node[above] {$P(x)$ (F CFA)};
    % Grid
    \draw[popline, step=1] (0,0) grid (6,2200);
    % Tick labels
    \foreach \x in {1,2,3,4,5,6} \draw (\x, -50) -- (\x, 50) node[below] {\x};
    \foreach \y in {500,1000,1500,2000} \draw (-0.2, \y) -- (0.2, \y) node[left] {\y};
    % Piece 1: horizontal
    \draw[popblue, very thick] (0,1000) -- (2,1000);
    \fill[popblue] (0,1000) circle (3pt);
    \fill[popblue] (2,1000) circle (3pt);
    % Piece 2: slope
    \draw[poporange, very thick] (2,1000) -- (6,2200);
    \fill[poporange] (2,1000) circle (3pt);
    \fill[poporange] (4,1600) circle (3pt);
    \fill[poporange] (6,2200) circle (3pt);
    % Labels
    \node[popblue] at (1, 1100) {$P(x)=1000$};
    \node[poporange] at (5, 1900) {$P(x)=300x+400$};
\end{tikzpicture}
\legende{Représentation graphique de la fonction tarifaire MTN/Orange (continuité en $x=2$).}
\end{popfigure}

\courssub{Modélisation, problèmes de synthèse}

\begin{methode}[Résoudre un problème concret par modélisation affine]
    \textbf{Objectif :} Traduire une situation réelle en fonction affine (simple ou par morceaux), résoudre et interpréter.
    \begin{enumerate}
        \item \textbf{Variables} : Définir clairement $x$ (indépendant, unité) et $y$ ou $f(x)$ (dépendant, unité).
        \item \textbf{Analyse} : Identifier les seuils, les coûts fixes (ordonnée à l'origine), les coûts unitaires (pente/coefficient directeur).
        \item \textbf{Modèle} : Écrire $f(x) = ax+b$ (ou par morceaux). Préciser le domaine de définition.
        \item \textbf{Résolution} : Calculs d'images, antécédents, inéquations, point d'intersection de deux droites (comparaison d'offres).
        \item \textbf{Interprétation} : Répondre par une phrase complète dans le contexte (« Il faut vendre au moins 50 poulets pour être rentable »).
        \item \textbf{Validation} : Vérifier la cohérence (prix positif, nombre entier si objets discrets).
    \end{enumerate}
\end{methode}

\begin{exoresolu}[Choix entre deux offres : Location de salle des fêtes]
    Pour un mariage à Yaoundé, deux salles proposent :
    \begin{itemize}
        \item \textbf{Salle A (Bastos)} : 150 000 F de location + 2 500 F par invité.
        \item \textbf{Salle B (Mvog-Ada)} : 250 000 F de location + 1 500 F par invité.
    \end{itemize}
    On note $x$ le nombre d'invités ($x \in \mathbb{N}, x \ge 1$). $P_A(x)$ et $P_B(x)$ les coûts totaux.
    \begin{enumerate}
        \item Exprimer $P_A(x)$ et $P_B(x)$. Sont-elles affines ? Linéaires ?
        \item Pour quel nombre d'invités les coûts sont-ils égaux ?
        \item Déterminer l'intervalle de $x$ où la Salle A est moins chère que la Salle B.
        \item Si le budget est limité à 600 000 F, quelle salle choisir pour maximiser le nombre d'invités ?
    \end{enumerate}
    \tcblower
    \textbf{Solution détaillée}
    \begin{enumerate}
        \item $P_A(x) = 2500x + 150000$. $P_B(x) = 1500x + 250000$.
        Ce sont des applications \textbf{affines} (forme $ax+b$, $b \neq 0$). Elles ne sont \textbf{pas linéaires} ($b \neq 0$).
        \item Égalité des coûts : $P_A(x) = P_B(x)$.
        \begin{align*}
            2500x + 150000 &= 1500x + 250000 \\
            1000x &= 100000 \\
            x &= 100
        \end{align*}
        Pour \textbf{100 invités}, les deux salles coûtent le même prix (400 000 F).
        \item Salle A moins chère : $P_A(x) < P_B(x)$.
        $2500x + 150000 < 1500x + 250000 \Leftrightarrow 1000x < 100000 \Leftrightarrow x < 100$.
        Comme $x \in \mathbb{N}^*$, la Salle A est moins chère pour \textbf{1 à 99 invités}.
        \item Budget 600 000 F. On cherche le max $x$ tel que $P(x) \le 600000$.
        \textbf{Salle A} : $2500x + 150000 \le 600000 \Rightarrow 2500x \le 450000 \Rightarrow x \le 180$. Max 180 invités.
        \textbf{Salle B} : $1500x + 250000 \le 600000 \Rightarrow 1500x \le 350000 \Rightarrow x \le 233,33$. Max 233 invités (entier).
        \textbf{Choix : La Salle B (Mvog-Ada) permet d'inviter 233 personnes contre 180 pour la Salle A.}
    \end{enumerate}
\end{exoresolu}

\begin{methode}[Approche globale : Modéliser, Tracer, Décider]
    \textbf{Objectif :} Mobiliser toutes les compétences de la leçon sur un problème complet.
    \begin{enumerate}
        \item \textbf{Lire et annoter} : Surligner les seuils, les prix unitaires, les abonnements.
        \item \textbf{Construire la fonction par morceaux} : Une formule par tranche. Vérifier la continuité aux seuils.
        \item \textbf{Tableau de variation global} : Croissante par morceaux (prix unitaires positifs).
        \item \textbf{Résolution equation/inéquation} : Par tranche. Ne pas oublier la condition d'appartenance à l'intervalle.
        \item \textbf{Représentation graphique} : Repère adapté (échelle), points de rupture, droites partielles.
        \item \textbf{Conclusion argumentée} : Phrase finale répondant à la question posée.
    \end{enumerate}
\end{methode}

\begin{exoresolu}[Facture ENEO : Tranches de consommation]
    Le tarif résidentiel ENEO (simplifié) est :
    \begin{itemize}
        \item Tranche 1 : 0 à 110 kWh : 50 F/kWh.
        \item Tranche 2 : 111 à 400 kWh : 79 F/kWh.
        \item Tranche 3 : Au-delà de 400 kWh : 95 F/kWh.
    \end{itemize}
    L'abonnement mensuel est de 1 000 F. Soit $x$ la consommation en kWh ($x \ge 0$) et $F(x)$ la facture totale.
    \begin{enumerate}
        \item Écrire $F(x)$ sous forme d'application affine par intervalles.
        \item Calculer la facture pour 100 kWh, 200 kWh et 500 kWh.
        \item Un ménage a une facture de 35 550 F. Quelle a été sa consommation ?
        \item Tracer la courbe de $F$ pour $x \in [0; 600]$.
    \end{enumerate}
    \tcblower
    \textbf{Solution détaillée}
    \begin{enumerate}
        \item \textbf{Construction par tranches} (on cumule) :
        \begin{itemize}
            \item $[0; 110]$ : $F(x) = 1000 + 50x$.
            \item $]110; 400]$ : Coût tranche 1 = $1000 + 50\times110 = 6500$ F.
            $F(x) = 6500 + 79(x-110) = 79x - 2190$.
            \item $]400; +\infty[$ : Coût tranches 1+2 = $6500 + 79\times290 = 6500 + 22910 = 29410$ F.
            $F(x) = 29410 + 95(x-400) = 95x - 8590$.
        \end{itemize}
        \[
        F(x) = \begin{cases}
        50x + 1000 & \text{si } 0 \le x \le 110 \\
        79x - 2190 & \text{si } 110 < x \le 400 \\
        95x - 8590 & \text{si } x > 400
        \end{cases}
        \]
        \textit{Vérif continuité :} $x=110 \to 50(110)+1000=6500$ ; $79(110)-2190=6500$. OK. $x=400 \to 79(400)-2190=29410$ ; $95(400)-8590=29410$. OK.
        \item $F(100) = 50(100)+1000 = 6000$ F.
        $F(200)$ : $200 \in ]110;400] \Rightarrow 79(200)-2190 = 15800-2190 = 13610$ F.
        $F(500)$ : $500 > 400 \Rightarrow 95(500)-8590 = 47500-8590 = 38910$ F.
        \item Facture 35 550 F.
        \begin{itemize}
            \item Tranche 1 max : $F(110)=6500 < 35550$.
            \item Tranche 2 max : $F(400)=29410 < 35550$.
            \item $\Rightarrow$ Tranche 3 ($x>400$). Résoudre $95x - 8590 = 35550$.
            $95x = 44140 \Rightarrow x = \frac{44140}{95} = 464,63...$
        \end{itemize}
        La consommation est d'environ \textbf{465 kWh} (arrondi au kWh supérieur si compteur entier, ou 464,6 kWh exact).
        \item \textbf{Graphique} : Repère $(O; \vec{\imath}, \vec{\jmath})$. 1 cm pour 50 kWh (abscisses), 1 cm pour 5000 F (ordonnées).
        Points clés : $(0;1000)$, $(110;6500)$, $(400;29410)$, $(600; 95\times600-8590=48410)$.
        Trois segments de droites de pentes croissantes (50, 79, 95) reliés entre eux (fonction continue, convexe).
    \end{enumerate}
\end{exoresolu}

\begin{popfigure}
\begin{tikzpicture}[x=0.012cm, y=0.00024cm, >=Stealth]
    % Axes
    \draw[->] (-10,0) -- (620,0) node[right] {$x$ (kWh)};
    \draw[->] (0,-500) -- (0,50000) node[above] {$F(x)$ (F CFA)};
    % Grid lines (light)
    \draw[popline, step=100] (0,0) grid (600,48000);
    % Tick labels x
    \foreach \x in {100,200,300,400,500,600} \draw (\x, -1000) -- (\x, 1000) node[below] {\x};
    % Tick labels y
    \foreach \y in {10000,20000,30000,40000} \draw (-20, \y) -- (20, \y) node[left] {\y};
    % Piece 1
    \draw[popblue, very thick] (0,1000) -- (110,6500);
    \fill[popblue] (0,1000) circle (3pt) node[left] {1000};
    \fill[popblue] (110,6500) circle (3pt);
    % Piece 2
    \draw[poporange, very thick] (110,6500) -- (400,29410);
    \fill[poporange] (400,29410) circle (3pt);
    % Piece 3
    \draw[popgreen, very thick] (400,29410) -- (600,48410);
    \fill[popgreen] (600,48410) circle (3pt);
    % Vertical dashed lines at thresholds
    \draw[dashed, popmuted] (110,0) -- (110,6500);
    \draw[dashed, popmuted] (400,0) -- (400,29410);
    \node[popmuted, below] at (110, -1500) {110};
    \node[popmuted, below] at (400, -1500) {400};
    % Labels
    \node[popblue] at (80, 4000) {T1: $50x+1000$};
    \node[poporange] at (250, 18000) {T2: $79x-2190$};
    \node[popgreen] at (500, 40000) {T3: $95x-8590$};
\end{tikzpicture}
\legende{Courbe représentative de la facture ENEO $F(x)$ (fonction continue, convexe, croissante par morceaux).}
\end{popfigure}

\begin{attention}[Les 5 erreurs qui coûtent des points au BEPC]
    \begin{enumerate}
        \item \textbf{Confondre Linéaire et Affine} : Dire « $f(x)=2x+3$ est linéaire » (FAUX, $b\neq0$). Dire « $f(x)=5$ n'est pas affine » (FAUX, $a=0, b=5$).
        \item \textbf{Oublier le changement de sens} en divisant une inéquation par un nombre \textbf{négatif} ($a<0$). Ex: $-2x < 4 \Rightarrow x > -2$ (pas $x < -2$).
        \item \textbf{Antécédent d'une fonction constante} : Pour $f(x)=3$, l'antécédent de 3 est $\mathbb{R}$ (pas « 0 » ou « 3 »), l'antécédent de 5 est $\varnothing$ (pas « n'existe pas » sans notation ensemble).
        \item \textbf{Fonction par morceaux : Oublier la condition d'intervalle}. Résoudre $300x+400=1000 \Rightarrow x=2$ mais dire « $x=2$ est solution pour $x>2$ » (FAUX, $2 \notin ]2;+\infty[$). Vérifier \textbf{toujours} $x \in I_k$.
        \item \textbf{Graphique : Échelle non respectée ou points non reportés}. Tracer une droite qui ne passe pas par $(0;b)$ ou dont la pente visuelle ne correspond pas au signe de $a$. Ne pas marquer les points de rupture (plein/vide) pour les fonctions par morceaux.
    \end{enumerate}
\end{attention}

\begin{aretenir}[Bilan de la leçon]
    \begin{itemize}
        \item \textbf{Linéaire} : $f(x)=ax$ (passe par origine). \textbf{Affine} : $f(x)=ax+b$ (droite quelconque non verticale).
        \item \textbf{Image} : $f(x_0)=ax_0+b$. \textbf{Antécédent} : résoudre $ax+b=y_0$.
        \item \textbf{Variation} : Signe de $a$ ($+$ croissante, $-$ décroissante, $0$ constante).
        \item \textbf{Graphique} : Droite $\mathcal{C}_f$. $b$ = ordonnée à l'origine. $a$ = pente.
        \item \textbf{Par morceaux} : Une formule par intervalle. Vérifier l'appartenance à l'intervalle pour chaque solution. Tracer chaque morceau sur son intervalle (point plein si borne incluse, vide si exclue).
        \item \textbf{Modélisation} : Identifier fixe ($b$) et variable ($a$). Comparer offres = intersection droites. Optimiser budget = inéquation.
    \end{itemize}
\end{aretenir}

\newpage

\coursec{Exercices}

\courssub{Je vérifie mes connaissances}

\begin{exercice}[QCM : linéaire ou affine ?]
Pour chacune des applications suivantes, indique si elle est \textbf{linéaire}, \textbf{affine non linéaire}, ou \textbf{ni l'une ni l'autre}. Justifie ta réponse en précisant, le cas échéant, les nombres $a$ et $b$ de l'écriture $f(x)=ax+b$.
\begin{enumerate}
\item $f(x)=5x$
\item $g(x)=-2x+3$
\item $h(x)=x^2+1$
\item $k(x)=\dfrac{x}{4}-7$
\item $l(x)=\dfrac{3}{x}$
\end{enumerate}
\end{exercice}

\begin{exercice}[Vrai ou faux ?]
Réponds par \textbf{Vrai} ou \textbf{Faux} en justifiant chaque réponse.
\begin{enumerate}
\item L'application affine $f(x)=-4x+1$ est croissante sur $\mathbb{R}$.
\item L'image de $2$ par l'application $g(x)=3x-5$ est $1$.
\item L'antécédent de $0$ par l'application $h(x)=2x+6$ est $-3$.
\item La représentation graphique d'une application linéaire est une droite qui passe par l'origine du repère.
\item Deux applications affines ayant le même coefficient directeur ont des représentations graphiques parallèles.
\end{enumerate}
\end{exercice}

\begin{exercice}[Images et antécédents : calculs directs]
Soit l'application affine $f$ définie par $f(x)=-3x+7$.
\begin{enumerate}
\item Calcule les images de $0$, de $-2$ et de $\dfrac{5}{3}$ par $f$.
\item Détermine les antécédents de $7$, de $1$ et de $-5$ par $f$.
\item Le point $A(4\,;-5)$ appartient-il à la représentation graphique de $f$ ? Justifie.
\end{enumerate}
\end{exercice}

\begin{exercice}[Lecture graphique]
Dans un repère orthonormé $(O\,;I,J)$, la droite $(d)$ représente une application affine $f$ et passe par les points $A(0\,;2)$ et $B(3\,;0)$.
\begin{enumerate}
\item Lis graphiquement $f(0)$ et $f(3)$.
\item Lis graphiquement l'antécédent de $1$ par $f$.
\item Sachant que $f(x)=ax+b$, détermine $b$, puis calcule $a$ à l'aide du point $B$.
\end{enumerate}
\end{exercice}

\courssub{Je m'entraîne}

\begin{exercice}[Sens de variation et inéquation]
On considère l'application affine $f$ définie sur $\mathbb{R}$ par $f(x)=-3x+7$.
\begin{enumerate}
\item Justifie, à l'aide du signe du coefficient directeur, que $f$ est décroissante sur $\mathbb{R}$.
\item Compare, sans les calculer, $f(12)$ et $f(15)$. Justifie.
\item Résous dans $\mathbb{R}$ l'inéquation $f(x)>1$.
\item Vérifie ta réponse à la question 3 en calculant $f(2)$ et $f(3)$.
\end{enumerate}
\end{exercice}

\begin{exercice}[Détermination d'une application affine]
\begin{enumerate}
\item Détermine l'application linéaire $g$ telle que $g(4)=10$.
\item Détermine l'application affine $f$ telle que $f(1)=3$ et $f(3)=9$.
\item La droite représentative de $f$ coupe-t-elle l'axe des ordonnées au-dessus ou en dessous de l'origine ? Justifie.
\end{enumerate}
\end{exercice}

\begin{exercice}[Représentations graphiques]
Le plan est muni d'un repère orthonormé $(O\,;I,J)$ d'unité $1$ cm.
\begin{enumerate}
\item Trace la droite $(D_1)$ représentative de $f(x)=2x-1$ et la droite $(D_2)$ représentative de $g(x)=-x+5$.
\item Les droites $(D_1)$ et $(D_2)$ se coupent en un point $I$. Calcule les coordonnées exactes de $I$.
\item Détermine, par le calcul, l'ensemble des valeurs de $x$ pour lesquelles le point de $(D_1)$ d'abscisse $x$ est situé \emph{au-dessus} du point de $(D_2)$ de même abscisse.
\item Contrôle ta réponse à la question 3 par une lecture graphique.
\end{enumerate}
\end{exercice}

\begin{exercice}[Application affine par intervalles]
Un parking du centre-ville de Yaoundé applique le tarif suivant pour une durée de stationnement $x$ (en heures) :
\begin{itemize}
\item pour $0\le x\le 2$ : $200$ F l'heure ;
\item pour $2<x\le 6$ : un forfait de $400$ F plus $100$ F par heure au-delà de la deuxième heure.
\end{itemize}
\begin{enumerate}
\item Calcule le prix payé pour $1$ h $30$ min, puis pour $5$ h de stationnement.
\item Exprime le prix $P(x)$ en fonction de $x$, en distinguant les deux intervalles.
\item Représente graphiquement $P$ pour $x$ compris entre $0$ et $6$ (1 cm pour 1 h en abscisse, 1 cm pour 100 F en ordonnée).
\end{enumerate}
\end{exercice}

\courssub{Je me prépare au BEPC}

\begin{exercice}[Type BEPC 1 : étude complète d'une application affine]
On considère l'application affine $f$ définie sur $\mathbb{R}$ par $f(x)=-3x+7$. Le plan est muni d'un repère orthonormé $(O\,;I,J)$ d'unité $1$ cm.
\begin{enumerate}
\item \begin{enumerate}
\item Calcule l'image de $0$ et l'image de $3$ par $f$.
\item Détermine l'antécédent de $-2$ par $f$.
\end{enumerate}
\item Démontre que $f$ est décroissante sur $\mathbb{R}$.
\item Résous dans $\mathbb{R}$ l'inéquation $f(x)>1$. Donne l'ensemble des solutions sous forme d'un intervalle.
\item Trace soigneusement la droite $(\mathcal{C}_f)$, représentation graphique de $f$, dans le repère $(O\,;I,J)$.
\item À l'aide du graphique, lis l'antécédent de $-2$ par $f$ (on laissera les traits de lecture apparents). Compare avec le résultat de la question 1.b.
\item Soit $g$ l'application linéaire définie par $g(x)=2x$. Trace sa représentation graphique $(\mathcal{C}_g)$ dans le même repère, puis calcule les coordonnées du point d'intersection de $(\mathcal{C}_f)$ et $(\mathcal{C}_g)$.
\end{enumerate}
\end{exercice}

\begin{exercice}[Type BEPC 2 : seuil de rentabilité d'une usine de briquettes]
Une usine de la région du Littoral fabrique des briquettes de terre compressée. Pour une production de $x$ briquettes par jour :
\begin{itemize}
\item les charges fixes (loyer, salaires, entretien) s'élèvent à $200\,000$ F ;
\item la fabrication de chaque briquette coûte $150$ F (matière première, énergie) ;
\item chaque briquette est vendue $250$ F.
\end{itemize}
On note $C(x)$ le coût total de production et $R(x)$ la recette (revenu) pour $x$ briquettes produites et vendues.
\begin{enumerate}
\item Exprime $C(x)$ et $R(x)$ en fonction de $x$. Précise la nature (linéaire ou affine) de chacune de ces applications.
\item Calcule $C(1\,000)$ et $R(1\,000)$. L'usine fait-elle un bénéfice pour $1\,000$ briquettes vendues ?
\item Représente graphiquement $C$ et $R$ dans un même repère orthogonal pour $x$ compris entre $0$ et $3\,000$ (1 cm pour 500 briquettes en abscisse, 1 cm pour $100\,000$ F en ordonnée).
\item On appelle \emph{seuil de rentabilité} le nombre de briquettes à partir duquel la recette devient supérieure ou égale au coût total.
\begin{enumerate}
\item Traduis cette condition par une inéquation d'inconnue $x$.
\item Résous cette inéquation et conclus : à partir de combien de briquettes vendues par jour l'usine réalise-t-elle un bénéfice ?
\item Retrouve ce résultat par lecture graphique (laisser les traits de construction apparents).
\end{enumerate}
\item Le directeur souhaite un bénéfice journalier d'au moins $100\,000$ F. Combien de briquettes, au minimum, doit-il produire et vendre chaque jour ?
\end{enumerate}
\end{exercice}

\begin{exercice}[Type BEPC 3 : tarification de l'eau par tranches]
Une société de distribution d'eau facture la consommation mensuelle $x$ (en m$^3$) d'un abonné selon les tranches suivantes, auxquelles s'ajoute un abonnement fixe de $1\,000$ F :
\begin{itemize}
\item tranche 1 : de $0$ à $10$ m$^3$ inclus, à $300$ F le m$^3$ ;
\item tranche 2 : au-delà de $10$ m$^3$ et jusqu'à $30$ m$^3$ inclus, à $500$ F le m$^3$ ;
\item tranche 3 : au-delà de $30$ m$^3$, à $800$ F le m$^3$.
\end{itemize}
On note $P(x)$ le prix total à payer, en francs, pour une consommation de $x$ m$^3$.
\begin{enumerate}
\item Vérifie qu'une consommation de $10$ m$^3$ est facturée $4\,000$ F au total, et qu'une consommation de $30$ m$^3$ est facturée $14\,000$ F au total.
\item Montre que $P$ s'écrit :
\[
P(x)=\begin{cases} 300x+1\,000 & \text{si } 0\le x\le 10\\ 500x-1\,000 & \text{si } 10<x\le 30\\ 800x-10\,000 & \text{si } x>30 \end{cases}
\]
\item Calcule $P(5)$, $P(15)$ et $P(40)$.
\item Représente graphiquement $P$ pour $x$ compris entre $0$ et $40$ (1 cm pour 5 m$^3$ en abscisse, 1 cm pour $5\,000$ F en ordonnée).
\item Observe ton graphique : peut-on le tracer sans lever le crayon ? Que remarques-tu aux points d'abscisses $10$ et $30$ ?
\item La famille Mbarga a payé $9\,000$ F ce mois-ci. Quelle a été sa consommation d'eau ? Justifie par le calcul en précisant la tranche utilisée.
\end{enumerate}
\end{exercice}



\newpage

\coursec{Corrigés des exercices}

\courssub{Je vérifie mes connaissances}

\begin{corrige}[Exercice 1 — QCM : linéaire ou affine ?]
\begin{enumerate}
\item $f(x)=5x$ est \textbf{linéaire} : elle s'écrit $f(x)=ax+b$ avec $a=5$ et $b=0$.
\item $g(x)=-2x+3$ est \textbf{affine non linéaire} : $a=-2$ et $b=3\neq 0$.
\item $h(x)=x^2+1$ n'est \textbf{ni linéaire ni affine} : le terme $x^2$ empêche toute écriture sous la forme $ax+b$.
\item $k(x)=\dfrac{x}{4}-7=\dfrac{1}{4}x-7$ est \textbf{affine non linéaire} : $a=\dfrac{1}{4}$ et $b=-7$.
\item $l(x)=\dfrac{3}{x}$ n'est \textbf{ni linéaire ni affine} : la variable $x$ est au dénominateur.
\end{enumerate}
\end{corrige}

\begin{corrige}[Exercice 2 — Vrai ou faux ?]
\begin{enumerate}
\item \textbf{Faux} : le coefficient directeur est $a=-4<0$, donc $f$ est \emph{décroissante} sur $\mathbb{R}$.
\item \textbf{Vrai} : $g(2)=3\times 2-5=6-5=1$.
\item \textbf{Vrai} : on résout $2x+6=0$, soit $2x=-6$, d'où $x=-3$.
\item \textbf{Vrai} : pour une application linéaire $f(x)=ax$, on a $f(0)=0$, donc sa représentation graphique passe par l'origine $O$.
\item \textbf{Vrai} : deux droites de même coefficient directeur sont parallèles (distinctes ou confondues).
\end{enumerate}
\end{corrige}

\begin{corrige}[Exercice 3 — Images et antécédents : calculs directs]
\begin{enumerate}
\item $f(0)=-3\times 0+7=7$ ; $f(-2)=-3\times(-2)+7=6+7=13$ ; $f\left(\dfrac{5}{3}\right)=-3\times\dfrac{5}{3}+7=-5+7=2$.
\item On résout $f(x)=y$ pour chaque valeur :
\begin{itemize}
\item $-3x+7=7 \iff -3x=0 \iff x=0$ : l'antécédent de $7$ est $0$ ;
\item $-3x+7=1 \iff -3x=-6 \iff x=2$ : l'antécédent de $1$ est $2$ ;
\item $-3x+7=-5 \iff -3x=-12 \iff x=4$ : l'antécédent de $-5$ est $4$.
\end{itemize}
\item On calcule $f(4)=-3\times 4+7=-12+7=-5$. Comme $f(4)=-5$, le point $A(4\,;-5)$ \textbf{appartient} à la représentation graphique de $f$.
\end{enumerate}
\end{corrige}

\begin{corrige}[Exercice 4 — Lecture graphique]
\begin{enumerate}
\item $A(0\,;2)\in(d)$ donc $f(0)=2$ ; $B(3\,;0)\in(d)$ donc $f(3)=0$.
\item On cherche le point de $(d)$ d'ordonnée $1$ : par lecture graphique, son abscisse est $1{,}5$. L'antécédent de $1$ est donc $1{,}5$.
\item $b=f(0)=2$ (ordonnée à l'origine). Ensuite $f(3)=3a+2=0$, donc $3a=-2$ et $a=-\dfrac{2}{3}$. Ainsi $f(x)=-\dfrac{2}{3}x+2$. Vérification : $f(1{,}5)=-\dfrac{2}{3}\times 1{,}5+2=-1+2=1$, ce qui confirme la lecture de la question 2.
\end{enumerate}
\end{corrige}

\courssub{Je m'entraîne}

\begin{corrige}[Exercice 5 — Sens de variation et inéquation]
\begin{enumerate}
\item Le coefficient directeur de $f$ est $a=-3$. Comme $a<0$, l'application affine $f$ est \textbf{décroissante} sur $\mathbb{R}$.
\item Comme $12<15$ et que $f$ est décroissante, l'ordre est renversé : $f(12)>f(15)$.
\item $f(x)>1 \iff -3x+7>1 \iff -3x>-6 \iff x<2$ (on divise par $-3<0$, le sens de l'inégalité change). L'ensemble des solutions est $]-\infty\,;2[$.
\item $f(2)=-3\times 2+7=1$ : la valeur $2$ est bien la frontière (non incluse). $f(3)=-9+7=-2<1$ : $3$ n'est pas solution, ce qui est cohérent puisque $3>2$. La réponse est confirmée.
\end{enumerate}
\end{corrige}

\begin{corrige}[Exercice 6 — Détermination d'une application affine]
\begin{enumerate}
\item $g$ est linéaire donc $g(x)=ax$. De $g(4)=4a=10$, on tire $a=\dfrac{10}{4}=\dfrac{5}{2}=2{,}5$. Donc $g(x)=2{,}5x$.
\item On cherche $f(x)=ax+b$. Le coefficient directeur est
\[ a=\frac{f(3)-f(1)}{3-1}=\frac{9-3}{2}=\frac{6}{2}=3. \]
Puis $f(1)=3\times 1+b=3$ donne $b=0$. Donc $f(x)=3x$ : c'est en fait une application \emph{linéaire}.
\item Comme $b=0$, la droite représentative de $f$ passe par l'origine $O$ : elle coupe l'axe des ordonnées \emph{en} l'origine, ni au-dessus ni en dessous.
\end{enumerate}
\end{corrige}

\begin{corrige}[Exercice 7 — Représentations graphiques]
\begin{enumerate}
\item $(D_1)$ passe par $(0\,;-1)$ et $(1\,;1)$ ; $(D_2)$ passe par $(0\,;5)$ et $(5\,;0)$. On trace chaque droite à la règle à partir de ses deux points.
\item On résout $2x-1=-x+5$, soit $3x=6$, d'où $x=2$ et $y=f(2)=2\times 2-1=3$. Donc $I(2\,;3)$.
\item « $(D_1)$ au-dessus de $(D_2)$ » se traduit par $f(x)>g(x)$, c'est-à-dire $2x-1>-x+5$, soit $3x>6$, d'où $x>2$. L'ensemble des solutions est l'intervalle $]2\,;+\infty[$.
\item Graphiquement, pour $x>2$, la droite $(D_1)$ est bien située au-dessus de $(D_2)$, ce qui confirme le calcul.
\end{enumerate}
\end{corrige}

\begin{corrige}[Exercice 8 — Application affine par intervalles]
\begin{enumerate}
\item Pour $1$ h $30$ min $=1{,}5$ h (premier intervalle) : $P=200\times 1{,}5=300$ F. Pour $5$ h (deuxième intervalle) : $P=400+100\times(5-2)=400+300=700$ F.
\item Pour $0\le x\le 2$ : $P(x)=200x$. Pour $2<x\le 6$ : $P(x)=400+100(x-2)=100x+200$. Donc
\[ P(x)=\begin{cases} 200x & \text{si } 0\le x\le 2\\ 100x+200 & \text{si } 2<x\le 6 \end{cases} \]
On remarque la continuité en $x=2$ : $200\times 2=400$ et $100\times 2+200=400$.
\item Le graphique est formé de deux segments : de $(0\,;0)$ à $(2\,;400)$, puis de $(2\,;400)$ à $(6\,;800)$.
\end{enumerate}
\end{corrige}

\courssub{Je me prépare au BEPC}

\begin{corrige}[Exercice 9 — Type BEPC 1 : étude complète d'une application affine]
\begin{enumerate}
\item \begin{enumerate}
\item $f(0)=-3\times 0+7=7$ ; $f(3)=-3\times 3+7=-9+7=-2$.
\item On résout $-3x+7=-2$, soit $-3x=-9$, d'où $x=3$. L'antécédent de $-2$ est $3$.
\end{enumerate}
\item Soient $x_1$ et $x_2$ deux réels tels que $x_1<x_2$. En multipliant par $-3$ (négatif), l'ordre change : $-3x_1>-3x_2$, puis en ajoutant $7$ : $-3x_1+7>-3x_2+7$, c'est-à-dire $f(x_1)>f(x_2)$. Donc $f$ est décroissante sur $\mathbb{R}$ (ce qui confirme la règle : $a=-3<0$).
\item $-3x+7>1 \iff -3x>-6 \iff x<2$ (on divise par $-3<0$, le sens de l'inégalité change). L'ensemble des solutions est $]-\infty\,;2[$.
\item $(\mathcal{C}_f)$ est la droite passant par $(0\,;7)$ et $(3\,;-2)$.
\item Graphiquement, le point de $(\mathcal{C}_f)$ d'ordonnée $-2$ a pour abscisse $3$ : on retrouve le résultat de la question 1.b.
\item $(\mathcal{C}_g)$ passe par $O$ et $(1\,;2)$. Intersection : $-3x+7=2x \iff 5x=7 \iff x=\dfrac{7}{5}=1{,}4$ ; alors $y=2\times 1{,}4=2{,}8$. Le point d'intersection est $\left(\dfrac{7}{5}\,;\dfrac{14}{5}\right)$, soit $(1{,}4\,;2{,}8)$.
\end{enumerate}

\textbf{Barème indicatif (sur 10)} : Q1.a : 1 pt ; Q1.b : 1 pt ; Q2 : 1,5 pt ; Q3 : 1,5 pt ; Q4 : 1,5 pt ; Q5 : 1 pt ; Q6 : 2,5 pts.

\textbf{Points de vigilance} : à la question 2, rédiger la démonstration complète avec $x_1<x_2$ (et pas seulement citer le signe de $a$) ; à la question 3, \emph{changer le sens de l'inégalité} en divisant par $-3$ et donner la réponse sous forme d'intervalle ; à la question 5, laisser les traits de lecture apparents ; à la question 6, donner les coordonnées \emph{exactes} $\left(\dfrac{7}{5}\,;\dfrac{14}{5}\right)$ avant la valeur décimale.
\end{corrige}

\begin{corrige}[Exercice 10 — Type BEPC 2 : seuil de rentabilité d'une usine de briquettes]
\begin{enumerate}
\item $C(x)=150x+200\,000$ (application \textbf{affine}, $a=150$, $b=200\,000$) ; $R(x)=250x$ (application \textbf{linéaire}, $a=250$).
\item $C(1\,000)=150\times 1\,000+200\,000=150\,000+200\,000=350\,000$ F ; $R(1\,000)=250\times 1\,000=250\,000$ F. Comme $R(1\,000)<C(1\,000)$, l'usine est en \emph{perte} de $350\,000-250\,000=100\,000$ F pour $1\,000$ briquettes.
\item La droite de $C$ passe par $(0\,;200\,000)$ et $(3\,000\,;650\,000)$ ; celle de $R$ passe par $(0\,;0)$ et $(3\,000\,;750\,000)$. On les trace avec les échelles indiquées.
\item \begin{enumerate}
\item La condition est $R(x)\ge C(x)$, c'est-à-dire $250x\ge 150x+200\,000$.
\item $250x-150x\ge 200\,000 \iff 100x\ge 200\,000 \iff x\ge 2\,000$. L'usine réalise un bénéfice à partir de $2\,000$ briquettes vendues par jour (bénéfice nul à exactement $2\,000$, strictement positif au-delà).
\item Graphiquement, les deux droites se coupent au point d'abscisse $2\,000$ (et d'ordonnée $500\,000$) ; pour $x>2\,000$, la droite de $R$ est au-dessus de celle de $C$.
\end{enumerate}
\item Le bénéfice est $B(x)=R(x)-C(x)=250x-(150x+200\,000)=100x-200\,000$. On veut $100x-200\,000\ge 100\,000$, soit $100x\ge 300\,000$, d'où $x\ge 3\,000$. Le directeur doit produire et vendre au minimum $3\,000$ briquettes par jour.
\end{enumerate}

\textbf{Barème indicatif (sur 10)} : Q1 : 2 pts ; Q2 : 1,5 pt ; Q3 : 2 pts ; Q4.a : 1 pt ; Q4.b : 1,5 pt ; Q4.c : 1 pt ; Q5 : 1 pt.

\textbf{Points de vigilance} : bien distinguer \emph{coût} (affine, avec charges fixes) et \emph{recette} (linéaire) ; à la question 2, conclure par une phrase (« perte de $100\,000$ F ») et pas seulement par des calculs ; au graphique, respecter les échelles et nommer les axes ; à la question 5, définir explicitement le bénéfice $B(x)=R(x)-C(x)$ avant de résoudre.
\end{corrige}

\begin{corrige}[Exercice 11 — Type BEPC 3 : tarification de l'eau par tranches]
\begin{enumerate}
\item Pour $10$ m$^3$ : $1\,000+300\times 10=1\,000+3\,000=4\,000$ F. Pour $30$ m$^3$ : $1\,000+300\times 10+500\times(30-10)=1\,000+3\,000+10\,000=14\,000$ F.
\item \textbf{Tranche 1} ($0\le x\le 10$) : $P(x)=1\,000+300x$. \textbf{Tranche 2} ($10<x\le 30$) : on paie $4\,000$ F pour les $10$ premiers m$^3$, puis $500$ F par m$^3$ au-delà : $P(x)=4\,000+500(x-10)=4\,000+500x-5\,000=500x-1\,000$. \textbf{Tranche 3} ($x>30$) : on paie $14\,000$ F pour les $30$ premiers m$^3$, puis $800$ F par m$^3$ au-delà : $P(x)=14\,000+800(x-30)=14\,000+800x-24\,000=800x-10\,000$.
\item $P(5)=300\times 5+1\,000=1\,500+1\,000=2\,500$ F ; $P(15)=500\times 15-1\,000=7\,500-1\,000=6\,500$ F ; $P(40)=800\times 40-10\,000=32\,000-10\,000=22\,000$ F.
\item Le graphique est formé de trois segments joignant les points $(0\,;1\,000)$, $(10\,;4\,000)$, $(30\,;14\,000)$ et $(40\,;22\,000)$.
\item Oui, le graphique se trace \emph{sans lever le crayon} : aux points d'abscisses $10$ et $30$, les segments se raccordent exactement (par exemple en $x=10$ : $300\times 10+1\,000=4\,000$ et $500\times 10-1\,000=4\,000$). On dit que l'application $P$ est \emph{continue} ; seule la pente (le prix au m$^3$) change à chaque tranche.
\item $9\,000$ F est compris entre $4\,000$ F et $14\,000$ F : on est donc dans la \textbf{tranche 2}. On résout $500x-1\,000=9\,000$, soit $500x=10\,000$, d'où $x=20$. La famille Mbarga a consommé $20$ m$^3$ d'eau.
\end{enumerate}

\textbf{Barème indicatif (sur 10)} : Q1 : 1,5 pt ; Q2 : 2,5 pts ; Q3 : 1,5 pt ; Q4 : 2 pts ; Q5 : 1 pt ; Q6 : 1,5 pt.

\textbf{Points de vigilance} : à la question 2, détailler le calcul tranche par tranche en partant du prix déjà acquis ($4\,000$ F puis $14\,000$ F) ; ne pas oublier l'abonnement de $1\,000$ F ; à la question 6, \emph{justifier le choix de la tranche} en comparant $9\,000$ F aux paliers $4\,000$ F et $14\,000$ F avant de résoudre l'équation ; vérifier que la réponse trouvée ($x=20$) appartient bien à l'intervalle de la tranche utilisée ($10<20\le 30$).
\end{corrige}

\newpage

\coursec{Fiche bilan}

\begin{popfigure}
\begin{tikzpicture}[mindmap, grow cyclic, every node/.style={concept, text=white, font=\small\bfseries}, root concept/.append style={concept color=popdark, font=\bfseries}, level 1/.append style={level distance=3.6cm, sibling angle=51.4}, level 2/.append style={level distance=2.2cm, font=\footnotesize}]
\node[root concept]{Applications linéaires et affines}
  child[concept color=popblue]{ node {Définitions}
    child { node {Linéaire : $f(x)=ax$} }
    child { node {Affine : $f(x)=ax+b$} }
  }
  child[concept color=poporange]{ node {Images et antécédents}
    child { node {Image : calculer $f(x)$} }
    child { node {Antécédent : résoudre $ax+b=y$} }
  }
  child[concept color=popgreen]{ node {Sens de variation}
    child { node {$a>0$ : croissante} }
    child { node {$a<0$ : décroissante} }
    child { node {$a=0$ : constante} }
  }
  child[concept color=poppurple]{ node {Représentation graphique}
    child { node {Droite passant par $(0;b)$} }
    child { node {2 points suffisent} }
  }
  child[concept color=poppink]{ node {Par intervalles}
    child { node {Un morceau par intervalle} }
  }
  child[concept color=popgold]{ node {Modélisation}
    child { node {Forfaits, tarifs, salaires} }
  }
  child[concept color=popmuted]{ node {Lecture graphique}
    child { node {$a$ : coefficient directeur} }
    child { node {$b$ : ordonnée à l'origine} }
  };
\end{tikzpicture}
\legende{Carte mentale de la leçon : applications linéaires et affines.}
\end{popfigure}

\begin{bilan}
\textbf{Définitions clés}
\begin{itemize}
  \item Une \cle{application linéaire} est une application de la forme $f(x)=ax$, où $a$ est un nombre réel fixé. Sa représentation graphique est une droite qui passe par l'\cle{origine} du repère.
  \item Une \cle{application affine} est une application de la forme $f(x)=ax+b$, où $a$ et $b$ sont des réels fixés. $a$ est le \cle{coefficient directeur}, $b$ est l'\cle{ordonnée à l'origine}.
  \item Une application linéaire est un cas particulier d'application affine (avec $b=0$).
  \item L'\cle{image} de $x$ est $f(x)$ ; un \cle{antécédent} de $y$ est un nombre $x$ tel que $f(x)=y$.
\end{itemize}

\textbf{Formules et propriétés à connaître par cœur}
\begin{center}
\begin{tabularx}{\linewidth}{|l|X|}\hline
\rowcolor{popblueL} \textbf{Notion} & \textbf{Résultat} \\ \hline
Image de $x$ & $f(x)=ax+b$ (on remplace $x$ par sa valeur) \\ \hline
Antécédent de $y$ & On résout l'équation $ax+b=y$, soit $x=\dfrac{y-b}{a}$ (si $a\neq 0$) \\ \hline
Sens de variation & $a>0$ : croissante ; $a<0$ : décroissante ; $a=0$ : constante \\ \hline
Coefficient directeur & $a=\dfrac{f(x_2)-f(x_1)}{x_2-x_1}$ pour $x_1\neq x_2$ \\ \hline
Ordonnée à l'origine & $b=f(0)$ : la droite coupe l'axe des ordonnées au point $(0\,;\,b)$ \\ \hline
Proportionnalité & Un tableau de valeurs de $f$ est un tableau de proportionnalité $\iff f$ est linéaire \\ \hline
\end{tabularx}
\end{center}

\textbf{Méthodes express}
\begin{itemize}
  \item \textbf{Tracer la droite} de $f(x)=ax+b$ : placer le point $(0\,;\,b)$, calculer l'image d'un autre nombre, placer le second point, tracer la droite.
  \item \textbf{Trouver $a$ et $b$} à partir de deux points : calculer $a=\dfrac{y_2-y_1}{x_2-x_1}$, puis $b=y_1-a x_1$.
  \item \textbf{Lecture graphique} : $b$ se lit à l'intersection avec l'axe des ordonnées ; $a$ se lit en avançant de $1$ en abscisse et en comptant la variation en ordonnée.
  \item \textbf{Application affine par intervalles} : découper l'ensemble de définition en intervalles, écrire une expression $a_i x+b_i$ sur chacun, tracer chaque morceau séparément (attention aux extrémités).
  \item \textbf{Modéliser} : repérer la partie fixe (c'est $b$) et la partie proportionnelle (c'est $ax$) dans un tarif ou un salaire.
\end{itemize}
\end{bilan}

\begin{aretenir}[Checklist avant le BEPC]
Avant l'épreuve, vérifie que tu sais :
\begin{itemize}
  \item[$\square$] Reconnaître si une application est linéaire, affine ou ni l'un ni l'autre.
  \item[$\square$] Identifier le coefficient directeur $a$ et l'ordonnée à l'origine $b$ dans $f(x)=ax+b$.
  \item[$\square$] Calculer l'image d'un nombre donné sans erreur de signe.
  \item[$\square$] Déterminer un antécédent en résolvant l'équation $ax+b=y$.
  \item[$\square$] Donner le sens de variation d'une application affine selon le signe de $a$.
  \item[$\square$] Calculer $a=\dfrac{f(x_2)-f(x_1)}{x_2-x_1}$ à partir de deux points.
  \item[$\square$] Tracer la représentation graphique d'une application affine avec deux points bien choisis.
  \item[$\square$] Lire graphiquement $a$, $b$, une image ou un antécédent sur une droite tracée.
  \item[$\square$] Vérifier qu'une application linéaire correspond à une situation de proportionnalité.
  \item[$\square$] Étudier et représenter une application affine définie par intervalles.
  \item[$\square$] Modéliser une situation concrète (forfait téléphonique, tarif de taxi, salaire) par une application affine et interpréter le résultat.
  \item[$\square$] Rédiger clairement les étapes de la démarche et conclure par une phrase.
\end{itemize}
\end{aretenir}

\findecours

\end{document}
